\documentclass[11pt]{article}
\usepackage[T1]{fontenc}
\usepackage[utf8]{inputenc}
\usepackage{amsmath,amsthm,mathtools}
\usepackage{newpxtext,newpxmath}
\usepackage[margin=1in,headheight=15pt]{geometry}
\usepackage{microtype}
\usepackage{xcolor,booktabs,array,longtable,enumitem}
\usepackage{fancyhdr}
\usepackage[most]{tcolorbox}
\usepackage{hyperref}
\definecolor{inkblue}{HTML}{16465A}
\definecolor{paleblue}{HTML}{F0F6F8}
\hypersetup{colorlinks=true,linkcolor=inkblue,citecolor=inkblue,urlcolor=inkblue,
 pdftitle={L-convex polyominoes and a colored-partition comparison},
 pdfauthor={Research manuscript prepared with ChatGPT for Vladimir Reshetnikov}}
\pagestyle{fancy}\fancyhf{}
\fancyhead[L]{\small\textsc{L-convex polyominoes}}
\fancyhead[R]{\small\textsc{OEIS A126764}}
\fancyfoot[C]{\thepage}
\setlength{\parskip}{0.35em}
\setlength{\emergencystretch}{2em}
\numberwithin{equation}{section}
\newtheorem{theorem}{Theorem}[section]
\newtheorem{lemma}[theorem]{Lemma}
\newtheorem{proposition}[theorem]{Proposition}
\newtheorem{corollary}[theorem]{Corollary}
\theoremstyle{definition}\newtheorem{definition}[theorem]{Definition}
\theoremstyle{remark}\newtheorem{remark}[theorem]{Remark}
\newcommand{\NN}{\mathbb N}
\newcommand{\CC}{\mathbb C}
\newcommand{\OO}{\mathrm O}
\newcommand{\coef}[1]{[q^{#1}]}
\newcommand{\pp}[1]{(q;q)_{#1}}
\newcommand{\Acal}{\mathcal A}
\newcommand{\Lcal}{\mathcal L}
\newcommand{\pre}{\preccurlyeq}
\newcommand{\ee}{\mathrm e}
\newcommand{\ii}{\mathrm i}% Part II (batch 77P5)
\newcommand{\dd}{\mathrm d}% Part II (batch 77P5)
\DeclareMathOperator{\ord}{ord}
\newenvironment{writenote}[1][]{\par\smallskip\begingroup\small
 \noindent\textbf{[write]}\if\relax\detokenize{#1}\relax\else~\textbf{#1.}\fi\ }%
 {\par\endgroup\smallskip}
\newtcolorbox{resultbox}{colback=paleblue,colframe=inkblue,boxrule=0.6pt,
 arc=1mm,left=3mm,right=3mm,top=2mm,bottom=2mm,breakable}
\title{\color{inkblue}\bfseries L-convex polyominoes\par
 and a colored-partition comparison\vspace{0.6em}\par
 \large A proof of the A126764 asymptotic, three correction terms,\par
 an all-orders radial expansion, and inverse asymptotics\vspace{0.4em}\par
 \normalsize Part~II: the all-orders coefficient theorem\par
 Part~III: a complete transfer proof}
\author{Research manuscript prepared for Vladimir Reshetnikov}
\date{October 1, 2026}
\begin{document}
\maketitle
\thispagestyle{empty}
\begin{abstract}
Let $a_n$ count fixed L-convex polyominoes of area $n$, with $a_0=1$.
Starting from the established area generating function, we prove the
Guttmann--Kot\v e\v sovec asymptotic recorded as a conjecture in OEIS A126764:
\[
 a_n\sim\frac{13\sqrt2}{768\,n^{3/2}}
       \exp\!\left(\pi\sqrt{\frac{13n}{6}}\right).
\]
The main device is a positive comparison with a colored-partition
Euler product $R(q)$, given explicitly below.
If $p_n=[q^n]R(q)$, we obtain the finite, coefficientwise bound
$0\le p_n-4a_n\le3(p_n-2p_{n-1}+p_{n-2})$ for every $n\ge3$.
A higher finite-cut identity then proves three explicit correction terms,
with relative remainder $\OO(n^{-7/4})$.
Independently, a finite recurrence computes the complete radial Poincar\'e
expansion of $\Acal(\ee^{-t})$ to any prescribed order.
We also derive an asymptotic inverse, including the necessary distinction
between a real interpolation and an integer threshold.
Exact integer checks through $n=2000$ and symbolic coefficient checks
accompany the manuscript. No numerical fit enters the proofs.
\end{abstract}
\begin{resultbox}
\textbf{Scope and status.} The proofs below are conventional mathematical
proofs, not Lean certificates or a claim of independent peer review.
The exact enumerative generating function is a cited input; its asymptotic
analysis is supplied here. The all-orders \emph{radial} expansion is proved.
An all-orders \emph{coefficient} expansion, and exponentially small
transseries sectors for $a_n$, are not claimed.
\end{resultbox}
\begin{writenote}[Added 2 October 2026, batch~77P5]
The box describes Part~I. Part~II (Sections~\ref{lcp:ao:sec:intro}--\ref{lcp:ao:sec:provenance}),
a second manuscript, proves the all-orders \emph{coefficient} expansion
(Theorem~\ref{lcp:ao:thm:main}). Exponentially small transseries sectors
remain unclaimed in both parts.
\end{writenote}

\begin{writenote}[Added 5 October 2026, batch~101]
Since batch~101 the report has a third part, Part~III
(Sections~\ref{lcp:tr:sec:intro}--\ref{lcp:tr:sec:provenance}). It proves
the all-orders coefficient expansion again (Theorem~\ref{lcp:tr:thm:main}),
by a third route whose estimates are all written out. Exponentially small
transseries sectors remain unclaimed in all three parts.
\end{writenote}
\clearpage
{\small\setlength{\parskip}{0pt}\tableofcontents}
\clearpage

\phantomsection\addcontentsline{toc}{part}{Part I.\texorpdfstring{\quad}{ } L-convex polyominoes and a colored-partition comparison}
\section{The problem and the main results}
\subsection{The sequence and the source boundary}
An L-convex polyomino is a convex polyomino in which any two cells can be
joined inside the polyomino by a path with at most one turn.
We count fixed polyominoes: translations are identified, but rotations
and reflections are not generally identified. This is the convention in
OEIS A126764 \cite{oeis}. The initial terms are
\[
1,1,2,6,15,35,76,156,310,590,1098,1984,3515,6094,\ldots.
\]
The zero-area term is a generating-function convention.

Castiglione, Frosini, Munarini, Restivo, and Rinaldi established an area
generating function \cite{castiglione}; we use its explicit recurrence form
as reproduced in Guttmann and Kot\v e\v sovec \cite{gk}, equation (1.1).
The latter authors proposed the leading asymptotic from a numerical
analysis. The OEIS entry inspected for this report still labels that
formula a conjecture \cite{oeis}. A targeted search did not locate an
earlier proof of the specific area asymptotic. This search is not a
comprehensive priority certification.

The ProveIt repository \cite{proveit} was inspected as the proposed research
context. A targeted repository search for A126764 returned no matching
file. The present manuscript is a standalone contribution: it does not
assume any unverified theorem from another repository report. Enumeration
by perimeter, and the asymptotics of 201-avoiding ascent sequences discussed
alongside L-convex polyominoes in \cite{gk}, are separate questions and are
not asserted to be solved here.

\begin{writenote}[Added 2 October 2026, batch~75]
This report is a single manuscript, printed in full: manuscript~05 of
batch~75 of ProveIt's incoming-reports intake, delivered as the archive
\texttt{A126764\_\allowbreak Lconvex\_\allowbreak Research.zip} in commit
\texttt{4b874cea0} (main file \texttt{lconvex\_asymptotics.tex}, now
\texttt{article.tex}) and placed in commit \texttt{6ea60e367} in the
research-report collection, among the OEIS-sequence asymptotics reports.
The manuscript names no ProveIt commit: its repository input was an
overview and a targeted search for A126764 on 1~October 2026, so it has no
pin. That search is still accurate for the repository it saw; the only
match now is this report. Its PDF metadata records it as ``prepared with
ChatGPT''. Nothing in this report is formalized in Lean or Rocq, and
placement in the collection confers no formal status on any statement here.
The writing step changed only the labels (every label now carries the
prefix \texttt{lcp:}), set the bibliography ragged-right, and added these
\textsf{[write]} notes and two references
\cite{lcp-tai,lcp-klarner}; no statement, proof or number of the manuscript
was altered. The polyomino report nearest in subject,
\texttt{polyomino-\allowbreak growth-\allowbreak finite-\allowbreak prefix-\allowbreak corrections}
\cite{lcp-klarner}, bounds Klarner's growth constant $\lambda\le4.498$ for
\emph{all} polyominoes (exponential growth, with a Lean development
\nolinkurl{Combinatorics/Polyominoes/KlarnerConstant} for the weaker bound
$9047/2000$); the L-convex class grows only like $\ee^{B\sqrt n}$, so the
two constants are unrelated and no theorem is shared. The report
\texttt{skew-\allowbreak partition-\allowbreak continued-\allowbreak fractions}
touches parallelogram polyominoes (A006958) only.
\end{writenote}

\begin{writenote}[Added 2 October 2026, batch~77P1]
The 201-avoiding ascent sequences that the preceding paragraph sets aside
are now treated in the sibling report
\texttt{a202062-\allowbreak ascent-\allowbreak 201-\allowbreak enumeration}
(manuscript~43 of batch~77): it proves the Guttmann--Kot\v e\v sovec
algebraic (cubic) generating function for A202062 and its all-orders
coefficient asymptotics. The two problems share only the paper
\cite{gk}; no theorem, constant or method of this report is used there,
and nothing here changes.
\end{writenote}

\begin{writenote}[Added 2 October 2026, batch~77P5]
The note above describes Part~I, Sections~1--12 and Appendices~A and~B.
Since batch~77 the report has a second part, Part~II
(Sections~\ref{lcp:ao:sec:intro}--\ref{lcp:ao:sec:provenance}), printed in
full from manuscript~31 of batch~77 (archive
\texttt{lconvex-\allowbreak asymptotics-\allowbreak reproducibility.zip},
arrival commit \texttt{096ee7b87}, placement commit \texttt{4f11bc9c0}; no
pin). It proves the all-orders coefficient expansion that Part~I leaves
open, by an exact decomposition $\Acal=R\,D^2+(\text{an exponentially
smaller remainder})$ in Part~I's letters. Its letters clash with Part~I's;
Table~\ref{lcp:ao:tab:notation} resolves them.
\end{writenote}

\begin{writenote}[Added 5 October 2026, batch~101]
Part~III (Sections~\ref{lcp:tr:sec:intro}--\ref{lcp:tr:sec:provenance})
prints bundle Report~95 of batch~101 in full (archive
\texttt{L\_\allowbreak Convex\_\allowbreak Polyomino\_\allowbreak Area\_\allowbreak Asymptotics\_\allowbreak Source.zip},
arrival commit \texttt{60f54ea06}, placement commit \texttt{f7c612c72}; no
pin). It proves the decomposition $\Acal=1+R^{\mathrm I}D^2+(\text{a
positive remainder})$ in Part~I's letters by a duality argument, and the
all-orders coefficient expansion by positivity, the common partition
factor $(q;q)_N$ and Chebyshev interpolation. It knows Part~I's manuscript
but not Part~II. Its letters clash with both parts;
Table~\ref{lcp:tr:tab:notation} resolves them.
\end{writenote}

\subsection{Constants and coefficient asymptotics}
Throughout, put
\begin{equation}\label{lcp:eq:constants}
 K=\frac{13\pi^2}{24},\qquad B=2\sqrt K=\pi\sqrt{\frac{13}{6}},
 \qquad c=\frac{13\sqrt2}{768}.
\end{equation}
The letter $K$ denotes a constant, not a generating function.

\begin{theorem}[Coefficient expansion]\label{lcp:thm:main}
As $n\to\infty$ through the positive integers,
\begin{equation}\label{lcp:eq:main}
 a_n=c\,n^{-3/2}\ee^{B\sqrt n}
 \left(1-\frac{d}{\sqrt n}+\frac{e}{n}
              +\frac{f}{n^{3/2}}+\OO(n^{-7/4})\right),
\end{equation}
where
\begin{align}
 d&=\frac{\sqrt K}{6}+\frac{3}{2\sqrt K}
     =\frac B{12}+\frac3B,\label{lcp:eq:d}\\
 e&=\frac12+\frac{3}{4K}-\frac{35K}{72},\label{lcp:eq:e}\\
 f&=\frac{755K^{3/2}}{1296}+\frac{175\sqrt K}{72}
                  -\frac{5}{8\sqrt K}.\label{lcp:eq:f}
\end{align}
Numerically,
\[
 d=1.0341052176262747\ldots,\quad
 e=-1.9584764928967362\ldots,\quad
 f=12.550445779146218\ldots.
\]
In particular, the leading asymptotic stated in the abstract holds.
\end{theorem}

There are two complementary proof routes. The first proves the leading
formula and its $n^{-1/2}$ correction by an explicit positive squeeze.
The second uses a finite rational approximation and a contour estimate;
it proves all the terms in Theorem~\ref{lcp:thm:main}.

\begin{theorem}[A finite coefficient squeeze]\label{lcp:thm:squeeze}
Define
\begin{equation}\label{lcp:eq:R}
 R(q)=\frac{(-q;q^2)_\infty}{\pp{\infty}^3}
     =\frac{(q^2;q^2)_\infty^2}
            {(q;q)_\infty^4(q^4;q^4)_\infty}
     =\sum_{n\ge0}p_nq^n,
 \qquad J(q)=(1-q)^2R(q)=\sum_{n\ge0}j_nq^n.
\end{equation}
Then, for every $n\ge3$,
\begin{equation}\label{lcp:eq:squeeze}
 0\le p_n-4a_n\le3j_n
   =3(p_n-2p_{n-1}+p_{n-2}).
\end{equation}
Both $R$ and $J$ have nonnegative integer coefficients. Moreover,
\begin{equation}\label{lcp:eq:deficit}
 n\left(1-\frac{4a_n}{p_n}\right)\longrightarrow\frac K2
 =\frac{13\pi^2}{48}.
\end{equation}
\end{theorem}

The product $R$ counts partitions with four colors available for an odd
part, two for a part congruent to $2$ modulo $4$, and three for a multiple
of $4$. Thus the constant $13$ comes from the average number $13/4$ of
colors, not from a fit to a decimal growth constant. The factor $1/4$
in the comparison is forced by two boundary conditions of the recurrence.
We do not assert a four-to-one combinatorial map; finding one is a natural
further problem.

\subsection{What the all-orders statement means}
Let $t>0$ and $q=\ee^{-t}$. We will prove
\begin{equation}\label{lcp:eq:radialintro}
 \Acal(\ee^{-t})\sim
 \frac{t^{3/2}}{4(2\pi)^{3/2}}
 \exp\!\left(\frac Kt-\frac t6\right)\mathcal F(t),
\end{equation}
where $\mathcal F$ is a formally specified series beginning
\begin{align}\label{lcp:eq:Fintro}
 \mathcal F(t)={}&1-\frac12t^2+\frac12t^3-\frac{41}{48}t^4
       +\frac74t^5-\frac{6317}{1440}t^6\notag\\
       &+\frac{1529}{120}t^7-\frac{1702513}{40320}t^8+\cdots.
\end{align}
Here $\sim$ is meant to all fixed algebraic orders along the positive
real axis. It is not a claim that this series converges, nor that a
real-axis expansion alone implies coefficient asymptotics. The complex
estimate needed for the three coefficient corrections is proved separately.

\section{An exact recurrence with a positive normalization}
Write $(a;q)_m=\prod_{j=0}^{m-1}(1-aq^j)$, and set $\pp0=1$.
The enumerative input is
\begin{equation}\label{lcp:eq:enum}
 \Acal(q)=1+\sum_{n\ge1}
 \frac{q^n h_n(q)}{\pp{n-1}^2(1-q^n)},
\end{equation}
where
\begin{equation}\label{lcp:eq:hrec}
 h_0=1,\qquad h_1=1,\qquad
 h_{n+1}=2h_n-(1-q^n)^2h_{n-1}\quad(n\ge1).
\end{equation}
Indeed, $h_n=f_{n-1}$ for $n\ge1$ in the notation of \cite{gk}, and
$h_2=1+2q-q^2$. Introducing $h_0$ makes the recurrence valid already
at $n=1$.

The auxiliary solution
\begin{equation}\label{lcp:eq:Qrec}
 Q_0=1,\qquad Q_1=2,\qquad
 Q_{n+1}=2Q_n-(1-q^n)^2Q_{n-1}
\end{equation}
will produce the comparison product. For a general solution $y_n$ of the
same recurrence, define
\[
 b_n=\frac{y_n}{\pp n^2},\qquad e_n=b_n-b_{n-1},\qquad
 T_n(q)=\frac1{(1-q^n)^2}-1.
\]
We reserve $d,e,f$ without subscripts for the constants in
Theorem~\ref{lcp:thm:main}.

\begin{lemma}[Positive transfer]\label{lcp:lem:positive}
For $n\ge2$,
\begin{align}\label{lcp:eq:positive}
 b_n&=\frac{b_{n-1}+e_{n-1}}{(1-q^n)^2},\\
 e_n&=e_{n-1}+T_n(q)(b_{n-1}+e_{n-1}).\notag
\end{align}
Every entry of this linear transfer rule is a power series with nonnegative
coefficients. Consequently, coefficientwise inequalities between two
pairs $(b_m,e_m)$ propagate to all later pairs and to their limiting slopes,
whenever the slopes exist.
\end{lemma}
\begin{proof}
Dividing the recurrence by $\pp{n-1}^2$ gives
$(1-q^n)^2b_n=2b_{n-1}-b_{n-2}=b_{n-1}+e_{n-1}$.
Subtract $b_{n-1}$. Finally,
$T_n=\sum_{r\ge1}(r+1)q^{nr}$ has nonnegative coefficients.
\end{proof}

For $h$, the initial normalized pair is
\[
 b_1=(1-q)^{-2},\qquad e_1=(1-q)^{-2}-1,
\]
and both entries are nonnegative. The same holds for $Q$, with
$b_1=2(1-q)^{-2}$ and $e_1=2(1-q)^{-2}-1$.
Throughout, $F\pre G$ means that every coefficient of $G-F$ is
nonnegative. This is a stronger relation than $F(q)\le G(q)$ for
$0<q<1$. We will be careful not to multiply a coefficientwise inequality
by a polynomial with negative coefficients.

\section{Limiting slopes and the exact Euler product}
\subsection{Existence and analyticity of the slopes}
\begin{lemma}\label{lcp:lem:slope}
For a solution with polynomial initial values and every $0<r<1$, there
are analytic functions $s(q),b(q)$ on $|q|<1$ such that, uniformly on
$|q|\le r$,
\[
 y_n(q)=s(q)n+b(q)+\OO_r(n^2r^n).
\]
In particular, the normalized limiting slope
\begin{equation}\label{lcp:eq:slope}
 \Lcal_y(q)=\lim_{n\to\infty}
 \left(\frac{y_n}{\pp n^2}-\frac{y_{n-1}}{\pp{n-1}^2}\right)
 =\frac{s(q)}{\pp\infty^2}
\end{equation}
exists locally uniformly. The same limit exists coefficientwise as a
formal power series.
\end{lemma}
\begin{proof}
Put $v_j=2q^j-q^{2j}$. Twice summing
$y_{j+1}-2y_j+y_{j-1}=v_jy_{j-1}$ gives
\[
 y_n=y_0+n(y_1-y_0)+\sum_{j=1}^{n-1}(n-j)v_jy_{j-1}.
\]
On $|q|\le r$, $|v_j|\le3r^j$. Divide by $n+1$ and apply the elementary
finite-product version of discrete Gronwall's inequality. The resulting
bound is a constant times $\prod_{j\ge1}(1+3jr^j)<\infty$.
Thus $y_n=\OO_r(n+1)$. The increments $y_{n+1}-y_n$ have a uniformly
convergent sum of successive differences, with exponentially small tails.
One more summation yields the displayed affine expansion.
Since $\pp n/\pp\infty=1+\OO_r(r^n)$, normalization gives
\eqref{lcp:eq:slope}. Finally, in \eqref{lcp:eq:positive} the change in $e_n$
has valuation at least $n$. Each fixed coefficient therefore stabilizes.
\end{proof}

Write $L=\Lcal_h$. If
\begin{equation}\label{lcp:eq:C}
 C(z;q)=\sum_{n\ge0}\frac{h_n(q)}{\pp n^2}z^n,
 \qquad C_j(q)=C(q^j;q),
\end{equation}
then $C$ has radius of convergence at least $1$ in $z$ and a possible
double pole at $z=1$ with leading coefficient $L$.
Equation~\eqref{lcp:eq:enum} becomes
\begin{equation}\label{lcp:eq:AC}
 \Acal=1+C_1-C_2.
\end{equation}
Multiplying the normalized recurrence by $z^n$ and summing, including the
initial condition $h_1=1$, yields
\begin{equation}\label{lcp:eq:Cdiff}
 (1-z)^2C(z;q)-2C(qz;q)+C(q^2z;q)=-z.
\end{equation}
Taking $z\to1$ is now legitimate and gives
\begin{equation}\label{lcp:eq:Lidentity}
 L=2C_1-C_2-1,\qquad \Acal=L-C_1+2.
\end{equation}
The pole term must not be dropped in this limit.

\subsection{Solving the auxiliary generating function}
Define $G(z)=\sum_{n\ge0}Q_nz^n/\pp n$. Directly from
\eqref{lcp:eq:Qrec},
\[
 (1-z)^2G(z)=(1+qz^2)G(qz),\qquad G(0)=1.
\]
Iteration gives the exact identity
\begin{equation}\label{lcp:eq:Gproduct}
 G(z)=\frac{(-qz^2;q^2)_\infty}{(z;q)_\infty^2}.
\end{equation}
At $z=1$, the coefficient of $(1-z)^{-2}$ is
$(-q;q^2)_\infty/\pp\infty^2$. It follows, either by subtracting the
principal part or by Lemma~\ref{lcp:lem:slope}, that
\begin{equation}\label{lcp:eq:Qproduct}
 \Lcal_Q=\frac{(-q;q^2)_\infty}{\pp\infty^3}=R(q).
\end{equation}
The second expression for $R$ in \eqref{lcp:eq:R} follows from
$(-q;q^2)_\infty=(q^2;q^2)_\infty^2/
((q;q)_\infty(q^4;q^4)_\infty)$.
In particular,
\begin{equation}\label{lcp:eq:colors}
 R(q)=\prod_{j\ge1}(1-q^j)^{-b_j},\qquad
 b_j=4-2\mathbf1_{2\mid j}+\mathbf1_{4\mid j}.
\end{equation}
This proves the colored-partition interpretation stated earlier.

\section{A coefficientwise proof of the leading conjecture}
\subsection{An exact positive defect}
Let $v_n=Q_n-2h_n$. Then
$v_0=-1$, $v_1=0$, $v_2=(1-q)^2$. Its normalized pair at index $2$ is
\[
 b_2=e_2=(1-q^2)^{-2}.
\]
The positive transfer lemma shows that its slope
\begin{equation}\label{lcp:eq:E}
 E=R-2L
\end{equation}
has nonnegative coefficients. Combining \eqref{lcp:eq:AC} and
\eqref{lcp:eq:Lidentity} gives the exact decomposition
\begin{equation}\label{lcp:eq:defectexact}
 R-4\Acal+6=E+2C_2.
\end{equation}
Both terms on the right are nonnegative power series. The constant $6$
only corrects the chosen zero-area terms.

\subsection{A positive majorant for the defect}
We prove the bounds
\begin{equation}\label{lcp:eq:majorants}
 0\pre E\pre J,
 \qquad 0\pre C_2\pre J+q+2q^2.
\end{equation}
These statements require a little care: a real-axis bound by
$(1-q)^2R$ would not by itself imply a bound on coefficients.

First multiply the normalized $Q$ solution by $(1-q)^2$, and denote the
resulting pair by $(\widetilde b_n,\widetilde e_n)$.
At $n=2$,
\begin{align*}
 \widetilde b_2&=\frac{3+2q-q^2}{(1-q^2)^2},\\
 \widetilde e_2&=\frac{1+2q+3q^2-2q^4}{(1-q^2)^2}.
\end{align*}
Subtracting the normalized pair for $v$ gives
\[
 \frac{2+2q-q^2}{(1-q^2)^2},\qquad
 \frac{2q+3q^2-2q^4}{(1-q^2)^2}.
\]
Both are nonnegative power series: use
$2+2q-q^2=(1+2q)+(1-q^2)$ and
$3q^2-2q^4=q^2+2q^2(1-q^2)$.
The transfer lemma and passage to slopes prove $E\pre J$.
The displayed pair for $\widetilde b_2,\widetilde e_2$ itself is
nonnegative by the same decompositions.

Next, every coefficient of $(1-q^n)h_n/\pp n^2$ is nonnegative.
For $n=1$ this is immediate; for $n\ge2$ it follows by removing only
one of the two factors $(1-q^n)$ in the denominator in
\eqref{lcp:eq:positive}. Hence $C_2-C_3\succeq0$.
Setting $z=q$ in \eqref{lcp:eq:Cdiff} gives
\[
 2C_2-C_3=(1-q)^2C_1+q,
 \qquad C_2\pre(1-q)^2C_1+q.
\]
It remains to compare the last expression with $J$.

Let $C^Q_j=\sum_{n\ge0}Q_nq^{jn}/\pp n^2$.
For the $Q$ solution the initial-condition term in \eqref{lcp:eq:Cdiff}
is zero, so $R=2C^Q_1-C^Q_2$.
The solution $u_n=Q_n-h_n$ has $u_0=0,u_1=1$; after normalization
and multiplication by $(1-q)^2$ its initial pair at $1$ is $(1,1)$.
Consequently,
\[
 (1-q)^2(C^Q_1-C_1)\succeq0.
\]
Also
\[
 (1-q)^2(C^Q_1-C^Q_2)
 =\sum_{n\ge1}q^n(1-q^n)\widetilde b_n\succeq-2q^2.
\]
For completeness, the $n=1$ summand is $2q-2q^2$.
For $n=2$ the remaining factor is
$(3+2q-q^2)/(1-q^2)=1+(2+2q)/(1-q^2)$, which is nonnegative.
For $n\ge3$, positivity follows from the nonnegative pair at index $2$
and \eqref{lcp:eq:positive}.
These two inequalities imply
$(1-q)^2C_1\pre J+2q^2$, establishing \eqref{lcp:eq:majorants}.

Taking coefficients of \eqref{lcp:eq:defectexact} now proves
\eqref{lcp:eq:squeeze}. In particular, the finite bound is independent of
any asymptotic approximation.

\subsection{Why the squeeze is asymptotically sharp enough}
The analytic estimates in Section~\ref{lcp:sec:saddle} imply
\begin{align}\label{lcp:eq:pfirst}
 p_n&=4c\,n^{-3/2}\ee^{B\sqrt n}
       \left(1-\frac d{\sqrt n}+\OO(n^{-1})\right),\\
 j_n&\sim\frac Kn p_n.\label{lcp:eq:j}
\end{align}
Therefore the squeeze alone gives
\[
 a_n=c\,n^{-3/2}\ee^{B\sqrt n}
       \left(1-\frac d{\sqrt n}+\OO(n^{-1})\right).
\]
This already proves the full proposed leading constant and the first
correction. The further work below is needed for $e$ and $f$, not for
the leading conjecture.

\section{Finite cuts: an exact rational approximation}
\subsection{General positive remainder series}
The discrete Wronskian of $h$ and $Q$ is
\begin{equation}\label{lcp:eq:wronskian}
 h_mQ_{m+1}-h_{m+1}Q_m=\pp m^2.
\end{equation}
At $m=0$ both sides are $1$; the recurrence multiplies each succeeding
Wronskian by $(1-q^m)^2$, proving the identity by induction.
Define
\begin{equation}\label{lcp:eq:Em}
 E_m=h_mR-Q_mL\qquad(m\ge1).
\end{equation}
Thus $E_1=E$.

\begin{lemma}[A positive tail with a small radial bound]\label{lcp:lem:tail}
The series $E_m$ has nonnegative coefficients. For real $0<q<1$,
\begin{equation}\label{lcp:eq:tailbound}
 0\le E_m(q)\le\pp m^2R(q).
\end{equation}
Furthermore, for each fixed $m$,
\begin{equation}\label{lcp:eq:Ctail}
 C_{m+1}(\ee^{-t})=\OO_m(t^{2m}R(\ee^{-t})+1)
 \quad(t\downarrow0).
\end{equation}
\end{lemma}
\begin{proof}
For $n\ge m$, the solution $w_n=h_mQ_n-Q_mh_n$ has
$w_m=0,w_{m+1}=\pp m^2$. Its normalized pair at $m+1$ is
$(1-q^{m+1})^{-2}$ in both coordinates. Hence it is positive and its
limiting slope, $E_m$, is positive.

For the real bound, let $z_m=0,z_{m+1}=1$ satisfy the original recurrence
for later indices. Then $w_n=\pp m^2z_n$.
For real $q$, the increments of $Q_n$ are positive and nondecreasing:
\[
 Q_{n+1}-Q_n=(Q_n-Q_{n-1})+(2q^n-q^{2n})Q_{n-1}.
\]
Thus $Q_n\ge n+1$ and $\Delta Q_m=Q_m-Q_{m-1}\ge1$.
A second Wronskian calculation yields
\[
 \frac{s_z}{s_Q}
 =Q_m\sum_{k=m}^\infty
 \frac{\prod_{j=m+1}^{k}(1-q^j)^2}{Q_kQ_{k+1}}
 \le\frac1{\Delta Q_m}\le1.
\]
To obtain the middle inequality, use
$Q_k\ge Q_m+(k-m)\Delta Q_m$ and telescope the reciprocals of the
resulting arithmetic progression. Since normalization divides both
slopes by $\pp\infty^2$, \eqref{lcp:eq:tailbound} follows.

Finally, positivity gives $C_1\le C^Q_1\le R$ on the real axis.
Indeed, $Q-h$ is positive after normalization, and
$R-C^Q_1=C^Q_1-C^Q_2\ge0$.
From \eqref{lcp:eq:Cdiff} and $C_{j+1}\ge C_{j+2}$ on the real axis,
\[
 C_{j+1}\le(1-q^j)^2C_j+q^j.
\]
Iterate this inequality a fixed number $m$ of times. The leading
coefficient is $\pp m^2=\OO_m(t^{2m})$, and the finitely many remaining
terms are bounded as $t\downarrow0$.
\end{proof}

\subsection{Eliminating the auxiliary evaluations}
Put $U_m=Q_m-h_m$, and let polynomials $P_m$ be specified by
\[
 P_0=P_1=0,\qquad
 P_{j+1}=2P_j-(1-q^j)^2P_{j-1}+q^j.
\]
Iteration of \eqref{lcp:eq:Cdiff} gives
\[
 C_{m+1}=Q_mC_1-U_m(L+1)-P_m.
\]
Solving this identity together with \eqref{lcp:eq:Em} and
$\Acal=L-C_1+2$ proves the following exact formula.

\begin{proposition}[Finite-cut identity]\label{lcp:prop:cut}
Where $Q_m(q)\ne0$,
\begin{equation}\label{lcp:eq:cut}
 \Acal(q)=\left(\frac{h_m}{Q_m}\right)^2R
 -\frac{h_mE_m}{Q_m^2}-\frac{C_{m+1}}{Q_m}
 +2-\frac{U_m+P_m}{Q_m}.
\end{equation}
After multiplication by $Q_m^2$, this is an identity of analytic
functions on the entire open unit disk and of formal power series.
\end{proposition}

The square in the first term is the precise source of the factor $1/4$:
for every fixed $m\ge1$, $h_m(1)=2^{m-1}$ and $Q_m(1)=2^m$.
The identity also provides an algebraic expansion algorithm and a
route to rigorous coefficient bounds.

For the coefficient theorem we use only $m=3$:
\begin{align}\label{lcp:eq:thirdcut}
 h_3&=1+4q-q^4,\qquad
 Q_3=4+4q+2q^2-2q^4
      =-2(q+1)(q^3-q^2-2),\notag\\
 U_3+P_3&=3+2q+3q^2-q^4,\qquad
 F_3(q)=\left(\frac{1+4q-q^4}{4+4q+2q^2-2q^4}\right)^2.
\end{align}
The cubic $q^3-q^2-2$ has no zero in $|q|\le1$.
A zero there would satisfy $|q|^2|q-1|=2$; equality in the bound
$|q|^2|q-1|\le2$ forces $q=-1$, which does not solve the equation.
Thus $F_3R$ is analytic in $|q|<1$, with only a possible additional
boundary pole at $q=-1$.

\section{Saddle estimates for the comparison product}\label{lcp:sec:saddle}
\subsection{The local modular expansion}
The classical eta transformation, in its Euler-product form, is
\begin{equation}\label{lcp:eq:eta}
 (\ee^{-z};\ee^{-z})_\infty
 =\sqrt{\frac{2\pi}{z}}
  \exp\!\left(-\frac{\pi^2}{6z}+\frac z{24}\right)
  (\ee^{-4\pi^2/z};\ee^{-4\pi^2/z})_\infty,
 \qquad \Re z>0.
\end{equation}
The square root is the branch positive for positive real $z$.
This follows from the usual modular transformation of Dedekind's eta
function; see \cite{dlmf}, equations 27.14.12--27.14.14.
Applying it at $z,2z,4z$ gives, uniformly in any fixed closed sector
$|\arg z|<\pi/2$,
\begin{equation}\label{lcp:eq:Rlocal}
 R(\ee^{-z})=\left(\frac z{2\pi}\right)^{3/2}
 \exp\!\left(\frac Kz-\frac z6\right)
 \left(1+\OO(\ee^{-\kappa/|z|})\right)
\end{equation}
for an appropriate sector-dependent $\kappa>0$.
In particular, there is no omitted algebraic power-series correction
in \eqref{lcp:eq:Rlocal} beyond $\exp(-z/6)$.

\subsection{A minor-arc bound rather than a Tauberian shortcut}
Let $r=\ee^{-t}$. The colored-product form \eqref{lcp:eq:colors} implies
\begin{equation}\label{lcp:eq:minor}
 \frac{|R(r\ee^{i\theta})|}{R(r)}
 \le\exp\!\left[-\frac{c_0}{t}
             \min\!\left\{1,\frac{\theta^2}{t^2}\right\}\right]
 \quad(-\pi\le\theta\le\pi)
\end{equation}
for all sufficiently small $t$ and some $c_0>0$.
Here and below a constant such as $c_0$ is unrelated to $c$ in
\eqref{lcp:eq:constants}.

To verify \eqref{lcp:eq:minor}, restrict the product to integers
$N\le j\le2N$, with $N$ comparable to $1/t$. Then $r^j$ stays in a
compact subinterval of $(0,1)$, and
\[
 \log\frac{|1-r^j\ee^{ij\theta}|}{1-r^j}
 \ge c_1(1-\cos(j\theta)).
\]
In the product for $|R(r\ee^{i\theta})|/R(r)$, every omitted
factor is at most $1$. The elementary
trigonometric-sum estimate
\[
 \sum_{j=N}^{2N}(1-\cos(j\theta))
 \ge c_2N\min\{1,N^2\theta^2\}
\]
finishes the proof. For $|\theta|\le1/N$ it follows from the quadratic
bound for $1-\cos x$; for an intermediate compact range of $N\theta$
it follows by comparison with its strictly positive integral; for larger
$N\theta$ use the geometric-sum bound for $\sum\ee^{ij\theta}$.
This argument also covers the neighborhood of $\theta=\pi$.

Consequently a rational factor with at most fixed-order poles on the
unit circle does not create a competing dominant saddle. The product's
exponential suppression outweighs such polynomial losses.

\subsection{The transfer calculation}
For a local expansion
\[
 H(\ee^{-z})\sim\lambda z^{3/2}\ee^{K/z}
                 \sum_{j\ge0}s_jz^j,
\]
with the preceding minor-arc control, Cauchy's integral is localized at
$t=\sqrt{K/n}$ and angular width $t^{3/2}$.
Termwise saddle integration gives
\begin{equation}\label{lcp:eq:transfer}
 [q^n]H(q)\sim
 \frac{\lambda K}{2\sqrt\pi}\,n^{-3/2}\ee^{2\sqrt{Kn}}
 \sum_{\ell\ge0}n^{-\ell/2}
 \sum_{j+k=\ell}s_jK^{j/2}
       \frac{\beta_k(j)}{(2\sqrt K)^k},
\end{equation}
where
\begin{equation}\label{lcp:eq:betak}
 \beta_k(j)=\frac{(-1)^k}{k!8^k}
 \prod_{r=1}^k\left((5+2j)^2-(2r-1)^2\right),\qquad
 \beta_0(j)=1.
\end{equation}
For clarity, \eqref{lcp:eq:transfer} asserts a finite-order expansion whenever
a corresponding finite-order local expansion is available.

Here is a direct way to check the constants. The model integral for
$z^{3/2+j}\ee^{K/z}$ is the Bessel-type contour expression
\[
 \left(\frac Kn\right)^{(5/2+j)/2}
 I_{-5/2-j}(2\sqrt{Kn}).
\]
Its large positive-argument saddle expansion is
$\ee^x/\sqrt{2\pi x}$ times the polynomial coefficients
\eqref{lcp:eq:betak} in powers of $x^{-1}$. Alternatively those same
coefficients follow by substituting a power series into the Bessel
differential equation. The coefficient $\lambda K/(2\sqrt\pi)$ in
\eqref{lcp:eq:transfer} follows immediately.
One need not identify the original generating function with a Bessel
function: only this local contour model is used.

For a fully controlled remainder, take a central arc
$|\theta|\le t^{3/2-\delta}$ with $0<\delta<1/6$.
On the complement, \eqref{lcp:eq:minor} makes the integral smaller than
any algebraic multiple of its leading saddle value. On the central arc,
put $\theta=t^{3/2}u$; the real part of the quadratic term is a negative
constant times $u^2$, and Taylor remainders are bounded by a polynomial
in $u$ times that Gaussian. Truncation and integration are therefore
valid to every fixed order. This supplies the error justification
behind \eqref{lcp:eq:transfer}.

For $H=R$, use $\lambda=(2\pi)^{-3/2}$ and $s_j=(-1/6)^j/j!$.
This proves \eqref{lcp:eq:pfirst}. Multiplication by
$(1-\ee^{-z})^2=z^2+\OO(z^3)$ gives \eqref{lcp:eq:j}.
It also gives
\begin{equation}\label{lcp:eq:pe2}
 p_n=4c\,n^{-3/2}\ee^{B\sqrt n}
 \left(1-\frac d{\sqrt n}
 +\frac{\frac12+\frac{3}{4K}+\frac K{72}}{n}
 +\OO(n^{-3/2})\right).
\end{equation}
The comparison product itself has an all-orders coefficient expansion;
this fact is not being inferred for $\Acal$ without additional estimates.

\section{Proof of the three correction terms}
\subsection{Bounding the finite-cut remainder on the whole circle}
\begin{proposition}\label{lcp:prop:approxerror}
With $F_3$ from \eqref{lcp:eq:thirdcut},
\begin{equation}\label{lcp:eq:approxerror}
 a_n=[q^n](F_3R)+\OO(p_n n^{-7/4}).
\end{equation}
\end{proposition}
\begin{proof}
Take $q=r\ee^{i\theta}$ with $r=\ee^{-t}$ and $t=\sqrt{K/n}$.
By positivity and Lemma~\ref{lcp:lem:tail},
\[
 |E_3(q)|\le E_3(r)=\OO(t^6R(r)),\qquad
 |C_4(q)|\le C_4(r)=\OO(t^6R(r)+1).
\]
Since the cubic factor of $Q_3$ is bounded away from zero on the closed
unit disk, there are uniform estimates
\begin{align*}
 \int_{-\pi}^{\pi}\frac{d\theta}{|Q_3(r\ee^{i\theta})|^2}
     &=\OO(t^{-1}),\\
 \int_{-\pi}^{\pi}\frac{d\theta}{|Q_3(r\ee^{i\theta})|}
     &=\OO(\log(1/t)).
\end{align*}
The first follows from
$\int_{-\pi}^{\pi}|1+r\ee^{i\theta}|^{-2}d\theta
=2\pi/(1-r^2)$; the second follows by integrating
$(t^2+(\theta-\pi)^2)^{-1/2}$ near the single boundary zero.
The polynomials $h_3,U_3,P_3$ are uniformly bounded.

Apply Cauchy's coefficient formula to the three remainder terms in
\eqref{lcp:eq:cut} with $m=3$. Their total is
\[
 \OO\!\left(r^{-n}
     \{t^5R(r)+t^6R(r)\log(1/t)+\log(1/t)+1\}\right).
\]
The terms without $R(r)$ are exponentially smaller than the main
coefficient scale. Since the saddle estimate gives
$p_n\asymp r^{-n}R(r)t^{3/2}$, the dominant error is
$\OO(p_nt^{7/2})=\OO(p_nn^{-7/4})$.
\end{proof}

\subsection{Extracting the coefficients}
A finite rational expansion at $q=1$ gives
\[
 4F_3(\ee^{-z})=1-\frac12z^2+\frac12z^3+\OO(z^4).
\]
Together with \eqref{lcp:eq:Rlocal}, this yields
\begin{align}\label{lcp:eq:locals}
 (F_3R)(\ee^{-z})={}&
 \frac{z^{3/2}\ee^{K/z}}{4(2\pi)^{3/2}}
 \left(1-\frac z6-\frac{35z^2}{72}
                +\frac{755z^3}{1296}+\OO(z^4)\right),
\end{align}
up to a locally exponentially small modular remainder.
The rational factor has only fixed-order boundary poles, so the minor-arc
argument applies. Using \eqref{lcp:eq:transfer} through $\ell=3$ gives
\eqref{lcp:eq:d}--\eqref{lcp:eq:f}; its own truncation error is $\OO(n^{-2})$
relative to the leading term. The larger remainder in
Proposition~\ref{lcp:prop:approxerror} is $\OO(n^{-7/4})$, proving
Theorem~\ref{lcp:thm:main}.

For an explicit audit of the algebra, the relative $n^{-1}$ coefficient is
\[
 \underbrace{\frac{3}{4K}}_{j=0,k=2}
 +\underbrace{\frac12}_{j=1,k=1}
 -\underbrace{\frac{35K}{72}}_{j=2,k=0}.
\]
The relative $n^{-3/2}$ coefficient is
\[
 \underbrace{-\frac{5}{8\sqrt K}}_{j=1,k=2}
 +\underbrace{\frac{175\sqrt K}{72}}_{j=2,k=1}
 +\underbrace{\frac{755K^{3/2}}{1296}}_{j=3,k=0}.
\]
There is no $j=0,k=3$ contribution: the Bessel polynomial for that
half-integer order terminates after $k=2$.

Subtracting four times \eqref{lcp:eq:main} from \eqref{lcp:eq:pe2} proves
\eqref{lcp:eq:deficit}. Thus the lost proportion in the colored-partition
comparison is asymptotically $K/(2n)$, substantially smaller than the
explicit upper bound $3K/n$ supplied by the elementary squeeze.

\section{A proved all-orders radial expansion}
\subsection{A compatible family of rational approximants}
Define $D(q)=L(q)/R(q)$ for real $0<q<1$.
For every fixed $m\ge1$, \eqref{lcp:eq:Em} and \eqref{lcp:eq:tailbound} imply
\begin{equation}\label{lcp:eq:Dapprox}
 D(\ee^{-t})=\frac{h_m(\ee^{-t})}{Q_m(\ee^{-t})}
                       +\OO_m(t^{2m}).
\end{equation}
The denominator tends to $2^m$, so it is nonzero near the positive real
limit. The Taylor expansions of these rational functions are therefore
compatible: the coefficient of $t^j$ is fixed once $2m>j$.
Let $\mathcal D(t)$ be their unique formal limit.

\begin{theorem}[Radial expansion to arbitrary fixed order]\label{lcp:thm:radial}
The expansion \eqref{lcp:eq:radialintro} holds with
\begin{equation}\label{lcp:eq:Dformal}
 \mathcal F(t)=4\mathcal D(t)^2,\qquad
 \mathcal D(t)=\lim_{m\to\infty}^{\text{formal}}
         \frac{h_m(\ee^{-t})}{Q_m(\ee^{-t})}.
\end{equation}
For each fixed $m$,
\begin{equation}\label{lcp:eq:radialfinite}
 \frac{\Acal(\ee^{-t})}{R(\ee^{-t})}
   =\left(\frac{h_m(\ee^{-t})}{Q_m(\ee^{-t})}\right)^2
          +\OO_m(t^{2m}).
\end{equation}
The coefficients through $t^8$ are those in \eqref{lcp:eq:Fintro}.
\end{theorem}
\begin{proof}
Divide the finite-cut identity \eqref{lcp:eq:cut} by $R$.
For fixed $m$, the denominator $Q_m(\ee^{-t})$ stays bounded away from
zero. The $E_m/R$ and $C_{m+1}/R$ terms are $\OO_m(t^{2m})$ by
Lemma~\ref{lcp:lem:tail}; the remaining rational expression divided by $R$
is smaller than every power of $t$ by \eqref{lcp:eq:Rlocal}.
This proves \eqref{lcp:eq:radialfinite}. Compatibility and the product
expansion give \eqref{lcp:eq:Dformal}. Ordinary rational series arithmetic
then gives the displayed coefficients.
\end{proof}

For example, setting $w=1-q$ gives the useful intermediate expression
\[
 \mathcal D=
 \frac12-\frac18w^2-\frac5{32}w^4+\frac18w^5
 -\frac{65}{128}w^6+\frac{37}{32}w^7
 -\frac{2053}{512}w^8+\OO(w^9).
\]
Replacing $w$ by $1-\ee^{-t}$ and squaring reproduces
\eqref{lcp:eq:Fintro}. The accompanying symbolic checker verifies these
operations exactly over the rationals.

\subsection{Why this is not yet an all-orders coefficient theorem}
Increasing $m$ makes the positive radial tails smaller, but the complex
behavior of $1/Q_m$ must still be controlled. For general $m$, one cannot
simply assume that all zeros of $Q_m$ lie outside the unit disk.
No general zero classification is assumed or established here.
A meromorphic rational approximant can be useful on a contour, but poles
inside that contour require attention if one interprets its integral as
an ordinary Taylor coefficient.

The choice $m=3$ deliberately avoids these issues: the factorization in
\eqref{lcp:eq:thirdcut} proves the required zero exclusion exactly.
Consequently the first three correction terms are supported by a full
complex argument. Coefficients at later orders obtained by purely formal
transfer of \eqref{lcp:eq:radialintro} are candidates until an equally valid
remainder argument is supplied. No exponentially small contribution from
other roots of unity, nor a complete transseries with all such sectors,
is claimed in this report.

\begin{writenote}[Added 2 October 2026, batch~77P5: the gap is closed in Part~II]
Part~II proves the all-orders coefficient theorem
(Theorem~\ref{lcp:ao:thm:main}): for every fixed $M$, $a_n$ has an
expansion in powers of $(n-1/6)^{-1/2}$ with remainder
$\OO(n^{-M/2})$ and algorithmically determined coefficients. It does not
control the zeros of $Q_m$. It avoids them: the series $D=L/R$ of this
section is a partial theta function (Remark~\ref{lcp:ao:rem:LR}), and
$\Acal=R\,D^2$ plus a remainder whose coefficients are exponentially
smaller (Proposition~\ref{lcp:ao:prop:square},
Corollary~\ref{lcp:ao:cor:exponential}), with whole-circle bounds for both
factors. The later coefficients called ``candidates'' above are thereby
proved. The statement about exponentially small sectors and a complete
transseries still stands for both parts.
\end{writenote}

\begin{writenote}[Added 5 October 2026, batch~101]
Part~III proves the same theorem again (Theorem~\ref{lcp:tr:thm:main}),
also without controlling the zeros of $Q_m$: it uses the positivity of the
generating function off the real axis, a damping factor from the finite
product $(q;q)_N$, and an interpolation lemma that carries this section's
radial expansion into the saddle window (Section~\ref{lcp:tr:sec:transfer}).
That transfer uses only a complete real-axis expansion, analyticity and
the positive representation of Section~\ref{lcp:tr:sec:transfer}, so this
section's radial expansion, combined with it, already gives every order.
The zeros of $Q_m$ remain a separate question
(Section~\ref{lcp:tr:sec:further}, item~6).
\end{writenote}

\section{Inverse asymptotics and integer thresholds}
\subsection{A real inverse requires an interpolation convention}
The sequence itself is defined only on integers. To state a real inverse
accurately, choose an interpolation that preserves the proved asymptotic
accuracy. Put
\[
 M(x)=c\,x^{-3/2}\ee^{B\sqrt x}
       (1-dx^{-1/2}+ex^{-1}+fx^{-3/2})
\]
for sufficiently large real $x$, where the parenthesis is positive.
For an integer $k$, set $r_k=\log a_k-\log M(k)$.
Theorem~\ref{lcp:thm:main} gives $r_k=\OO(k^{-7/4})$.
For $x=k+\theta$, $0\le\theta\le1$, define
\begin{equation}\label{lcp:eq:interpolation}
 \log\widetilde a(x)=\log M(x)+(1-\theta)r_k+\theta r_{k+1}.
\end{equation}
This continuous interpolation agrees with $a_k$ at integers and obeys
\[
 \log\widetilde a(x)=\log M(x)+\OO(x^{-7/4}).
\]
It is eventually strictly increasing: the derivative of $\log M$ is
$B/(2\sqrt x)+\OO(x^{-1})$, while the piecewise-linear residual has
slope $\OO(x^{-7/4})$. Denote its eventual inverse by $\nu(y)$.
Other interpolations with the same asymptotic error give the same
expansion below to the stated accuracy.

This convention matters beyond low orders. Direct linear interpolation
of $\log a_k$, without first subtracting $\log M$, introduces a curvature
error of order $x^{-3/2}$, which can obscure the third logarithmic
correction. Equation~\eqref{lcp:eq:interpolation} avoids that ambiguity.

\subsection{Lambert \texorpdfstring{$W$}{W} and the correction terms}
Define
\begin{align}\label{lcp:eq:inverseconstants}
 D_0&=Bd=\frac{B^2}{12}+3,\notag\\
 \alpha_2&=B^2(e-d^2/2),\qquad
 \alpha_3=B^3(f+de-d^3/3),\\
 \ell&=\log\frac{y}{cB^3},\qquad
 s_0=-3W_{-1}\!\left(-\frac13\ee^{-\ell/3}\right).\notag
\end{align}
Here $W_{-1}$ is the real lower branch, which tends to $-\infty$ as its
argument tends to $0$ from below. It is selected because the inverse
index tends to $+\infty$.

\begin{theorem}[Asymptotic inverse]\label{lcp:thm:inverse}
For the interpolation \eqref{lcp:eq:interpolation}, as $y\to\infty$,
\begin{align}\label{lcp:eq:inverse}
 \nu(y)=\frac1{B^2}\bigg[
 s_0^2+2D_0+\frac{2(3D_0-\alpha_2)}{s_0}
 +\frac{18D_0-6\alpha_2-D_0^2-2\alpha_3}{s_0^2}
 \bigg]+\OO(s_0^{-5/2}).
\end{align}
In particular, a formula involving only logarithms is
\begin{equation}\label{lcp:eq:inverse-logs}
 \nu(y)=\frac{\ell^2+6\ell\log\ell+9(\log\ell)^2
                    +18\log\ell+2D_0}{B^2}
       +\OO\!\left(\frac{(\log\ell)^2}{\ell}\right).
\end{equation}
\end{theorem}
\begin{proof}
Set $s=B\sqrt{\nu(y)}$. Taking logarithms in the coefficient expansion
and using the interpolation error gives
\begin{equation}\label{lcp:eq:loginverse}
 \ell=s-3\log s-\frac{D_0}{s}
             +\frac{\alpha_2}{s^2}+\frac{\alpha_3}{s^3}
             +\OO(s^{-7/2}).
\end{equation}
The definition of $s_0$ is equivalent to
$s_0-3\log s_0=\ell$.
Put $s=s_0+\delta$ and compare successive inverse powers of $s_0$.
Since $1-3/s_0$ is bounded away from zero,
\[
 \delta=\frac{D_0}{s_0}
 +\frac{3D_0-\alpha_2}{s_0^2}
 +\frac{9D_0-3\alpha_2-D_0^2-\alpha_3}{s_0^3}
 +\OO(s_0^{-7/2}).
\]
Squaring and dividing by $B^2$ gives \eqref{lcp:eq:inverse}.
For \eqref{lcp:eq:inverse-logs}, direct substitution in
\eqref{lcp:eq:loginverse} yields
\[
 s=\ell+3\log\ell+\frac{9\log\ell+D_0}{\ell}
       +\OO((\log\ell)^2/\ell^2),
\]
and another squaring completes the calculation.
\end{proof}

If $N(y)=\min\{n\ge n_0:a_n\ge y\}$ for a sufficiently large fixed
$n_0$, then exactly
\[
 N(y)=\lceil\nu(y)\rceil.
\]
The error term belongs \emph{inside} that ceiling. Uniformly discarding
it can produce an incorrect integer when the approximation is very close
to a threshold. A safe coarser statement is that the right-hand side of
\eqref{lcp:eq:inverse-logs}, without its remainder, differs from $N(y)$ by
$\OO(1)$. Away from integers by more than a proved error bound, rounding
would be unambiguous; this article does not supply a numerical universal
constant for that error bound.

\begin{writenote}[Added 2 October 2026, batch~75: an instance of repository results]
The inversion of this section is an instance of the Lambert-centred
inversion apparatus of the transseries-and-inversion volume
\cite{lcp-tai}, and no novelty is claimed for the method. The core
$s_0-3\log s_0=\ell$ is \texttt{p0:thm:lambert-core} with $a=1$, $b=-3$;
its branch rule selects $W_{-1}$, as above. Equation
\eqref{lcp:eq:loginverse} is the core plus a power series in $1/s$, and the
correction $\delta$ is the all-orders reversion around the core,
\texttt{p0:thm:lambert-centered}, truncated at the order the forward
remainder allows. The shape is that of the partition-number chapter of the
same volume, \texttt{p3:sec:lambert} with \texttt{p3:thm:core-correction},
where the factor $n^{-1}$ in $p(n)$ gives $b=-2$; here $n^{-3/2}$ gives
$b=-3$. The rule $N(y)=\lceil\nu(y)\rceil$ and the warning that follows it
are part~(1) and part~(2) (the separation condition) of
\texttt{p0:thm:staircase}, specialized as in \texttt{p3:thm:staircase}. The
generic arithmetic of those two parts is formalized in Lean as
\texttt{Fabius.staircase\_ceil} and \texttt{Fabius.staircase\_separation}
in \nolinkurl{Analysis/FabiusFunction/Lean/FabiusFunction/StaircaseInversion.lean};
those statements concern an arbitrary monotone interpolation and give no
formal status to Theorem~\ref{lcp:thm:inverse}. What is specific to this
report is the forward expansion that feeds the inversion and the explicit
interpolation \eqref{lcp:eq:interpolation}.
\end{writenote}

\section{Reproducible exact and numerical checks}
\subsection{What was computed}
The attached \texttt{verify.py} uses only Python's standard library and
unbounded integers. It computes $a_0,\ldots,a_{2000}$ from the normalized
recurrence, and independently computes $p_n$ from the colored Euler
product. Division by $(1-q^m)^2$ is performed by two exact cumulative-sum
passes. Thus the computation avoids floating-point cancellation in the
original polynomial recurrence.

The checker confirms the initial 37 OEIS terms and four separately read
b-file anchors, at $n=50,100,200,500$ \cite{bfile}. It then verifies the
slope identity, the positive defect decomposition, both majorants in
\eqref{lcp:eq:majorants}, and the squeeze for every integer $3\le n\le2000$.
The third finite-cut identity, after clearing its polynomial denominator,
and positivity of its auxiliary tails are checked as well.
The output CSV contains all computed coefficients, not just a selected
plot. This is an exact finite test, not a formal proof of the infinite
statements.

\begin{writenote}[Added 5 October 2026, batch~101]
Since batch~101 all 2001 terms of the b-file have been compared: Part~III's
checker matches $a_0,\ldots,a_{2000}$ with a snapshot that is
byte-identical to the live b-file (checked at intake on 5~October 2026),
and this checker's coefficients agree with them
(Section~\ref{lcp:tr:sec:rechecked}). Run with \texttt{--nmax 8000} on a
copy at intake, it also passed all its checks, the squeeze included, through
$n=8000$, and reproduced Part~III's coefficients there.
\end{writenote}

The optional \texttt{verify\_symbolic.py} uses SymPy for rational symbolic
arithmetic and mpmath for high-precision numerical evaluations. It checks
Wronskians, the factorization of $Q_3$, the radial coefficients through
$t^8$, and the transfer coefficients $d,e,f$.
All included assertions passed in the supplied run.

\subsection{Accuracy of the asymptotic approximants}
Let
\[
 A_0(n)=cn^{-3/2}\ee^{B\sqrt n},\quad
 A_1=A_0(1-d/\sqrt n),\quad
 A_2=A_1+A_0e/n,\quad
 A_3=A_2+A_0f/n^{3/2}.
\]
The table reports $a_n/A_j(n)$; a ratio approaching $1$ is the desired
behavior. Rounded display values are not used by the checker.
\begin{center}
\begin{tabular}{r r r r r}
\toprule
$n$ & $a_n/A_0$ & $a_n/A_1$ & $a_n/A_2$ & $a_n/A_3$\\
\midrule
50   & 0.836729727 & 0.980057860 & 1.027184142 & 0.984290725\\
100  & 0.885790139 & 0.987955090 & 1.010017535 & 0.995767525\\
200  & 0.920485946 & 0.993103974 & 1.003708062 & 0.998875075\\
500  & 0.950776376 & 0.996878614 & 1.000989560 & 0.999807957\\
1000 & 0.965688969 & 0.998335826 & 1.000361244 & 0.999950135\\
2000 & 0.976025253 & 0.999128395 & 1.000130943 & 0.999987161\\
\bottomrule
\end{tabular}
\end{center}
More terms need not improve a short truncation at a modest index:
$n=50$ illustrates this. At $n=2000$, the four-term relative error is
approximately $1.284\times10^{-5}$.
The data also show
\[
 n(1-4a_n/p_n)=
 1.739699,\ 1.940375,\ 2.110009,\ 2.286402,\ 2.386995,\ 2.463821
\]
at the six listed indices, tending toward the proved limit
$K/2=2.673017858628\ldots$.

\subsection{Testing the inverse at known sequence values}
For $y=a_n$, the compatible inverse is exactly $n$.
The leading Lambert approximation is $s_0^2/B^2$; the corrected one
is the bracketed expression in \eqref{lcp:eq:inverse}.
\begin{center}
\begin{tabular}{r r r}
\toprule
$n$ & Leading inverse minus $n$ & Corrected inverse minus $n$\\
\midrule
50   & $-0.598588435$ & $-0.028604645$\\
100  & $-0.560169154$ & $-0.011441486$\\
200  & $-0.530798365$ & $-0.004479985$\\
500  & $-0.502617905$ & $-0.001257271$\\
1000 & $-0.487447241$ & $-0.000471103$\\
2000 & $-0.476247422$ & $-0.000174032$\\
\bottomrule
\end{tabular}
\end{center}
The nonvanishing limiting offset in the uncorrected Lambert approximation
is explained by the constant term $2D_0/B^2$ in the inverse theorem.
These values are numerical checks of already derived formulas.

\section{Further research questions and concrete next steps}
\begin{enumerate}[leftmargin=*,label=\textbf{\arabic*.},itemsep=0.65em]
\item \textbf{An all-orders coefficient expansion.}
Use the general finite-cut identity \eqref{lcp:eq:cut} to extend the
$\OO(n^{-7/4})$ remainder to arbitrary powers. The outstanding issue is
not generating formal coefficients; it is uniform contour control of
$Q_m^{-1}$ and treatment of possible interior poles of rational
approximants. A theorem organizing admissible cuts would turn the radial
algorithm into an all-orders coefficient algorithm.
\begin{writenote}[Added 2 October 2026, batch~77P5: answered]
Part~II proves the all-orders coefficient expansion
(Theorem~\ref{lcp:ao:thm:main}) by a different route, an exact
product-square decomposition, without organizing finite cuts. A theorem on
admissible cuts, or on the zeros of $Q_m$, remains a separate question.
\end{writenote}

\item \textbf{Explicit numerical error constants.}
Make the product minor-arc constants and the cubic lower bound in the
$m=3$ proof effective. This would give certified numerical intervals for
$a_n$ and safe integer inversion, rather than asymptotic $\OO$ statements
with unspecified constants.

\item \textbf{A combinatorial explanation of the factor four.}
Construct a direct relation between four copies of the L-convex class
and the colored partitions counted by $R$. The proven inequality
$p_n\ge4a_n$ for $n\ge1$ is suggestive, but an injection or decomposition
is not supplied by the analytic proof. A description of the defect class
could explain the limit $K/2$ combinatorially.

\item \textbf{The singularities of the boundary series.}
Study $D(q)=L/R$ as an analytic object, not just a radial formal series.
The recurrence, Wronskians, and positive tails give several exact
representations. Determine its domain of continuation and the role of
zeros of the finite denominators $Q_m$.
\begin{writenote}[Added 2 October 2026, batch~77P5: partly answered]
$D=L/R$ equals the partial theta function
$1-\sum_{k\ge0}q^{k(k+1)}(1-q^{2k+1})/(1+q^{2k+1})$ on the whole disk
$|q|<1$ (Remark~\ref{lcp:ao:rem:LR}, from Part~II's
Proposition~\ref{lcp:ao:prop:square}). Its continuation and the zeros of
$Q_m$ remain open; Part~II's open problem~5 continues the question.
\end{writenote}

\item \textbf{Beyond-all-orders terms.}
Separate exponentially small contributions to $\Acal(\ee^{-t})$ from
the modular product and from the recurrence boundary corrections.
Determine which roots of unity contribute to a coefficient transseries
and whether the radial series has a useful summability or resurgence
structure. None of these properties follows just from the computed
rational coefficients.
\begin{writenote}[Added 2 October 2026, batch~77P5]
Part~II separates an exact remainder whose coefficients are
$\OO(n^{1/4}\ee^{\pi\sqrt{2n}})$, exponentially smaller than $a_n$
(Corollary~\ref{lcp:ao:cor:exponential}); the bound is deliberately
coarse. The actual exponentially small contributions remain open (Part~II,
open problem~2).
\end{writenote}

\item \textbf{A reusable theorem for recurrence-defined area series.}
Abstract the positive-transfer and finite-cut mechanism to other
recurrences $y_{n+1}=a(q^n)y_n-b(q^n)y_{n-1}$ whose auxiliary generating
functions are Euler products. State conditions under which a positive
boundary defect is small enough for coefficient transfer. The present
proof gives a worked example with all analytic and positivity obligations
visible.

\item \textbf{Z-convex and higher-convexity area enumeration.}
Investigate whether the numerically studied Z-convex area asymptotics
\cite{zconvex} admit a related positive comparison. This is a proposed
application, not an assertion that the L-convex recurrence or constant
extends unchanged. Area and perimeter must remain distinct parameters.
\begin{writenote}[Added 2 October 2026, batch~77P5]
Part~II's open problem~4 asks the analogous question for its
product-square decomposition.
\end{writenote}

\begin{writenote}[Added 5 October 2026, batch~101]
The series data of~\cite{zconvex} are flagged as wrong by its first author
(the note in that entry and in Section~\ref{lcp:tr:sec:context}), so any
comparison would have to start from corrected Z-convex data.
\end{writenote}

\item \textbf{Formalization and OEIS integration.}
Formalize the normalized transfer, coefficientwise order, Wronskian, and
finite-cut identity first. These are algebraic and should be separated
from the eta transformation and complex saddle analysis. An OEIS update
should cite a reviewed version of the coefficient proof, rather than only
a numerical output or the all-orders radial statement.
\end{enumerate}

\section{Conclusion and claim ledger}
The principal result is a proof of a specific recorded conjecture, not
a conjecture inferred from additional terms. The established generating
function leads to a positive normalized recurrence; one auxiliary initial
condition leads exactly to a colored-partition Euler product. Two boundary
conditions produce the factor $1/4$. A positive coefficient squeeze proves
the leading asymptotic and its first correction. A finite cut at index
three, together with a full-circle error estimate, proves two further
corrections and a relative $\OO(n^{-7/4})$ remainder.

The following boundaries summarize what can be reused without confusion.
\begin{center}
\begin{tabular}{p{0.54\textwidth}p{0.36\textwidth}}
\toprule
Statement & Status in this manuscript\\
\midrule
Area generating function and its enumerative interpretation
& Cited input from prior work\\
Colored-partition product, positive squeeze, finite-cut identity
& Proved here from the recurrence\\
Leading coefficient asymptotic and three correction terms
& Proved with a complex remainder estimate\\
All-orders expansion for $\Acal(\ee^{-t})$, $t>0$
& Proved to every fixed radial order\\
Inverse expansion for a specified compatible interpolation
& Proved; integer rounding treated separately\\
All-orders coefficient expansion and exponentially small sectors
& Further research; not claimed\\
Finite checks through $n=2000$ and symbolic arithmetic
& Executed checks, not formal certification\\
Independent peer review or Lean verification
& Not performed in this report\\
\bottomrule
\end{tabular}
\end{center}
\begin{writenote}[Added 2 October 2026, batch~77P5]
The ledger is Part~I's. Its row ``All-orders coefficient expansion and
exponentially small sectors'' is now half answered: Part~II proves the
all-orders coefficient expansion (Theorem~\ref{lcp:ao:thm:main});
exponentially small sectors are still not claimed. The last two rows hold
for Part~II as well.
\end{writenote}

%% ====================================================================
%% Batch-77P5 addition (manuscript 31, "All order asymptotics and
%% inversion for L convex polyominoes").  All labels of Part II carry the
%% prefix lcp:ao: (all orders).  Section lcp:ao:sec:intro, the closing
%% section lcp:ao:sec:provenance and every text marked [write] were added
%% when Part II was written into the report; the other Part II sections
%% print the manuscript.
%% ====================================================================
\clearpage
\phantomsection\part*{Part II.\quad All-order asymptotics and inversion}
\addcontentsline{toc}{part}{Part II.\texorpdfstring{\quad}{ } All-order asymptotics and inversion}

\noindent\emph{Sections~1--12 and Appendices~A and~B form the original
report (Part~I) of 1~October 2026. Part~I is unchanged except for dated
\textsf{[write]} pointers to this addition (after the scope box, in
Section~1.1, at the end of Section~8.2, in research questions~1, 4, 5
and~7 of Section~11, and after the claim ledger of Section~12), a line in
the title, contents entries for the two parts, two macros in the preamble,
one bibliography entry used only by Part~II, and notes in two existing
entries. No number of Part~I moves.}

\section{Part II: source, scope and notation}\label{lcp:ao:sec:intro}
This section was written when Part~II was added. Sections~\ref{lcp:ao:sec:main}--\ref{lcp:ao:sec:open}
print the manuscript; Section~\ref{lcp:ao:sec:provenance} records the
provenance of both parts and the choices the merge made.

\subsection{Source and status}
Sections~\ref{lcp:ao:sec:main}--\ref{lcp:ao:sec:open} print, in full,
manuscript~31 of batch~77 of ProveIt's incoming-reports intake, titled
``All order asymptotics and inversion for L convex polyominoes'' and dated
1~October 2026. It arrived as the archive
\texttt{lconvex-\allowbreak asymptotics-\allowbreak reproducibility.zip}
in commit \texttt{096ee7b87} and was placed in commit \texttt{4f11bc9c0}.
Its author line is empty and it names no tool; like every manuscript of
the intake it is treated as AI-assisted research output. It names no
ProveIt commit and does not mention the repository, so it has no pin, and
it was not written against Part~I.

It is not a version of Part~I's manuscript. Only 40 of its 380 distinct
non-blank source lines occur in Part~I's source. Its proof architecture is
different: an exact product-square decomposition and a Hankel--Bessel
transfer, where Part~I uses a positive coefficient squeeze and finite cuts.
The proof note that its shipped audit reviewed is a different file from
Part~I's manuscript. The package ships three internal reviews
(\texttt{02-all-orders-mathematical-audit.md},
\texttt{02-all-orders-final-source-review.md},
\texttt{02-all-orders-computational-review.md}). They are the delivery's
own checks, not external peer review. Part~II is unrefereed. Nothing in
it is formalized in Lean or Rocq, and its place in the collection confers
no formal status on it.

\paragraph{What was rechecked at intake and in this write.} Independently
of the delivery:
\begin{itemize}[leftmargin=*,itemsep=0.15em]
\item $a_0,\ldots,a_{120}$ from the original form \eqref{lcp:ao:eq:known}
 agree with the shipped coefficient file, and that file agrees entry by
 entry with the column $a_n$ of Part~I's \texttt{data/coefficients.csv}
 for all $0\le n\le2000$;
\item $c_0,\ldots,c_4$ follow from \eqref{lcp:ao:eq:cr} with
 $b_0,\ldots,b_4$ read from \eqref{lcp:ao:eq:Dexp}, the $h_r$ of
 Section~\ref{lcp:ao:sec:inverse} are $(2\sqrt\kappa)^rc_r$, and the
 unshifted re-expansion gives exactly Part~I's $d,e,f$
 (Remark~\ref{lcp:ao:rem:partI}), all with exact SymPy arithmetic;
\item $B=D$ and $A=PB^2+R$, with $R$ from the literal $E(q)$, at
 $q=0.3$, $0.6$ and $0.3+0.4\ii$, to 40 digits (discrepancies
 $10^{-38}$ to $10^{-41}$);
\item the identity of Remark~\ref{lcp:ao:rem:LR} between Part~I's
 $L/R^{\mathrm I}$ and Part~II's $B$, at the same three points;
\item the entries of Table~\ref{lcp:ao:tab:errors}, the leading constant,
 the monotonicity injection (by hand), and the log-concavity expansion
 numerically (Corollary~\ref{lcp:ao:cor:logconcavity});
\item the shipped replay, core and optional parts
 (Section~\ref{lcp:ao:sec:checks}).
\end{itemize}
The whole-circle estimates (Lemmas~\ref{lcp:ao:lem:global}
and~\ref{lcp:ao:lem:minor}) and the contour deformation in
Section~\ref{lcp:ao:sec:extraction} were read but not re-derived in
detail; they are written at the level of standard saddle-point estimates.

\subsection{What Part II adds, and what stays in Part I}
Part~II proves, from the same cited area generating function:
\begin{itemize}[leftmargin=*,itemsep=0.15em]
\item the exact decomposition $A=PBD+R$ and the identity $B=D$, a partial
 theta function, hence $A=PB^2+R$ exactly
 (Propositions~\ref{lcp:ao:prop:split} and~\ref{lcp:ao:prop:square});
\item whole-circle bounds and an exponentially smaller exact remainder
 (Lemma~\ref{lcp:ao:lem:global}, Corollary~\ref{lcp:ao:cor:exponential});
\item \textbf{the all-orders coefficient theorem}
 (Theorem~\ref{lcp:ao:thm:main}) in the shifted index $N=n-1/6$, with
 coefficients given by a finite algorithm. This closes the gap that
 Part~I states in Section~8.2 and research question~1 of Section~11;
\item eventual strict log-concavity with a two-term expansion of
 $\log(a_n^2/(a_{n-1}a_{n+1}))$ (Corollary~\ref{lcp:ao:cor:logconcavity});
\item an injection proving that $a_n$ is nondecreasing;
\item a Lambert-$W$ inverse in shifted coordinates, logarithmic inverse
 polynomials, and shrinking integer brackets for every order
 (Section~\ref{lcp:ao:sec:inverse}).
\end{itemize}
Only Part~I has the coefficientwise squeeze $0\le p_n-4a_n\le3j_n$ for
every $n\ge3$ with its colored-partition reading, the deficit limit $K/2$,
the finite-cut identity and its $m=3$ route, the interpolation $\nu(y)$
with the exact rule $N(y)=\lceil\nu(y)\rceil$, and the tables of
Section~10.

Several results are proved in both parts, by different arguments; each
part keeps its own proof.
\begin{itemize}[leftmargin=*,itemsep=0.15em]
\item The leading asymptotic and the corrections $d,e,f$ of
 Theorem~\ref{lcp:thm:main}: Part~I by squeeze and finite cut, with
 remainder $\OO(n^{-7/4})$; Part~II as the case $M=4$ of
 Theorem~\ref{lcp:ao:thm:main}, with remainder $\OO(n^{-2})$
 (Remark~\ref{lcp:ao:rem:partI}).
\item The radial expansion \eqref{lcp:eq:radialintro} with the
 coefficients \eqref{lcp:eq:Fintro} of Theorem~\ref{lcp:thm:radial}: it
 follows from \eqref{lcp:ao:eq:exactmain}, and the coefficients are four
 times those of $D(\ee^{-z})^2$ in \eqref{lcp:ao:eq:Dexp}
 (Remark~\ref{lcp:ao:rem:LR}).
\item The local eta expansion \eqref{lcp:eq:Rlocal} is
 \eqref{lcp:ao:eq:Plocal}; the minor-arc bound \eqref{lcp:eq:minor} is
 Lemma~\ref{lcp:ao:lem:minor}, for the same product; the transfer
 \eqref{lcp:eq:transfer}--\eqref{lcp:eq:betak} is
 \eqref{lcp:ao:eq:Bessel}--\eqref{lcp:ao:eq:cr} in the shifted index.
\item The generating function \eqref{lcp:eq:enum} is
 \eqref{lcp:ao:eq:known}, with Part~I's $h_n$ equal to the manuscript's
 $f_{n-1}$; the initial terms are printed in both parts.
\end{itemize}

\subsection{Notation across the two parts}\label{lcp:ao:sec:notation}
Part~II keeps the manuscript's letters. The shipped data files and reviews
use them (for example the keys \texttt{A\_equals\_P\_B\_squared\_plus\_R}
and \texttt{initial\_R} of \texttt{02-all-orders-exact-identities.json},
and \texttt{kappa} in the forward and inverse series), and renaming would
break that correspondence. Many of these letters mean something else in
Part~I. In Sections~\ref{lcp:ao:sec:main}--\ref{lcp:ao:sec:open} a letter
has its Part~II meaning. Where a \textsf{[write]} note in Part~II refers to
a Part~I object whose letter clashes, it adds the superscript
$\mathrm I$: $R^{\mathrm I}$, $K^{\mathrm I}$, $B^{\mathrm I}$,
$D^{\mathrm I}$, $L^{\mathrm I}$, $y^{\mathrm I}$. No symbol of either
manuscript was renamed and no normalization was changed.
Table~\ref{lcp:ao:tab:notation} lists every clash. The last column gives
the tempting false reading.

{\small
\begin{longtable}{@{}>{\raggedright}p{0.1\textwidth}>{\raggedright}p{0.3\textwidth}>{\raggedright}p{0.22\textwidth}>{\raggedright\arraybackslash}p{0.28\textwidth}@{}}
\caption{Letters of Part~II, the same object in Part~I, and the clash to avoid.}\label{lcp:ao:tab:notation}\\
\toprule
Part II & Meaning in Part II & Same object in Part I & Not to be confused with\\
\midrule
\endfirsthead
\toprule
Part II & Meaning in Part II & Same object in Part I & Not to be confused with\\
\midrule
\endhead
\bottomrule
\endfoot
$\kappa$ & $13\pi^2/24$ & $K$, \eqref{lcp:eq:constants} & Part~I's $\kappa$ in \eqref{lcp:eq:Rlocal}, an unspecified sector-dependent exponent\\
$C$ & $13\sqrt2/768$ & $c$ & Part~I's series $C(z;q)$, $C_j$, \eqref{lcp:eq:C}\\
$N$ & $n-1/6$ & --- & Part~I's threshold $N(y)$, and its block length $N\asymp1/t$ in the minor-arc proof\\
$P(q)$ & $(-q;q^2)_\infty/(q;q)_\infty^3$, \eqref{lcp:ao:eq:PBD}, \eqref{lcp:ao:eq:eta} & $R(q)$, \eqref{lcp:eq:R} & Part~I's polynomials $P_m$; Part~II's $P_j$\\
$R(q)$ & exact remainder $1+E(q)/(q;q)_\infty$, Proposition~\ref{lcp:ao:prop:split} & none & Part~I's $R(q)$, which is Part~II's $P$: $A=PB^2+R$ does \emph{not} say $\Acal=R^{\mathrm I}D^2+R^{\mathrm I}$\\
$p_n$ & $[q^n]P(q)B(q)^2$, Corollary~\ref{lcp:ao:cor:exponential}; $p_n\sim a_n$ & none & Part~I's $p_n=[q^n]R^{\mathrm I}(q)=[q^n]P(q)$, with $p_n\sim4a_n$\\
$E(q)$ & the double sum \eqref{lcp:ao:eq:E} & none & Part~I's defect $E=R-2L$, \eqref{lcp:eq:E}, and tails $E_m$, \eqref{lcp:eq:Em}\\
$B(q)$ & $1-\sum_lu_l$, the partial theta function \eqref{lcp:ao:eq:theta} & $D(q)=L/R$ of Section~8.1, exactly (Remark~\ref{lcp:ao:rem:LR}) & Part~I's constant $B=2\sqrt K=\pi\sqrt{13/6}$\\
$D(q)$ & the series of \eqref{lcp:ao:eq:PBD}; $D=B$ & $D(q)=L/R$: the same function, not only the same letter & Part~II's bracket constants $D_M$\\
$K$ & $C(2\sqrt\kappa)^3$, \eqref{lcp:ao:eq:inversevars} & $cB^3$ & Part~I's $K=13\pi^2/24$, which is Part~II's $\kappa$\\
$L$, $w$ & $\log(X/K)$, $\log L$ & $\ell=\log(y/(cB^3))$ of \eqref{lcp:eq:inverseconstants} & Part~I's slope series $L=\Lcal_h$, \eqref{lcp:eq:slope}\\
$X$ & target value of the inverse & $y$ of Section~9 & ---\\
$y$ & $2\sqrt{\kappa(n-1/6)}=B^{\mathrm I}\sqrt N$ & unshifted analogue $s=B\sqrt\nu$ & Part~I's $y$, the target value\\
$y_0$ & Lambert root \eqref{lcp:ao:eq:W} & $s_0$ of \eqref{lcp:eq:inverseconstants}, the same number when $X=y^{\mathrm I}$ & ---\\
$n_*(X)$ & $\min\{n:a_n\ge X\}$ & $N(y)$ of Section~9 & ---\\
$F_M$, $x_M$ & real model \eqref{lcp:ao:eq:model} and its large inverse & analogues of $M(x)$ and $\nu(y)$ & Part~I's rational $F_3(q)$, \eqref{lcp:eq:thirdcut}, and series $\mathcal F(t)$\\
$M$ & truncation order & --- & Part~I's model function $M(x)$\\
$h_n$, $H$, $G$ & $f_{n-1}/(q;q)_n$, its generating function, the product \eqref{lcp:ao:eq:HG} & Part~I's $h_n=f_{n-1}$ lacks the factor $1/(q;q)_n$; $G$ is Part~I's $G$, \eqref{lcp:eq:Gproduct} & $h_r$ of Section~\ref{lcp:ao:sec:inverse}\\
$h_r$, $\ell_r$ & coefficients of \eqref{lcp:ao:eq:yforward} and of its logarithm, \eqref{lcp:ao:eq:ell} & $D_0$, $\alpha_2$, $\alpha_3$ play this role for $s$ & $h_n$ above; Part~I's $\ell$\\
$b_m$ & Taylor coefficients of $D(\ee^{-z})^2$, \eqref{lcp:ao:eq:bdef} & $\mathcal F(t)=4\sum b_mt^m$, \eqref{lcp:eq:Dformal} & Part~I's $b_n=y_n/(q;q)_n^2$ and color multiplicities $b_j$, \eqref{lcp:eq:colors}\\
$c_r$ & coefficients of \eqref{lcp:ao:eq:allorders} & unshifted: $1,-d,e,f$ & Part~I's $c$ and $c_0$; the local constants $c$, $c_1,\ldots,c_4$, $C_1$ of Part~II's proofs\\
$u_l$, $T_j$ & summands of $B$ in \eqref{lcp:ao:eq:PBD}; summands of $D$ (proof of Lemma~\ref{lcp:ao:lem:localD}) & --- & Part~I's $u_n=Q_n-h_n$ and $T_n(q)$\\
$P_j$ & logarithmic inverse polynomials & --- & $P(q)$; Part~I's $P_m$\\
$a$, $b$, $S$, $V$ & local to the proof of Proposition~\ref{lcp:ao:prop:square} ($a=\ii\sqrt q$) & --- & $a_n$; $b_m$; Part~I's $s$\\
$z$, $t$ & auxiliary variable in Section~\ref{lcp:ao:sec:decomposition}; from Section~\ref{lcp:ao:sec:local} on $q=\ee^{-z}$, $t=\Re z$ & same two uses in Part~I & saddle radius $t=\sqrt{\kappa/N}$ here, $t=\sqrt{K/n}$ in Part~I\\
$\lambda$, $\alpha$ & $\sqrt{\kappa N}$; $3/2+m$ & --- & Part~I's amplitude $\lambda$ in \eqref{lcp:eq:transfer}; $\alpha_2$, $\alpha_3$\\
$d_r$ & local to the proof of Corollary~\ref{lcp:ao:cor:logconcavity} & --- & Part~I's $d$\\
\end{longtable}}

Two further letters are shared with the same meaning: $f_n$ (the
recurrence \eqref{lcp:ao:eq:rec}; Part~I's $h_{n+1}$) and the branch
$W_{-1}$.

\subsection{The manuscript's abstract}
\begin{quote}\small
Let $a_n$ count fixed $L$-convex polyominoes of area $n$, with the empty object included. Starting from the known exact area generating function, we prove the asymptotic equivalent conjectured by Guttmann and Kot\v{e}\v{s}ovec and recorded in OEIS A126764. An exact decomposition into an eta quotient times the square of a partial theta function, plus an exponentially smaller coefficient remainder, supplies every fixed algebraic order. The correction coefficients are given by a finite algorithm. We also give a branch-explicit Lambert $W$ inverse, logarithmic inverse polynomials, and shrinking intervals for the integer threshold inverse before rounding. The proof includes whole-circle estimates; a positive-real singular expansion alone is not used as a coefficient-transfer argument. Exact coefficient checks and a portable replay accompany the analysis.
\end{quote}

\section{The sequence and the main theorem}\label{lcp:ao:sec:main}
An $L$-convex polyomino is a convex polyomino in which any two cells can be joined internally by a path with at most one right-angle turn. We count fixed polyominoes, identifying translations but keeping rotations and reflections distinct. Its area generating function is
\begin{equation}\label{lcp:ao:eq:known}
 A(q)=\sum_{n\ge0}a_nq^n
 =1+\sum_{n\ge1}\frac{q^n f_{n-1}(q)}{(q;q)_{n-1}^2(1-q^n)},
\end{equation}
where
\begin{equation}\label{lcp:ao:eq:rec}
 f_{-1}=f_0=1,\qquad f_n=2f_{n-1}-(1-q^n)^2f_{n-2}\quad(n\ge1).
\end{equation}
Here $(a;q)_j=\prod_{h=0}^{j-1}(1-aq^h)$. Formula \eqref{lcp:ao:eq:known} is the Castiglione--Frosini--Munarini--Restivo--Rinaldi area formula~\cite{castiglione}, in the normalization reproduced in~\cite[Eq.~(1.1)]{gk}. The first terms are
\[
1,1,2,6,15,35,76,156,310,590,1098,1984,3515,6094,\ldots.
\]
Set
\begin{equation}\label{lcp:ao:eq:constants}
 \kappa=\frac{13\pi^2}{24},\qquad C=\frac{13\sqrt2}{768},\qquad N=n-\frac16.
\end{equation}

\begin{theorem}\label{lcp:ao:thm:main}
For every fixed integer $M\ge1$,
\begin{equation}\label{lcp:ao:eq:allorders}
 a_n=C N^{-3/2}\ee^{2\sqrt{\kappa N}}
 \left(\sum_{r=0}^{M-1}c_r N^{-r/2}+O(N^{-M/2})\right).
\end{equation}
The coefficients are algorithmically determined as follows. Define the formal asymptotic coefficients $b_m$ by
\begin{equation}\label{lcp:ao:eq:bdef}
 \left(\sum_{j\ge0}\frac{(-1)^j\ee^{-zj(j-1)/2}(\ee^{-z};\ee^{-z})_j}
 {(-\ee^{-z};\ee^{-2z})_{j+1}}\right)^2
 \sim\sum_{m\ge0}b_m z^m.
\end{equation}
To determine $b_0,\ldots,b_{M-1}$, it suffices to retain $0\le j<M$ on the left and take ordinary Taylor coefficients. Then
\begin{equation}\label{lcp:ao:eq:cr}
 c_r=\sum_{\substack{m+k=r\\0\le k\le m+2}}
 4b_m\kappa^{m/2}\frac{(-1)^k(m+2+k)!}
 {(m+2-k)!\,k!\,(4\sqrt\kappa)^k}.
\end{equation}
In particular,
\begin{align}\label{lcp:ao:eq:firstc}
 c_0&=1,&c_1&=-\frac{3}{2\sqrt\kappa},&
 c_2&=\frac{3}{4\kappa}-\frac\kappa2,\notag\\
 c_3&=\frac{\sqrt\kappa(\kappa+5)}2,&
 c_4&=-\frac{45}{8}-\frac{15\kappa}{4}-\frac{41\kappa^2}{48}.
\end{align}
\end{theorem}
The $M=1$ case proves the OEIS conjecture
\begin{equation}\label{lcp:ao:eq:leading}
 a_n\sim \frac{13\sqrt2}{768}\,n^{-3/2}\exp\!\left(\pi\sqrt{\frac{13n}{6}}\right).
\end{equation}
For readers preferring the unshifted variable, the first correction is
\begin{equation}\label{lcp:ao:eq:unshifted}
 a_n=Cn^{-3/2}\ee^{2\sqrt{\kappa n}}
 \left[1-\left(\frac{\sqrt\kappa}{6}+\frac{3}{2\sqrt\kappa}\right)n^{-1/2}
 +O(n^{-1})\right].
\end{equation}
The shift $N=n-1/6$ keeps the higher formulas shorter.

\begin{remark}[write, 2 October 2026: consistency with Part~I]\label{lcp:ao:rem:partI}
Take $M=4$ in Theorem~\ref{lcp:ao:thm:main} and re-expand in the unshifted
index. With $\kappa=K^{\mathrm I}$ this gives
\[
 a_n=Cn^{-3/2}\ee^{2\sqrt{\kappa n}}
 \left(1-\frac d{\sqrt n}+\frac en+\frac f{n^{3/2}}+\OO(n^{-2})\right),
\]
with exactly the $d,e,f$ of \eqref{lcp:eq:d}--\eqref{lcp:eq:f}: the three
symbolic differences vanish (exact SymPy arithmetic, at intake and again in
this write). So Theorem~\ref{lcp:ao:thm:main} re-proves
Theorem~\ref{lcp:thm:main} with the sharper relative remainder
$\OO(n^{-2})$, by a route independent of Part~I's squeeze and finite cut,
and continues it to every order. Equation \eqref{lcp:ao:eq:unshifted} is
Part~I's first correction $-d$, and \eqref{lcp:ao:eq:leading} is the
leading asymptotic proved in both parts. Part~I keeps its own statement
and proof.
\end{remark}

\paragraph{Literature status and prefactor.}
The inspected OEIS entry~\cite{oeis} explicitly labels \eqref{lcp:ao:eq:leading} conjectural. The published 2023 paper~\cite{gk} describes its coefficient analysis as conjectural, and the 2024 paper~\cite[p.~19]{zconvex} repeats that status. The introduction of~\cite{gk} and its quoted numerical value $0.023938510821419\ldots$ agree with $C$ above, while its final display in Section~2 prints the reciprocal. Our derivation fixes the normalization independently. Targeted searches through 1 October 2026 did not locate an earlier proof or an all-orders expansion; this is a scoped search result, not an exhaustive claim of priority.

\begin{writenote}[Added 2 October 2026, batch~77P5]
The statement that the final display of Section~2 of~\cite{gk} prints the
reciprocal of $C$ was \emph{not} checked at intake. The shipped
mathematical audit (item~13) says the same, but it is the delivery's own
review. Part~I cites~\cite{gk} for the conjecture and its equation~(1.1)
only. As in Part~I, the search statement is scoped, and no priority is
claimed. Part~I's own proof of the leading term and three corrections
(Theorem~\ref{lcp:thm:main}) predates this manuscript in the repository,
though the two were written independently.
\end{writenote}

\begin{writenote}[Added 5 October 2026, batch~101]
Now checked: at intake on 5~October 2026 the final display of Section~2
of~\cite{gk} (arXiv v3) was found to have the amplitude $1/C$, while the
decimal constant quoted there is $C$, as Part~II says. See the note in
Section~\ref{lcp:tr:sec:main}.
\end{writenote}

\section{An exact decomposition of the generating function}\label{lcp:ao:sec:decomposition}
All $q$-series in this section converge for $|q|<1$. Introduce
\begin{equation}\label{lcp:ao:eq:HG}
 h_n=\frac{f_{n-1}}{(q;q)_n},\quad H(z)=\sum_{n\ge0}h_nz^n\quad(|z|<1),\quad
 G(z)=\frac{(-qz^2;q^2)_\infty}{(z;q)_\infty^2}.
\end{equation}
The letter $z$ here is an auxiliary generating-function variable; later $q=\ee^{-z}$ will be explicitly stated.

\begin{lemma}\label{lcp:ao:lem:H}
The following identity holds for $|z|<1$:
\begin{equation}\label{lcp:ao:eq:Hsolution}
 H(z)=G(z)\left[1-\sum_{j\ge0}\frac{q^jz(q^{j+1}z;q)_\infty^2}
 {(-q^{2j+1}z^2;q^2)_\infty}\right].
\end{equation}
\end{lemma}
\begin{proof}
The recurrence gives
\[
(1-z)^2H(z)=(1+qz^2)H(qz)-z,
\]
whereas $G$ satisfies the corresponding homogeneous equation. Division by $G$ and iteration toward $q^jz\to0$ give \eqref{lcp:ao:eq:Hsolution}. To justify convergence directly, put $\epsilon_k=2q^k-q^{2k}$. First differences in \eqref{lcp:ao:eq:rec} yield
\[
f_n=1+\sum_{k=1}^n(n-k+1)\epsilon_k f_{k-2}.
\]
Divide by $n+2$ and apply the discrete Gronwall inequality; the relevant product converges because $\sum k|\epsilon_k|<\infty$. Thus $f_n=O_q(n)$, uniformly on compact $q$-subdisks. Since $(q;q)_n$ has a nonzero limit, $H$ is analytic for $|z|<1$. All subsequent rearrangements are absolutely convergent, initially on compact subdisks.
\end{proof}
Define
\begin{align}\label{lcp:ao:eq:PBD}
 P(q)&=\frac{(-q;q^2)_\infty}{(q;q)_\infty^3},\notag\\
 u_l(q)&=\frac{q^l(q^{l+1};q)_\infty^2}{(-q^{2l+1};q^2)_\infty},
 &B(q)&=1-\sum_{l\ge0}u_l(q),\notag\\
 D(q)&=\sum_{j\ge0}\frac{(-1)^j q^{j(j-1)/2}(q;q)_j}{(-q;q^2)_{j+1}}.
\end{align}
Formula \eqref{lcp:ao:eq:Hsolution} gives
\begin{equation}\label{lcp:ao:eq:Hqj}
 H(q^{j+1})=G(q^{j+1})\left(B+\sum_{l=0}^j u_l\right).
\end{equation}

\begin{proposition}\label{lcp:ao:prop:split}
There is an exact identity
\begin{equation}\label{lcp:ao:eq:splitBD}
 A(q)=P(q)B(q)D(q)+R(q),
\end{equation}
where $R(q)=1+(q;q)_\infty^{-1}E(q)$ and
\begin{equation}\label{lcp:ao:eq:E}
 E(q)=\sum_{j\ge0}\frac{(-1)^j q^{j(j-1)/2}}{(q;q)_j}
 \sum_{l=0}^j
 \frac{q^l(q^{l+1};q)_{j-l}^2}{(-q^{2l+1};q^2)_{j-l+1}}.
\end{equation}
\end{proposition}
\begin{proof}
Euler's $q$-binomial theorem gives
\[
\frac1{(q;q)_{n-1}}=\frac{(q^n;q)_\infty}{(q;q)_\infty}
=\frac1{(q;q)_\infty}\sum_{j\ge0}
\frac{(-1)^jq^{j(j-1)/2+nj}}{(q;q)_j}.
\]
Substitute this in $A-1=\sum_{n\ge1}q^nh_n/(q;q)_{n-1}$ and interchange sums. The constant piece vanishes because
\[
\sum_{j\ge0}\frac{(-1)^jq^{j(j-1)/2}}{(q;q)_j}=(1;q)_\infty=0.
\]
For the homogeneous part of \eqref{lcp:ao:eq:Hqj}, use
\[
G(q^{j+1})=\frac{(-q;q^2)_\infty}{(q;q)_\infty^2}
\frac{(q;q)_j^2}{(-q;q^2)_{j+1}}.
\]
The remaining product simplifies, for $l\le j$, to
\[
G(q^{j+1})u_l=
\frac{q^l(q^{l+1};q)_{j-l}^2}{(-q^{2l+1};q^2)_{j-l+1}},
\]
which proves the proposition.
\end{proof}

\section{The square identity and a partial theta representation}\label{lcp:ao:sec:square}
\begin{proposition}\label{lcp:ao:prop:square}
One has
\begin{equation}\label{lcp:ao:eq:theta}
 B(q)=D(q)=1-\sum_{k\ge0}q^{k(k+1)}\frac{1-q^{2k+1}}{1+q^{2k+1}}.
\end{equation}
Consequently
\begin{equation}\label{lcp:ao:eq:exactmain}
 A(q)=P(q)B(q)^2+R(q).
\end{equation}
\end{proposition}
\begin{proof}
Set $a=\ii\sqrt q$, $b=-\ii\sqrt q$, so $a+b=0$, $ab=q$. Both choices of square root give the same symmetric expressions. We use the convention
\[
{}_2\phi_1(a,b;c;q,z)=\sum_{n\ge0}\frac{(a,b;q)_n}{(c,q;q)_n}z^n.
\]
For the transformation argument assume $q\ne0$; the identities at $q=0$ follow by continuity.
Heine's first transformation~\cite[17.6.6]{lcp-ao-dlmf} applied to ${}_2\phi_1(a,b;q;q,q)$ gives
\begin{equation}\label{lcp:ao:eq:Lambert}
 S:=1-B=\sum_{n\ge0}\frac{b^n}{1-aq^n}
 =\sum_{n,k\ge0}a^k b^nq^{nk}.
\end{equation}
Indeed, the defining sum of the $u_l$ is
\[
\frac{(q;q)_\infty^2}{(a,b;q)_\infty}\,{}_2\phi_1(a,b;q;q,q).
\]
After Heine's transformation, the series ratio is $(a;q)_n/(aq;q)_n=(1-a)/(1-aq^n)$; the product prefactor cancels with $1-a$. Formula \eqref{lcp:ao:eq:Lambert} is symmetric in $a,b$.

Let $V=(1+q)D$ and let $V_q$ denote the same summation with an extra factor $q^n$ in term $n$. Fine's second transformation~\cite[17.6.10]{lcp-ao-dlmf}, in the form
\[
(1-z){}_2\phi_1(q,\alpha q;\beta q;q,z)
=\sum_{n\ge0}\frac{(\beta/\alpha;q)_n(-\alpha z)^nq^{n(n+1)/2}}
 {(\beta q,zq;q)_n},
\]
with $\alpha=a/q^2$, $\beta=a/q$, $z=b$, yields
\begin{equation}\label{lcp:ao:eq:Fine}
\frac{V-aV_q}{1-a}=(1-b){}_2\phi_1(q,a/q;a;q,b)
=(1-b)\bigl[1+(b-1)S\bigr].
\end{equation}
For the second equality, split off $n=0$ and use
\[
\frac{(a/q;q)_n}{(a;q)_n}=\frac{1-a/q}{1-aq^{n-1}},\qquad ab=q.
\]
Swap $a,b$ in \eqref{lcp:ao:eq:Fine}. Multiplying the original and swapped equations by $1-a$ and $1-b$ and adding gives $2V=2(1+q)(1-S)$. Hence $D=B$.

Finally split the double sum \eqref{lcp:ao:eq:Lambert} into its diagonal and two off-diagonal triangles. On a diagonal difference $d$, $a^d+b^d$ vanishes for odd $d$, while it is $2(-1)^m q^m$ for $d=2m$. Therefore
\[
S=\sum_{k\ge0}q^{k(k+1)}\left(1+2\sum_{m\ge1}(-1)^mq^{m(2k+1)}\right),
\]
and summing the geometric progression gives \eqref{lcp:ao:eq:theta}.
\end{proof}

\begin{remark}[write, 2 October 2026: Part~I's boundary series is the partial theta function]\label{lcp:ao:rem:LR}
For $|q|<1$, Part~I's series $D^{\mathrm I}(q)=L^{\mathrm I}(q)/R^{\mathrm I}(q)$
of Section~8.1, with $L^{\mathrm I}$ the normalized limiting slope
\eqref{lcp:eq:slope} of the solution $h^{\mathrm I}_n=f_{n-1}$, equals
$B(q)=D(q)$. Indeed, in \eqref{lcp:ao:eq:Hsolution} the bracket is
analytic for $|z|<|q|^{-1/2}$ (the zeros of $(-q^{2j+1}z^2;q^2)_\infty$
have $|z|\ge|q|^{-1/2}$, and the factor $q^j$ makes the sum converge) and
equals $B(q)$ at $z=1$, while $G$ has in $|z|<|q|^{-1}$ only the double
pole at $z=1$, with $(1-z)^2G(z)\to(-q;q^2)_\infty/(q;q)_\infty^2$. Hence
\[
 \frac{f_{n-1}}{(q;q)_n}=[z^n]H(z)
 =\frac{(-q;q^2)_\infty B(q)}{(q;q)_\infty^2}\,n+\OO_q(1),
\]
so the slope of Lemma~\ref{lcp:lem:slope} is
$s(q)=(-q;q^2)_\infty B(q)/(q;q)_\infty$, and
$L^{\mathrm I}=s/(q;q)_\infty^2=R^{\mathrm I}B$. A 40-digit evaluation at
$q=0.3$, $0.6$ and $0.3+0.4\ii$ agrees to $10^{-37}$. Consequences: Part~I's
formal limit $\mathcal D(t)$ in \eqref{lcp:eq:Dformal} is the Poincar\'e
expansion of $B(\ee^{-t})$ given by Lemma~\ref{lcp:ao:lem:localD}, and the
coefficients \eqref{lcp:eq:Fintro} are four times those of $D(\ee^{-z})^2$
in \eqref{lcp:ao:eq:Dexp}; Part~I's rational approximants $h_m/Q_m$ in
\eqref{lcp:eq:Dapprox} approximate this partial theta function; and
Theorem~\ref{lcp:thm:radial} also follows from \eqref{lcp:ao:eq:exactmain},
since $P=R^{\mathrm I}$ and $R$ is exponentially smaller on the radius by
Lemma~\ref{lcp:ao:lem:global}. This answers part of Part~I's research
question~4 (Section~11): $D^{\mathrm I}$ has the closed form
\eqref{lcp:ao:eq:theta} on the whole unit disk. Its continuation beyond
$|q|<1$ and the zeros of the finite denominators $Q_m$ are not addressed.
\end{remark}

\section{Global estimates and the exponentially smaller part}\label{lcp:ao:sec:global}
These bounds hold on the entire coefficient-extraction circle, not only on the positive real radius. Let $r=\ee^{-t}$, $t>0$, and $|q|=r$.

\begin{lemma}\label{lcp:ao:lem:global}
As $t\downarrow0$, uniformly on $|q|=r$,
\begin{align}\label{lcp:ao:eq:globalbounds}
 |B(q)|&=O\!\left(t^{-1}\log\frac1t\right),\notag\\
 |R(q)|&=O\!\left(t^{-1/2}\ee^{\pi^2/(2t)}\right).
\end{align}
\end{lemma}
\begin{proof}
From \eqref{lcp:ao:eq:theta},
\[
 |B(q)|\le1+\sum_{k\ge0}r^{k(k+1)}\frac{1+r^{2k+1}}{1-r^{2k+1}}.
\]
Split at $k\asymp t^{-1/2}$ and $k\asymp t^{-1}$. The small-$k$ part is bounded by a harmonic sum times $t^{-1}$, and the Gaussian tails give the stated estimate.

In \eqref{lcp:ao:eq:E}, bound the finite numerator products by $(-r;r)_\infty^2$, the finite denominator products from below by $(r;r^2)_\infty$, and $\sum_{l=0}^jr^l$ by $(1-r)^{-1}$. Euler's identity gives
\[
\sum_{j\ge0}\frac{r^{j(j-1)/2}}{(r;r)_j}=(-1;r)_\infty=2(-r;r)_\infty.
\]
Thus
\begin{equation}\label{lcp:ao:eq:explicitR}
 |R(q)|\le1+\frac{2(-r;r)_\infty^3}
 {(1-r)(r;r)_\infty(r;r^2)_\infty}.
\end{equation}
The standard eta-product estimates
\begin{align*}
(r;r)_\infty&\sim\sqrt{\frac{2\pi}{t}}\ee^{-\pi^2/(6t)},\\
(-r;r)_\infty&\sim2^{-1/2}\ee^{\pi^2/(12t)},\qquad
(r;r^2)_\infty\sim\sqrt2\ee^{-\pi^2/(12t)}
\end{align*}
prove the second bound. Crucially, $\pi^2/2<\kappa=13\pi^2/24$.
\end{proof}

\begin{corollary}\label{lcp:ao:cor:exponential}
Let $p_n=[q^n]P(q)B(q)^2$. Then
\begin{equation}\label{lcp:ao:eq:expcoefficient}
 a_n-p_n=O\!\left(n^{1/4}\ee^{\pi\sqrt{2n}}\right).
\end{equation}
After Theorem~\ref{lcp:ao:thm:main}, this also gives
\begin{equation}\label{lcp:ao:eq:exprelative}
 \frac{a_n-p_n}{a_n}
 =O\!\left(n^{7/4}\ee^{-\pi(\sqrt{13/6}-\sqrt2)\sqrt n}\right).
\end{equation}
\end{corollary}
\begin{proof}
Cauchy's estimate and Lemma~\ref{lcp:ao:lem:global} bound the coefficient of $R$ by
$O(t^{-1/2}\exp(nt+\pi^2/(2t)))$. Choose its own optimizing radius
$t=\pi/\sqrt{2n}$. This is independent of the principal saddle used below.
\end{proof}
This is an exponentially smaller remainder for an \emph{exact principal generating function}. It is not a claim that the algebraic series in Theorem~\ref{lcp:ao:thm:main} converges, nor a complete transseries resolving all smaller exponential scales.

\begin{lemma}\label{lcp:ao:lem:minor}
There are constants $c,t_0>0$ such that for $0<t<t_0$ and $-\pi\le\theta\le\pi$,
\begin{equation}\label{lcp:ao:eq:minor}
 \frac{|P(\ee^{-t+\ii\theta})|}{P(\ee^{-t})}
 \le\exp\!\left[-c\min\left(\frac{\theta^2}{t^3},\frac1t\right)\right].
\end{equation}
\end{lemma}
\begin{proof}
Each numerator factor satisfies $|1+q^{2j+1}|\le1+r^{2j+1}$. For denominator factors on the consecutive block $t^{-1}\le j\le2t^{-1}$,
\[
 |1-r^j\ee^{\ii j\theta}|^2=(1-r^j)^2+2r^j(1-\cos j\theta).
\]
Since $r^j$ stays in a compact subinterval of $(0,1)$, taking logarithms gives an upper bound $-c_1\sum_j(1-\cos j\theta)$ for the product ratio. For $|\theta|\le c_2t$, the quadratic cosine bound gives $c_3\theta^2/t^3$. For $|\theta|/t$ in a fixed compact interval away from zero, rescaled Riemann sums converge to the strictly positive integral $\int_1^2(1-\cos(ux))\dd x$. For $|\theta|\ge C_1t$, choose $C_1$ large and use $|\sum_j\ee^{\ii j\theta}|\le1/|\sin(\theta/2)|$ to obtain a lower bound $c_4/t$. These ranges cover the circle.
\end{proof}

\begin{writenote}[Added 2 October 2026, batch~77P5]
The proofs of Lemmas~\ref{lcp:ao:lem:global} and~\ref{lcp:ao:lem:minor}
are sketched at the level of standard estimates (the split of the harmonic
sum, and the Riemann-sum compactness step); items 6--8 of the shipped
mathematical audit spell them out. Since $P=R^{\mathrm I}$,
Lemma~\ref{lcp:ao:lem:minor} is Part~I's bound \eqref{lcp:eq:minor} for the
same product, proved by the same argument; both texts are kept.
\end{writenote}

\begin{writenote}[Added 5 October 2026, batch~101]
Part~III proves Theorem~\ref{lcp:ao:thm:main} again
(Theorem~\ref{lcp:tr:thm:main}) without either lemma. Off the real axis it
uses only the positivity of the coefficients of $A$, a low-index part that
is exponentially smaller, and a damping bound for the finite product
$(q;q)_N$ whose three ranges are written out (\eqref{lcp:tr:tr:5},
\eqref{lcp:tr:tr:7}, \eqref{lcp:tr:tr:8}); in the saddle window it uses an
interpolation lemma instead of local expansions of $B$ and $P$
(Section~\ref{lcp:tr:sec:cross}).
\end{writenote}

\section{Uniform local expansions}\label{lcp:ao:sec:local}
From now on put $q=\ee^{-z}$, with $t=\Re z>0$ and $|\Im z|\le t/8$.
\begin{lemma}\label{lcp:ao:lem:localD}
For every fixed $M$, $D(\ee^{-z})$ has a uniform Poincar\'e expansion through degree $M-1$, obtained by Taylor expanding only its summands $0\le j<M$. Its remaining tail is $O(z^M)$.
\end{lemma}
\begin{proof}
Write $T_j$ for the $j$th summand of $D$. Then
\[
\frac{T_{j+1}}{T_j}=-q^j\frac{1-q^{j+1}}{1+q^{2j+3}}.
\]
Assume $0<t\le0.1$ and put $x=(j+1)t$. If $x\le1$, then
$|(2j+3)\Im z|\le(2+t)/8<\pi/2$, so the denominator has real part at least one. Also $|1-\ee^{-(j+1)z}|\le(j+1)|z|$. Therefore the ratio has modulus at most
\[
\ee^{0.1}\frac{\sqrt{65}}8xe^{-x}\le\frac{\ee^{0.1}\sqrt{65}}{8\ee}<\frac23.
\]
If $x\ge1$, the bound is
\[
\ee^t\frac{\ee^{-x}(1+\ee^{-x})}{1-\ee^{-2x}}
=\frac{\ee^t}{\ee^x-1}\le\frac{\ee^{0.1}}{\ee-1}<\frac23.
\]
Since $T_M=O(z^M)$ for fixed $M$, geometric summation proves the uniform tail estimate.
\end{proof}
The first coefficients are
\begin{align}\label{lcp:ao:eq:Dexp}
D(\ee^{-z})={}&\frac12-\frac{z^2}{8}+\frac{z^3}{8}-\frac{11z^4}{48}
 +\frac{15z^5}{32}-\frac{6737z^6}{5760}
 +\frac{6451z^7}{1920}-\frac{888521z^8}{80640}+O(z^9),\notag\\
D(\ee^{-z})^2={}&\frac14-\frac{z^2}{8}+\frac{z^3}{8}-\frac{41z^4}{192}
 +\frac{7z^5}{16}-\frac{6317z^6}{5760}
 +\frac{1529z^7}{480}-\frac{1702513z^8}{161280}+O(z^9).
\end{align}
These are exact rational coefficients, derived from finite Taylor arithmetic.

The eta quotient identity
\begin{equation}\label{lcp:ao:eq:eta}
P(q)=\frac{(q^2;q^2)_\infty^2}{(q;q)_\infty^4(q^4;q^4)_\infty}
\end{equation}
and eta inversion~\cite[23.18.5]{lcp-ao-dlmf} imply uniformly in the same wedge
\begin{equation}\label{lcp:ao:eq:Plocal}
 P(\ee^{-z})=\left(\frac{z}{2\pi}\right)^{3/2}
 \exp\!\left(\frac\kappa z-\frac z6\right)
 \left(1+O(\ee^{-c/t})\right),
\end{equation}
with the principal power of $z$. For example, the standard product formula is
\[
(\ee^{-z};\ee^{-z})_\infty
=\sqrt{\frac{2\pi}{z}}\exp\!\left(-\frac{\pi^2}{6z}+\frac z{24}\right)
\bigl(1+O(\ee^{-c/t})\bigr).
\]
Substitution at $z,2z,4z$ proves both the constant and the shift in \eqref{lcp:ao:eq:Plocal}.

\section{Coefficient extraction to every fixed order}\label{lcp:ao:sec:extraction}
\begin{proof}[Proof of Theorem~\ref{lcp:ao:thm:main}]
Use Cauchy's coefficient integral on $q=\ee^{-t+\ii\theta}$ with
$t=\sqrt{\kappa/N}$ and separate the exact remainder $R$. Lemma~\ref{lcp:ao:lem:global} makes its coefficient exponentially smaller than the result. Let
\[
\theta_0=t^{3/2}\log(1/t).
\]
On $|\theta|>\theta_0$, Lemma~\ref{lcp:ao:lem:minor}, together with the polynomial bound for $B$, makes $PB^2$ smaller than every algebraic relative order. On the remaining arc set $z=t-\ii\theta$. Lemma~\ref{lcp:ao:lem:localD} and \eqref{lcp:ao:eq:Plocal} give, for any fixed $M$,
\begin{equation}\label{lcp:ao:eq:localintegrand}
 P(\ee^{-z})B(\ee^{-z})^2\ee^{nz}
 =(2\pi)^{-3/2}\ee^{\kappa/z+Nz}
 \left(\sum_{m=0}^{M-1}b_mz^{3/2+m}+O(|z|^{3/2+M})\right),
\end{equation}
up to an exponentially small term. The real part of the phase has Gaussian decrease $-c\theta^2/t^3$; hence the integrated remainder is a relative $O(t^M)=O(N^{-M/2})$.

For explicit coefficients, consider a monomial $z^\alpha$ in \eqref{lcp:ao:eq:localintegrand}, with $\alpha=3/2+m$. Deform the truncated vertical saddle segment to the corresponding arc of $|z|=t$, adding radial connectors. Put $\lambda=\sqrt{\kappa N}$. On the circle the phase has real part $2\lambda\cos\phi$. On either connector, writing $|z|/t=s\in[1,\sec\phi]$, it is $\lambda\cos\phi(s+s^{-1})$. At the endpoint angle both are at most $2\lambda-c\log^2(1/t)$. Extending the circle to a Hankel contour cut along the negative real axis therefore changes the integral only by a superalgebraically small relative amount. The negative rays $z=-s$, $s\ge t$, have phase $-\kappa/s-Ns\le-2\lambda$.

The resulting Hankel integral is
\begin{equation}\label{lcp:ao:eq:Bessel}
\frac1{2\pi\ii}\int_{\mathcal H}z^\alpha\ee^{\kappa/z+Nz}\dd z
=\left(\frac\kappa N\right)^{(\alpha+1)/2}
 I_{-\alpha-1}(2\sqrt{\kappa N}).
\end{equation}
One can verify \eqref{lcp:ao:eq:Bessel} by expanding $\ee^{\kappa/z}$ and applying the reciprocal-Gamma Hankel formula term by term; $|z|\ge t$ ensures convergence. For fixed half-integral order the growing-exponential expansion of $I$~\cite[10.40.1]{lcp-ao-dlmf} terminates. Its coefficient at order $k$ is
\[
\frac{(-1)^k(m+2+k)!}{(m+2-k)!\,k!\,(4\sqrt{\kappa N})^k},
\qquad 0\le k\le m+2.
\]
Collecting total order $m+k=r$ gives \eqref{lcp:ao:eq:cr}. The leading constant is
\[
\frac{(2\pi)^{-3/2}b_0\kappa}{2\sqrt\pi}
=\frac{13\sqrt2}{768},\qquad b_0=\frac14,
\]
as required.
\end{proof}

\begin{writenote}[Added 2 October 2026, batch~77P5]
The transfer \eqref{lcp:ao:eq:Bessel}--\eqref{lcp:ao:eq:cr} is Part~I's
\eqref{lcp:eq:transfer}--\eqref{lcp:eq:betak} written in $N=n-1/6$: with
$m=j$, $\beta_k(j)/(2\sqrt{K^{\mathrm I}})^k$ equals the coefficient
displayed above, and the factor $\exp(-z/6)$ of \eqref{lcp:ao:eq:Plocal},
which Part~I expands into its $s_j$, is absorbed by the shift. Part~I
applies the transfer to $R^{\mathrm I}$ and to $F_3R^{\mathrm I}$; Part~II
applies it to $PB^2$, whose complement $R$ is exponentially smaller, and
that is what yields every order.
\end{writenote}

\begin{corollary}\label{lcp:ao:cor:logconcavity}
The sequence is eventually strictly log-concave: $a_n^2>a_{n-1}a_{n+1}$ for all sufficiently large $n$. More precisely,
\begin{equation}\label{lcp:ao:eq:logconcavity}
 \log\frac{a_n^2}{a_{n-1}a_{n+1}}
 =\frac{\sqrt\kappa}{2N^{3/2}}-\frac{3}{2N^2}+O(N^{-5/2}).
\end{equation}
\end{corollary}
\begin{proof}
Take $M=5$ in Theorem~\ref{lcp:ao:thm:main} and expand the logarithm. For fixed real constants $d_1,\ldots,d_4$ this gives
\[
 \log a_n=g(N)+O(N^{-5/2}),\qquad
 g(u)=\log C+2\sqrt{\kappa u}-\frac32\log u
       +\sum_{r=1}^4d_ru^{-r/2}.
\]
Taylor expansion of this explicit smooth function gives
\[
 g(N+1)-2g(N)+g(N-1)
 =-\frac{\sqrt\kappa}{2N^{3/2}}+\frac{3}{2N^2}+O(N^{-5/2}).
\]
The second difference of the three remainder values is still $O(N^{-5/2})$ by the triangle inequality; no derivative bound on the asymptotic remainder is assumed. Negating the resulting second difference proves \eqref{lcp:ao:eq:logconcavity}, whose leading term is positive.
\end{proof}

\begin{writenote}[Added 2 October 2026, batch~77P5]
From the exact coefficients through $n=2000$: $a_n^2>a_{n-1}a_{n+1}$ for
$3\le n\le1999$, and it fails only at $n=1,2$; the difference between
$\log(a_n^2/(a_{n-1}a_{n+1}))$ and the two displayed terms, times
$N^{5/2}$, is $0.674$ at $n=500$ and $0.593$ at $n=1999$. These are
numerical observations, not part of the proof; no effective threshold is
claimed.
\end{writenote}

\section{Continuous inverse models and integer thresholds}\label{lcp:ao:sec:inverse}
All inverse statements below concern the increasing large branch. Put
\begin{equation}\label{lcp:ao:eq:inversevars}
 y=2\sqrt{\kappa(n-1/6)},\qquad K=C(2\sqrt\kappa)^3,\qquad
 L=\log(X/K),\quad w=\log L.
\end{equation}
The transformed expansion has the form
\begin{equation}\label{lcp:ao:eq:yforward}
 X\sim K\frac{\ee^y}{y^3}\left(1+\sum_{r\ge1}\frac{h_r}{y^r}\right),
 \qquad h_r=(2\sqrt\kappa)^rc_r.
\end{equation}
Its first coefficients are
\begin{align*}
h_1&=-3,&h_2&=3-2\kappa^2,\\
h_3&=4\kappa^3+20\kappa^2,&
h_4&=-\frac{41\kappa^4}{3}-60\kappa^3-90\kappa^2.
\end{align*}
If $\log(1+\sum h_r y^{-r})\sim\sum\ell_r y^{-r}$, then
\begin{align}\label{lcp:ao:eq:ell}
\ell_1&=-3,&\ell_2&=-\frac32-2\kappa^2,\notag\\
\ell_3&=4\kappa^3+14\kappa^2,&
\ell_4&=\frac94-42\kappa^2-48\kappa^3-\frac{47\kappa^4}{3}.
\end{align}

\subsection{A branch explicit Lambert inverse}
The leading equation $X=K\ee^{y_0}/y_0^3$ has the exact large solution
\begin{equation}\label{lcp:ao:eq:W}
 y_0=-3W_{-1}\!\left(-\frac13(K/X)^{1/3}\right).
\end{equation}
The branch $W_{-1}$ is essential; it gives $y_0\to\infty$. Reverting \eqref{lcp:ao:eq:yforward} around $y_0$ gives
\begin{equation}\label{lcp:ao:eq:Winverse}
 y=y_0+\frac3{y_0}+\frac{21/2+2\kappa^2}{y_0^2}
 +\frac{45/2-8\kappa^2-4\kappa^3}{y_0^3}+O(y_0^{-4}),
 \qquad n=\frac16+\frac{y^2}{4\kappa}.
\end{equation}
These are inverse asymptotic models, not an assertion that a discrete sequence has a canonical noninteger inverse.

\subsection{Logarithmic inverse polynomials}
Substitute
\[
y=L+3\log L+\sum_{j\ge1}P_j(\log L)L^{-j}
\]
into $L=y-3\log y+\sum_{r\ge1}\ell_r y^{-r}$. At each new order $P_j$ enters linearly with coefficient one. This determines every polynomial recursively. In particular,
\begin{align}\label{lcp:ao:eq:logpolys}
 y={}&L+3w+\frac{9w+3}{L}
 +\frac{-27w^2/2+18w+21/2+2\kappa^2}{L^2}\notag\\
 &+\frac{27w^3-189w^2/2-9w+45/2-12\kappa^2w-8\kappa^2-4\kappa^3}{L^3}
 +O\!\left(\frac{w^4}{L^4}\right).
\end{align}
After any fixed order $J$, the remainder in the continuous inverse model for $y$ is
\[
O\!\left(\frac{(\log L)^{J+1}}{L^{J+1}}\right).
\]
Squaring and dividing by $4\kappa$ converts this to an index-scale error. The Lambert model and the logarithmic polynomials describe the same reversion in different coordinates.

\subsection{Shrinking integer brackets before rounding}
For $M\ge1$, define the explicit real function
\begin{equation}\label{lcp:ao:eq:model}
 F_M(x)=C(x-1/6)^{-3/2}\ee^{2\sqrt{\kappa(x-1/6)}}
 \sum_{r=0}^{M-1}c_r(x-1/6)^{-r/2}.
\end{equation}
It is positive and increasing for sufficiently large $x$, with
\[
\frac{\dd}{\dd x}\log F_M(x)=\frac{\sqrt\kappa}{\sqrt x}+O(x^{-1}).
\]
Let $x_M(X)$ be its unique large solution of $F_M(x)=X$, and let
$n_*(X)=\min\{n:a_n\ge X\}$.

\begin{corollary}\label{lcp:ao:cor:brackets}
For each fixed $M\ge2$, there are constants $D_M,X_M>0$ such that, for $X\ge X_M$,
\begin{equation}\label{lcp:ao:eq:brackets}
\left\lceil x_M(X)-\delta_M(X)\right\rceil
\le n_*(X)\le
\left\lceil x_M(X)+\delta_M(X)\right\rceil,
\qquad
\delta_M(X)=D_Mx_M(X)^{-(M-1)/2}.
\end{equation}
Thus the real bracket width tends to zero. If it does not cross an integer, the threshold is determined exactly by rounding its endpoints.
\end{corollary}
\begin{proof}
Theorem~\ref{lcp:ao:thm:main} gives $a_n/F_M(n)=1+O(n^{-M/2})$. The derivative above converts this relative error to an index displacement $O(n^{-(M-1)/2})$. Enlarge its implied constant to obtain strict comparisons on either side of $x_M\pm\delta_M$.

The sequence is nondecreasing: choose the topmost maximum-width row and adjoin one cell immediately to its left outside the bounding box. Row intervals of an $L$-convex polyomino are nested, since otherwise opposite extreme cells have no internal $L$ path. Thus the chosen row spans all columns, and the new cell connects to every old cell through that row using at most one turn. The enlarged row is uniquely longest, so removing its new leftmost cell recovers the original. This is an injection. Eventual strict increase also follows from the two-term expansion, because the positive order-$n^{-1/2}$ term in $a_{n+1}/a_n$ dominates its order-$n^{-1}$ uncertainty. Monotonicity and the two comparisons prove \eqref{lcp:ao:eq:brackets}.
\end{proof}
The constants in \eqref{lcp:ao:eq:brackets} are asymptotic existence constants. This report does not supply numerically certified values or a finite starting index. A small numerical model error by itself is therefore not a certified rounding rule, especially at a target $X=a_n$.

\begin{writenote}[Added 2 October 2026, batch~77P5: an instance of repository results]
As for Part~I's inverse (the note in Section~9.2), this section is an
instance of the transseries-and-inversion volume~\cite{lcp-tai}, and no
novelty is claimed for the method. The core $y_0-3\log y_0=L$ is
\texttt{p0:thm:lambert-core} with $a=1$, $b=-3$ and the branch $W_{-1}$;
the corrections in \eqref{lcp:ao:eq:Winverse} and the polynomials
\eqref{lcp:ao:eq:logpolys} are the reversion of
\texttt{p0:thm:perturbed-inversion}, in the Lambert-centred form
\texttt{p0:thm:lambert-centered}; Corollary~\ref{lcp:ao:cor:brackets} is
part~(2), the separation condition, of \texttt{p0:thm:staircase}, applied to
the model $F_M$ with the model error moved into $\delta_M$.

Relation to Part~I's Theorem~\ref{lcp:thm:inverse}. Since
$K=c(B^{\mathrm I})^3$, Part~II's $L$ at $X=y^{\mathrm I}$ is Part~I's
$\ell$, and $y_0=s_0$. Part~II's $y$ is $B^{\mathrm I}\sqrt{n-1/6}$, while
Part~I's $s$ is $B^{\mathrm I}\sqrt\nu$; hence
$s=y+(B^{\mathrm I})^2/(12y)+\cdots$, and $y=y_0+3/y_0+\cdots$ gives
$s=s_0+D_0/s_0+\cdots$ with $D_0=(B^{\mathrm I})^2/12+3$, Part~I's first
correction. The two inverses invert different real functions (Part~II's
model $F_M$, Part~I's interpolation $\widetilde a$), so they agree to the
orders both control, not term by term beyond them. Part~II's $n_*(X)$ is
Part~I's $N(y)$; with monotonicity proved for every $n$, Part~II needs no
starting index $n_0$ in the definition. The shipped
\texttt{data/02-all-orders-inverse-series.json} also records a fourth
Lambert-centred coefficient, $(47\kappa^4/3+36\kappa^3+45/4)/y_0^4$, and a
fourth logarithmic polynomial. They come from formal substitution
(residual zero through order~4) and are not printed in the manuscript.
\end{writenote}

\begin{writenote}[Added 5 October 2026, batch~101]
Part~III sharpens the monotonicity: the same injection, with the cell
adjoined on the right, misses the vertical bar of area $n+1$, so
$a_{n+1}\ge a_n+1$ for every $n\ge1$ (Section~\ref{lcp:tr:sec:inverse}).
It also proves that $n-n_0(a_n)\to1/6+36/(13\pi^2)<1/2$ for the leading
Lambert root $n_0=y_0^2/(4\kappa)$ (in this part's letters, $y_0$ at
$X=a_n$), so rounding $y_0^2/(4\kappa)$ recovers $n$ for all large $n$
(Remark~\ref{lcp:tr:rem:eta}). Its brackets
\eqref{lcp:tr:eq:threshold-bracket} are those of
Corollary~\ref{lcp:ao:cor:brackets} with the order shifted by one.
\end{writenote}

\section{Exact checks and reproducible calculations}\label{lcp:ao:sec:checks}
The proofs above do not infer their constants from numerical fits. For a separate check, set $g_n=f_{n-1}/(q;q)_n^2$. Then
\begin{equation}\label{lcp:ao:eq:compute}
(1-q^n)^2g_n=2g_{n-1}-g_{n-2},\quad g_{-1}=g_0=1,
\qquad A(q)=1+\sum_{n\ge1}q^n(1-q^n)g_n.
\end{equation}
Truncated power series division by $1-q^n$ is a simple shifted addition; using it twice gives exact integer coefficients in $O(N^2)$ arithmetic operations up to area $N$. The replay computes through $N=2000$ and matches the first 37 displayed OEIS values.

Table~\ref{lcp:ao:tab:errors} gives $(F_M(n)/a_n)-1$ for selected truncations. More terms improve these examples substantially, but a Poincar\'e expansion is a fixed-order statement as $n\to\infty$; indefinitely increasing $M$ at a fixed small $n$ need not improve accuracy.
\begin{table}[ht]
\centering\small
\begin{tabular}{r r r r r}
\toprule
$n$ & $M=1$ & $M=3$ & $M=5$ & $M=9$\\
\midrule
100  & $8.8963\,10^{-2}$ & $-9.3690\,10^{-3}$ & $-1.7838\,10^{-3}$ & $-5.2974\,10^{-4}$\\
200  & $5.8496\,10^{-2}$ & $-3.4962\,10^{-3}$ & $-3.4199\,10^{-4}$ & $-2.6476\,10^{-5}$\\
500  & $3.4317\,10^{-2}$ & $-9.3753\,10^{-4}$ & $-3.7809\,10^{-5}$ & $-4.9192\,10^{-7}$\\
1000 & $2.3243\,10^{-2}$ & $-3.4290\,10^{-4}$ & $-7.0411\,10^{-6}$ & $-2.3623\,10^{-8}$\\
2000 & $1.5900\,10^{-2}$ & $-1.2444\,10^{-4}$ & $-1.2961\,10^{-6}$ & $-1.1150\,10^{-9}$\\
\bottomrule
\end{tabular}
\caption{Signed relative error of the shifted forward model.}\label{lcp:ao:tab:errors}
\end{table}
The package contains exact coefficient generation, rational expansion arithmetic, inverse reversion checks, and independent evaluations of the exact decomposition at real and nonreal points. These checks test transcription and normalization; the complex bounds and coefficient-transfer proof supply the asymptotic theorem. A separate audit records the precise source version reviewed. Neither these computational checks nor that review constitute machine-checked formal verification or external peer review.

\begin{writenote}[Added 2 October 2026, batch~77P5: files and rerun]
In the repository the replay is \texttt{code/02-all-orders-replay.py}
(delivered as \texttt{code/replay.py}); its recorded outputs are the files
\texttt{data/02-all-orders-*.json}, and the ``separate audit'' is the three
\texttt{02-all-orders-*.md} reviews. The replay reads its reference
fixture from \texttt{../data/coefficients-2000.txt} relative to itself, a
file shipped as \texttt{data/02-all-orders-coefficients-2000.txt}, so it
runs only in the delivered layout; the report's README recreates that
layout on a scratch copy. The fixture agrees entry by entry with the
column $a_n$ of Part~I's \texttt{data/coefficients.csv}, $0\le n\le2000$:
the two parts computed the same 2001 integers by different recurrences.
On 2~October~2026 that recipe passed here: the core replay in 6~s, and
the replay with the optional SymPy and mpmath parts in 78~s. All five
outputs agreed with the shipped records apart from line endings (Windows
writes CRLF).
\end{writenote}

\begin{writenote}[Added 5 October 2026, batch~101]
Since batch~101 the fixture is known to be the complete OEIS b-file: Part~III's
snapshot of that b-file, byte-identical to the live file on 5~October 2026,
equals the fixture plus a final newline (Section~\ref{lcp:tr:sec:rechecked}).
Part~III extends the exact coefficients to $n=8000$.
\end{writenote}

\section{What remains open}\label{lcp:ao:sec:open}
\begin{enumerate}
\item \textbf{Effective constants and thresholds.} Quantifying all remainder bounds would turn the shrinking asymptotic brackets into certified finite algorithms with explicit starting ranges.
\item \textbf{Smaller exponential scales.} The exact remainder bound \eqref{lcp:ao:eq:expcoefficient} is deliberately coarse. Determining the actual dominant singular behavior of $R$, and then secondary contributions from $PB^2$, would be needed for a complete exponentially improved expansion.
\item \textbf{Large correction order.} The growth of the $b_m$ and $c_r$, optimal truncation, and possible resummation are separate problems. The fixed-order proof does not settle them.
\item \textbf{Further convexity classes.} The decisive step here is the exact product-square decomposition. Whether comparable transformations can be derived from area generating functions for other constrained polyomino classes is a concrete direction; this result alone makes no assertion about their asymptotics.
\item \textbf{The partial theta factor.} The Lambert, partial theta, and rapidly convergent $q$-series forms of $B$ may permit sharper arithmetic or analytic descriptions. No modularity claim for $B$ is required here.
\end{enumerate}

\begin{writenote}[Added 2 October 2026, batch~77P5]
These questions overlap Part~I's research questions of Section~11:
item~1 continues Part~I's question~2 (explicit error constants), item~2
its question~5 (beyond-all-orders terms), item~4 its question~7 (Z-convex
and other classes), and item~5 its question~4, since $B$ is Part~I's
$D^{\mathrm I}$ (Remark~\ref{lcp:ao:rem:LR}). Part~I's question~1 is
answered by Theorem~\ref{lcp:ao:thm:main}.
\end{writenote}

\section{Provenance of the two parts, and where the merge had to choose}\label{lcp:ao:sec:provenance}
This section was written when Part~II was added.

\begin{writenote}[Added 5 October 2026, batch~101]
Since batch~101 the report is built from three manuscripts. The third,
Part~III, has its own provenance section, Section~\ref{lcp:tr:sec:provenance}.
\end{writenote}

\paragraph{Two manuscripts.} The report is built from two manuscripts, each
printed in full.
\begin{itemize}[leftmargin=*,itemsep=0.15em]
\item \textbf{Part~I} (Sections~1--12, Appendices~A and~B): manuscript~05
 of batch~75, ``L-convex polyominoes and a colored-partition comparison'',
 archive \texttt{A126764\_\allowbreak Lconvex\_\allowbreak Research.zip},
 arrival commit \texttt{4b874cea0}, placed in \texttt{6ea60e367} and
 written in \texttt{936896776}. It contributed the squeeze, the finite-cut
 proof of three corrections, the radial theorem, the interpolated inverse
 and the exact and numerical checks of Section~10. It names no ProveIt
 commit (its repository input was an overview and a targeted search), so it
 has no pin.
\item \textbf{Part~II} (Sections~\ref{lcp:ao:sec:main}--\ref{lcp:ao:sec:open}):
 manuscript~31 of batch~77, ``All order asymptotics and inversion for L
 convex polyominoes'', archive
 \texttt{lconvex-\allowbreak asymptotics-\allowbreak reproducibility.zip},
 arrival commit \texttt{096ee7b87}, placed in \texttt{4f11bc9c0}. It
 contributed the exact product-square decomposition, the all-orders
 coefficient theorem, log-concavity, monotonicity by injection, and the
 shifted inverse with integer brackets. It names no ProveIt commit and does
 not mention the repository, so it has no pin either.
\end{itemize}

\paragraph{Where the merge had to choose.}
\begin{itemize}[leftmargin=*,itemsep=0.15em]
\item \emph{Placement.} Part~II is printed after Part~I's closing
 Section~12 and before Part~I's appendices, so no section, theorem or
 equation number of Part~I moves. Part~I's Appendix~A remains the
 reproduction appendix of Part~I; Part~II's reproduction notes are in
 Section~\ref{lcp:ao:sec:checks}.
\item \emph{Letters.} Part~II keeps the manuscript's letters, which the
 shipped data and reviews use; the clashes with Part~I are resolved by
 Table~\ref{lcp:ao:tab:notation} and by the superscript $\mathrm I$ in
 write notes, not by renaming (Section~\ref{lcp:ao:sec:notation}). The
 manuscript's macros \verb|\e|, \verb|\ii| and \verb|\dd| were mapped to
 Part~I's \verb|\ee| and two new preamble macros; three unused macros of the
 manuscript were dropped.
\item \emph{Duplicated results.} Results proved in both manuscripts keep
 both proofs, with cross-references (Section~\ref{lcp:ao:sec:intro}); no
 proof was shortened.
\item \emph{Bibliography.} The manuscript's entries for
 Castiglione--Frosini--Munarini--Restivo--Rinaldi, Guttmann--Kot\v e\v
 sovec, Guttmann--Massazza and OEIS A126764 cite the same works as Part~I's
 entries \cite{castiglione,gk,zconvex,oeis} and were merged into them (the
 page reference for Guttmann--Massazza is noted in that entry). Its DLMF
 entry cites other sections and a stated version, so it is kept as a
 separate entry \cite{lcp-ao-dlmf}.
\item \emph{Additions.} Section~\ref{lcp:ao:sec:intro}, this section, the
 remarks headed ``write'' and the \textsf{[write]} notes in Part~II are
 editorial; Remark~\ref{lcp:ao:rem:LR} is a write-phase observation, with
 its proof and a numerical check. No statement, proof or number of either
 manuscript was changed; labels of Part~II carry the prefix
 \texttt{lcp:ao:}.
\item \emph{Not shipped.} The delivered PDF of manuscript~31, its build
 script, its delivery README, its checksum manifest
 \texttt{MANIFEST.json} with the package verifier, the replay's own copy of
 its coefficient output (equal to the shipped fixture plus a final
 newline) and the record of its PDF build. The archive survives in the
 arrival commit; retrieve it with \texttt{git show
 096ee7b87:\allowbreak docs/\allowbreak incoming/\allowbreak lconvex-\allowbreak asymptotics-\allowbreak reproducibility.zip}.
\end{itemize}


%% ====================================================================
%% Batch-101 addition (bundle Report 95, "All Order Area Asymptotics for
%% L Convex Polyominoes").  All labels of Part III carry the prefix
%% lcp:tr: (transfer).  Sections lcp:tr:sec:intro, lcp:tr:sec:further and
%% lcp:tr:sec:provenance and every text marked [write] were added when
%% Part III was written into the report; the other Part III sections
%% print the manuscript.
%% ====================================================================
\clearpage
\phantomsection\part*{Part III.\quad All-order area asymptotics: a complete transfer proof}
\addcontentsline{toc}{part}{Part III.\texorpdfstring{\quad}{ } All-order area asymptotics: a complete transfer proof}

\noindent\emph{Parts~I and~II (Sections~1--23 and Appendices~A and~B) are
unchanged except for dated \textsf{[write]} notes of 5~October 2026 (after
the scope box, in Section~1.1, at the end of Section~8.2, in Section~10.1,
in research question~7 of Section~11, in Sections~14, 17, 20, 21 and~23),
a line in the title, a contents entry for this part, and a note in the
bibliography entries \texttt{castiglione} and \texttt{zconvex}. No number of
Part~I or Part~II moves.}

\section{Part III: source, scope and notation}\label{lcp:tr:sec:intro}
This section was written when Part~III was added.
Sections~\ref{lcp:tr:sec:main}--\ref{lcp:tr:sec:lambert} print the
manuscript; Section~\ref{lcp:tr:sec:further} collects its unproved claims
and open directions, and Section~\ref{lcp:tr:sec:provenance} records the
provenance of Part~III and the choices the write made.

\subsection{Source and status}\label{lcp:tr:sec:source}
Sections~\ref{lcp:tr:sec:main}--\ref{lcp:tr:sec:lambert} print, in full,
bundle Report~95 of ProveIt's incoming-reports intake (batch~101), titled
``All Order Area Asymptotics for L Convex Polyominoes'' and dated
2~October 2026. It arrived as the archive
\texttt{L\_\allowbreak Convex\_\allowbreak Polyomino\_\allowbreak Area\_\allowbreak Asymptotics\_\allowbreak Source.zip}
in commit \texttt{60f54ea06} and was placed in commit \texttt{f7c612c72}.
Its author line, and the author field of its PDF, read ``Research note
prepared with OpenAI''; like every manuscript of the intake it is treated
as AI-assisted research output. It names no ProveIt commit and does not
mention the repository, so it has no pin. Its only input from this
report's history is Part~I's manuscript, which it cites as a separate
unpublished manuscript of 1~October 2026~\cite{lcp-tr-earlier}. It does not
know Part~II, which was written into the report on 2~October 2026 (commits
\texttt{4f11bc9c0} and \texttt{424734dc4}). Part~III is unrefereed. Nothing
in it is formalized in Lean or Rocq, and its place in the collection
confers no formal status on it. The manuscript's Sections~1--7 and its
Appendix~A are Sections~\ref{lcp:tr:sec:main}--\ref{lcp:tr:sec:lambert}
here; a section number written out in its text (``Sections~2--4'' in
Section~\ref{lcp:tr:sec:checks}) uses its own numbering, so that its
Sections~2--4 are Sections~\ref{lcp:tr:sec:product}--\ref{lcp:tr:sec:transfer}.

The manuscript's own limitations are printed with it and remain in force:
it gives no complete transseries and identifies no exponentially small
term; its remainder constants are not effective, so its rounding
statements are eventual and the onset~$556$ of Section~\ref{lcp:tr:sec:checks}
is an observation, not a proved cutoff; thresholds for arbitrary values
need the ceiling brackets of Section~\ref{lcp:tr:sec:inverse}; the
coefficients $a_{2001},\ldots,a_{8000}$ are generated, not compared with
published data; its literature search was targeted, and W.~R.~G.~James's
2005 thesis was not checked in full; it does not claim the first proof of
the leading equivalent; its Bessel expressions are asymptotic shorthands,
not convergent sums; and it claims no journal publication.

\subsection{What was rechecked at intake and in this write}\label{lcp:tr:sec:rechecked}
Independently of the delivery, on 5~October 2026:
\begin{itemize}[leftmargin=*,itemsep=0.15em]
\item the delivered suite was rerun on a copy in the delivered layout
 (Python~3.14.4, SymPy~1.14.0): the multiplier expansion
 $d_0,\ldots,d_8$ and the agreement of the two formulas for $h$ through
 order~8; the Gaussian-saddle check through $z^{-4}$ and the formal inverse
 through $Y^{-4}$; the exact coefficients through $n=8000$ with all 2001
 b-file terms and the literal expansion through $n=400$; the diagnostics
 through $n=8000$; and the rounding scan. All passed, and every recorded
 output was reproduced apart from line endings and timing fields;
\item the b-file snapshot that the suite compares with is byte-identical to
 the live OEIS b-file (downloaded at intake, SHA-256
 \texttt{d036795c\ldots}), and it equals Part~II's fixture
 \texttt{data/\allowbreak 02-all-orders-coefficients-2000.txt} plus a final
 newline;
\item Part~I's \texttt{code/verify.py}, a different implementation, was run
 on a copy with \texttt{--nmax 8000}: all its assertions, among them the
 squeeze $0\le p_n-4a_n\le3j_n$ of Theorem~\ref{lcp:thm:squeeze}, hold
 through $n=8000$, and its $a_n$ equal the shipped
 \texttt{data/\allowbreak 03-transfer-exact-exact\_coefficients\_8000.json}
 for every $0\le n\le8000$;
\item the identities of Section~\ref{lcp:tr:sec:product} were evaluated at
 $q=0.3$ and $0.6$ to 40 digits: $h$ against the alternating series
 \eqref{lcp:tr:eq:T} and against Part~II's $B$; $A-1-Mh^2-R$; the
 remainder $R$ of \eqref{lcp:tr:eq:Rformula} against Part~II's
 $E(q)/(q;q)_\infty$; and the bound \eqref{lcp:tr:eq:Rbound}. The
 discrepancies are at most $4\cdot10^{-28}$;
\item in this write: $c_0,\ldots,c_4$ from \eqref{lcp:tr:eq:cr} and
 \eqref{lcp:tr:eq:dvalues}, and their exact (SymPy) equality with Part~I's
 $1,-d,e,f$ and with Part~II's $c_0,\ldots,c_4$ re-expanded at
 $N=n-1/6$ (Remark~\ref{lcp:tr:rem:shift}); the Al-Salam--Chihara
 identification (Remark~\ref{lcp:tr:rem:asc}) against the basic
 hypergeometric definition; the contour identity behind
 Lemma~\ref{lcp:tr:lem:hankel}, numerically; the value of $\eta$; and
 $a_{n+1}-a_n\ge1$ for $1\le n<8000$ in the shipped data, with equality
 only at $n=1$.
\end{itemize}
The transfer argument of Section~\ref{lcp:tr:sec:transfer} was checked line
by line at intake (the positivity induction, the low-index bound, the
damping product, the interpolation constants, the saddle window and the
Bessel collection) and again in this write, which spells out the last step
in Lemma~\ref{lcp:tr:lem:hankel}.

\subsection{What Part III adds, what is proved twice, and three proofs of the main theorem}\label{lcp:tr:sec:cross}
Part~III proves, from the same cited area generating function:
\begin{itemize}[leftmargin=*,itemsep=0.15em]
\item the exact decomposition $A=1+Mh^2+R$ with a positive series for the
 remainder and two real-axis bounds
 (Proposition~\ref{lcp:tr:prop:factor}). The identity itself is Part~II's
 \eqref{lcp:ao:eq:exactmain}: Part~III's $M$ is Part~II's $P$, its $h$ is
 Part~II's $B=D$, and its $R$ is Part~II's $R-1=E/(q;q)_\infty$.
 \textbf{New} are the duality proof \eqref{lcp:tr:eq:duality}, the positive
 formula \eqref{lcp:tr:eq:Rformula}, the sign $R\ge0$ on $(0,1)$ and the
 bounds \eqref{lcp:tr:eq:Rbound};
\item the multiplier $h$ by reduction of order, with the explicit tail
 bound \eqref{lcp:tr:eq:htail} and a proof that $h>0$, and by the
 alternating series of Proposition~\ref{lcp:tr:prop:T}, which is Part~II's
 series $D$ in \eqref{lcp:ao:eq:PBD} because
 $(-q;q^2)_{j+1}=(1+q)(-q^3;q^2)_j$. Its alternating error bound is new;
\item the complete real-axis expansion \eqref{lcp:tr:eq:realexp}: Part~I's
 Theorem~\ref{lcp:thm:radial} again, by a third route;
\item \textbf{a third proof of the all-orders coefficient theorem}
 (Theorem~\ref{lcp:tr:thm:main}), by positivity, a common partition factor
 and Chebyshev interpolation (Section~\ref{lcp:tr:sec:transfer}). Every
 estimate of this route is written out;
\item \textbf{the strict increase} $a_{n+1}\ge a_n+1$ for every $n\ge1$
 (Section~\ref{lcp:tr:sec:inverse}); Part~II proves that $a_n$ is
 nondecreasing, and strictly increasing only eventually;
\item \textbf{the limit} $n-n_0(a_n)\to\eta=1/6+36/(13\pi^2)$ for the
 leading Lambert model, hence its eventual exact-range rounding rule
 (Remark~\ref{lcp:tr:rem:eta}), with a finite scan through $n=8000$;
\item the coefficients in the variable $z=\beta\sqrt n$, the logarithmic
 coefficients $\lambda_1,\ldots,\lambda_4$, a power--log inverse in
 $Y=\log y$, and threshold brackets (the brackets are Part~II's
 Corollary~\ref{lcp:ao:cor:brackets}, with the order shifted by one);
\item exact coefficients through $n=8000$ and a comparison with all 2001
 b-file terms.
\end{itemize}

\paragraph{Three proofs of the coefficient asymptotics.} The leading
equivalent and the corrections $d,e,f$ of Theorem~\ref{lcp:thm:main} are
now proved three times, and the all-orders expansion twice. The routes
differ at every step.

{\small
\begin{longtable}{@{}>{\raggedright}p{0.14\textwidth}>{\raggedright}p{0.24\textwidth}>{\raggedright}p{0.25\textwidth}>{\raggedright\arraybackslash}p{0.27\textwidth}@{}}
\caption{The three proofs of the coefficient asymptotics.}\label{lcp:tr:tab:proofs}\\
\toprule
Step & Part~I & Part~II & Part~III\\
\midrule
\endfirsthead
\toprule
Step & Part~I & Part~II & Part~III\\
\midrule
\endhead
\bottomrule
\endfoot
Orders proved & leading term and $d,e,f$; relative remainder $\OO(n^{-7/4})$ (Theorem~\ref{lcp:thm:main}) & every order, in $N=n-1/6$ (Theorem~\ref{lcp:ao:thm:main}); $\OO(n^{-2})$ at $M=4$ & every order, in $n$ (Theorem~\ref{lcp:tr:thm:main}); $\OO(n^{-2})$ at $m=4$\\
Decomposition & squeeze \eqref{lcp:eq:squeeze} and finite cut \eqref{lcp:eq:cut} at $m=3$ & $A=PB^2+R$, by Heine and Fine (Proposition~\ref{lcp:ao:prop:square}) & $A=1+Mh^2+R$, by duality (Proposition~\ref{lcp:tr:prop:factor})\\
Whole circle & positive tails $E_3$, $C_4$ and explicit integrals of $|Q_3|^{-1}$, $|Q_3|^{-2}$ (Proposition~\ref{lcp:prop:approxerror}) & bounds for $|B|$ and $|R|$ (Lemma~\ref{lcp:ao:lem:global}) and minor arcs for $P$ (Lemma~\ref{lcp:ao:lem:minor}), both sketched & positivity of $A$, a low part $L_N$ and a damping factor $D_N$ from the common factor $(q;q)_N$ \eqref{lcp:tr:tr:7}, \eqref{lcp:tr:tr:8}, proved in full\\
Saddle window & eta expansion of $R^{\mathrm I}$ and Taylor expansion of the rational $F_3$ & local expansions of $D$ and $P$ in a wedge (Lemma~\ref{lcp:ao:lem:localD}, \eqref{lcp:ao:eq:Plocal}) & Chebyshev interpolation of the real expansion \eqref{lcp:tr:tr:9}--\eqref{lcp:tr:tr:11}: no complex estimate of $h$ or $M$\\
Extraction & Bessel transfer \eqref{lcp:eq:transfer} & Hankel--Bessel \eqref{lcp:ao:eq:Bessel} & Bessel \eqref{lcp:tr:tr:12}--\eqref{lcp:tr:tr:14}, completed by Lemma~\ref{lcp:tr:lem:hankel}\\
\end{longtable}}

\paragraph{Proved in two or three parts.} Each part keeps its own proof.
\begin{itemize}[leftmargin=*,itemsep=0.15em]
\item The leading equivalent: Part~I (Theorems~\ref{lcp:thm:main}
 and~\ref{lcp:thm:squeeze}), Part~II (Theorem~\ref{lcp:ao:thm:main},
 $M=1$), Part~III (Theorem~\ref{lcp:tr:thm:main}, $m=1$).
\item The corrections $d,e,f$: the three parts, with the coefficients
 equal exactly (Remark~\ref{lcp:tr:rem:shift}).
\item The all-orders coefficient expansion: Parts~II and~III. Their
 coefficients determine each other (Remark~\ref{lcp:tr:rem:shift}).
\item The radial expansion: Part~I's Theorem~\ref{lcp:thm:radial} (rational
 approximants $h_m/Q_m$), Part~II (from \eqref{lcp:ao:eq:exactmain} and
 Lemma~\ref{lcp:ao:lem:localD}; Remark~\ref{lcp:ao:rem:LR}) and Part~III
 (from \eqref{lcp:tr:eq:factor} and the tail bound \eqref{lcp:tr:eq:htail}).
 The coefficients agree: $\sum_jd_jt^j=\ee^{-t/6}\mathcal F(t)$ with
 Part~I's $\mathcal F$ of \eqref{lcp:eq:Fintro}.
\item The product generating function: Part~I's \eqref{lcp:eq:Gproduct},
 Part~II's $G$ in \eqref{lcp:ao:eq:HG}, Part~III's \eqref{lcp:tr:eq:P}.
 The discrete Wronskian: \eqref{lcp:eq:wronskian} and
 \eqref{lcp:tr:eq:reduction}.
\item Part~I's $D^{\mathrm I}=L/R^{\mathrm I}$, Part~II's $B=D$ and Part~III's
 $h$ are one function: Part~II's Remark~\ref{lcp:ao:rem:LR}, and
 Part~III's definition \eqref{lcp:tr:eq:h} with
 Proposition~\ref{lcp:tr:prop:T}. Part~III's Lambert series
 \eqref{lcp:tr:eq:LambertH} is Part~II's \eqref{lcp:ao:eq:Lambert}, by the
 same Heine transformation.
\item The eta expansion: \eqref{lcp:eq:Rlocal}, \eqref{lcp:ao:eq:Plocal},
 \eqref{lcp:tr:eq:Mexp}. The minor-arc damping: Part~I's \eqref{lcp:eq:minor}
 and Part~II's Lemma~\ref{lcp:ao:lem:minor} for the infinite product,
 Part~III's \eqref{lcp:tr:tr:7} for the finite product $(q;q)_N$, by the
 same three-range argument.
\item Monotonicity by an injection: Part~II adjoins a cell on the left and
 proves $a_{n+1}\ge a_n$; Part~III adjoins it on the right and, noting
 that the vertical bar is not in the image, proves $a_{n+1}\ge a_n+1$.
\item The Lambert inverse: Part~I's $s_0$, Part~II's $y_0$ and Part~III's
 $z_0$ are the same number at the same target; all three are instances of
 \texttt{p0:thm:lambert-core}~\cite{lcp-tai}. Part~III's brackets
 \eqref{lcp:tr:eq:threshold-bracket} are Part~II's
 \eqref{lcp:ao:eq:brackets}.
\end{itemize}
Only Part~I has the coefficientwise squeeze with its colored-partition
reading, the deficit limit, the finite-cut identity, the interpolation
$\nu(y)$ and the tables of Section~10. Only Part~II has the partial theta
form \eqref{lcp:ao:eq:theta}, the coefficient bound
\eqref{lcp:ao:eq:expcoefficient} for its remainder, and eventual
log-concavity.

\subsection{Notation across the three parts}\label{lcp:tr:sec:notation}
Part~III keeps the manuscript's letters, as Part~II did: its shipped
scripts and data use them (for example \texttt{b\_in\_K} and
\texttt{lambda\_in\_K} in \texttt{data/\allowbreak 03-transfer-symbolic\_results.json}).
In Sections~\ref{lcp:tr:sec:main}--\ref{lcp:tr:sec:lambert} a letter has its
Part~III meaning. Where a \textsf{[write]} note in Part~III refers to a
clashing object of another part, it adds the superscript $\mathrm I$ or
$\mathrm{II}$. No symbol was renamed and no normalization was changed.
Table~\ref{lcp:tr:tab:notation} lists every clash; the last column gives
the tempting false reading. The manuscript itself reuses several letters
(rows marked ``internal'').

{\small
\begin{longtable}{@{}>{\raggedright}p{0.1\textwidth}>{\raggedright}p{0.27\textwidth}>{\raggedright}p{0.25\textwidth}>{\raggedright\arraybackslash}p{0.3\textwidth}@{}}
\caption{Letters of Part~III, the same object in Parts~I and~II, and the clash to avoid.}\label{lcp:tr:tab:notation}\\
\toprule
Part III & Meaning in Part III & Same object in Parts I and II & Not to be confused with\\
\midrule
\endfirsthead
\toprule
Part III & Meaning in Part III & Same object in Parts I and II & Not to be confused with\\
\midrule
\endhead
\bottomrule
\endfoot
$A(q)$ & the area generating function & Part~I's $\Acal(q)$; Part~II's $A(q)$ & Part~III's constant $\mathcal A$ below\\
$K$ & $13\pi^2/24$ & Part~I's $K$; Part~II's $\kappa$ & Part~II's $K=C(2\sqrt\kappa)^3$, which is Part~III's $\mathcal A$\\
$C$ & $1/(4(2\pi)^{3/2})$, the real-axis amplitude \eqref{lcp:tr:eq:constants} & written out in Part~I's \eqref{lcp:eq:radialintro} & Part~II's $C=13\sqrt2/768$, which is Part~III's $\delta$; Part~I's series $C(z;q)$; the constants $C_0$, $C_M$ of Sections~\ref{lcp:tr:sec:transfer}--\ref{lcp:tr:sec:inverse}\\
$\delta$ & $13\sqrt2/768$ & Part~I's $c$; Part~II's $C$ & Part~I's correction $\delta$ in the proof of Theorem~\ref{lcp:thm:inverse}; Part~II's $\delta_M$\\
$\beta$ & $2\sqrt K=\pi\sqrt{13/6}$ & Part~I's constant $B$; Part~II's $2\sqrt\kappa$ & Part~II's function $B(q)$; Part~III's constant $B$ (Section~\ref{lcp:tr:sec:transfer})\\
$\mathcal A$ & $\delta\beta^3$ (Section~\ref{lcp:tr:sec:inverse}) & Part~I's $cB^3$; Part~II's $K$ & Part~I's generating function $\Acal(q)$\\
$p_n$, $p_\infty$ & $(q;q)_n$, $(q;q)_\infty$ & Part~I's $(q;q)_n$ & Part~I's $p_n=[q^n]R^{\mathrm I}$ ($\sim4a_n$); Part~II's $p_n=[q^n]PB^2$ ($\sim a_n$)\\
$u_n$ & $f_{n-1}$, with $u_0=1$ & Part~I's $h_n$ & Part~I's $u_n=Q_n-h_n$; Part~II's $u_l$\\
$Q_n$ & the auxiliary solution \eqref{lcp:tr:eq:Qrec} & Part~I's $Q_n$, the same & the inverse corrections $Q_r$ of Section~\ref{lcp:tr:sec:inverse} (internal)\\
$V_n$ & $2q^n-q^{2n}$ & Part~I's $v_j$ in the proof of Lemma~\ref{lcp:lem:slope}; Part~II's $\epsilon_k$ & Part~I's $v_n=Q_n-2h_n$; Part~III's $v_n$ of Section~\ref{lcp:tr:sec:transfer}\\
$P(z)$ & $\sum Q_nz^n/p_n$, \eqref{lcp:tr:eq:P} & Part~I's and Part~II's $G(z)$ & Part~II's $P(q)$, which is Part~III's $M$; Part~I's $P_m$; Part~III's $P_r(T)$ and its order $P$ in Section~\ref{lcp:tr:sec:transfer} (internal)\\
$D(q)$ & $(-q;q^2)_\infty/p_\infty$, the limit of the increments of $Q_n$, \eqref{lcp:tr:eq:D} & the slope $s(q)$ of Lemma~\ref{lcp:lem:slope} for $y=Q$ & \textbf{Part~I's $D=L/R$ and Part~II's $D=B$, which are Part~III's $h$}; the constant $D$ in $k=\lceil D\log(1/t)\rceil$ (internal); Part~I's $D_0$; Part~II's $D_M$\\
$h(q)$ & the multiplier \eqref{lcp:tr:eq:h} & Part~I's $D=L/R$; Part~II's $B=D$ & Part~I's $h_n$; Part~II's $h_n$ and $h_r$\\
$B_Q(z)$, $A_Q$ & $\sum Q_nz^n/p_n^2$; $B_Q(q)-B_Q(q^2)=Mh$ & $B_Q(q^j)$ is Part~I's $C^Q_j$ & Part~II's $B(q)$; Part~III's constant $B$ (internal)\\
$M(q)$ & the eta quotient \eqref{lcp:tr:eq:eta} & Part~I's $R(q)$; Part~II's $P(q)$ & Part~I's model $M(x)$; Part~II's order $M$; the bound $M$ of the interpolation lemma and the orders $M$ of Section~\ref{lcp:tr:sec:inverse} (internal)\\
$R(q)$ & the remainder \eqref{lcp:tr:eq:Rformula}, $\ge0$ on $(0,1)$ & Part~II's $R-1=E/(q;q)_\infty$ & Part~I's $R(q)$, which is Part~III's $M$; Part~II's $R$ itself; the integer $R$ in \eqref{lcp:tr:eq:coef-full} (internal)\\
$w_r$ & $B_Q(q^{r+1})$, Section~\ref{lcp:tr:sec:product} & --- & $w_n=v_n-v_{n-1}$ of Section~\ref{lcp:tr:sec:transfer} (internal); Part~II's $w=\log L$\\
$d_j$ & coefficients of the real expansion \eqref{lcp:tr:eq:realexp}: $\sum d_jt^j=4\ee^{-t/6}h(\ee^{-t})^2$ & $\ee^{-t/6}\mathcal F(t)$ of Part~I; $4\ee^{-t/6}\sum b_mt^m$ with Part~II's $b_m$ & the variable $d=1-q$ of Section~\ref{lcp:tr:sec:real} (internal); Part~I's constant $d$; Part~II's local $d_r$\\
$c_r$ & unshifted coefficients \eqref{lcp:tr:eq:cr} & Part~I's $1,-d,e,f$ for $r\le3$ & Part~II's $c_r$, which are shifted (Remark~\ref{lcp:tr:rem:shift}); Part~I's $c$\\
$v_n$, $w_n$ & $u_n/(q;q)_n$ and $v_n-v_{n-1}$ (Section~\ref{lcp:tr:sec:transfer}) & $v_n$ is Part~II's $h_n$ & Part~I's $v_n$; Part~III's $w_r$\\
$a_j$ & $2r^j-r^{2j}$ in Section~\ref{lcp:tr:sec:transfer} & $V_j$ at $q=r$ & the coefficients $a_n$ (internal)\\
$L_N$, $H_N$, $G_N$, $D_N$ & low and high parts of \eqref{lcp:tr:tr:3}, the numerator of $H_N$, the damping factor & --- & the truncation index $L$ of Section~\ref{lcp:tr:sec:transfer} (internal); Part~I's slope $L$; Part~II's $L$; $G(z)$ of Parts~I and~II; $D(q)$\\
$N$ & $\lfloor1/t\rfloor$ in Section~\ref{lcp:tr:sec:transfer}; the threshold $N(y)$ in Section~\ref{lcp:tr:sec:inverse} & $N(y)$ is Part~I's $N(y)$ and Part~II's $n_*$ & Part~II's $N=n-1/6$\\
$E$ & $3+\log2$, \eqref{lcp:tr:tr:5} & --- & Part~I's defect $E$; Part~II's $E(q)$\\
$B$, $D$, $L$, $P$ & constants of the interpolation step in Section~\ref{lcp:tr:sec:transfer} & --- & the functions $B_Q$, $D(q)$, $L_N$, $P(z)$ (internal)\\
$z$ & $\beta\sqrt n$ in Section~\ref{lcp:tr:sec:inverse}; $q=\ee^{-z}$ in Section~\ref{lcp:tr:sec:transfer}; a generating-function variable in Sections~\ref{lcp:tr:sec:product} and~\ref{lcp:tr:sec:lambert} & Part~II's $y=\beta\sqrt{n-1/6}$ is the shifted analogue & the three uses (internal)\\
$b_r$ & $c_r\beta^r$, \eqref{lcp:tr:eq:b-rule} & --- & $b_n(q)$ of Section~\ref{lcp:tr:sec:checks} (internal; the source says so); Part~II's $b_m$; Part~I's $b_n$, $b_j$\\
$\lambda_r$ & coefficients of $\log(1+\sum b_ru^r)$ & the unshifted analogue of Part~II's $\ell_r$ (not equal) & Part~II's $\lambda=\sqrt{\kappa N}$; Part~I's amplitude $\lambda$\\
$y$, $Y$, $T$ & target value; $\log y$; $3\log Y-\log\mathcal A$ & Part~I's $y$; Part~II's $X$; $Y=L+\log\mathcal A$ with Part~II's $L$ & Part~II's $y$ and $w$; Proposition~\ref{lcp:tr:prop:T}; the Chebyshev $T_k$; $y=\Im z$ in \eqref{lcp:tr:tr:10} (internal); Part~I's $T_n(q)$\\
$P_r(T)$ & power--log inverse polynomials \eqref{lcp:tr:eq:z-inverse} & --- & Part~II's $P_j(\log L)$, polynomials in another variable\\
$z_0$, $n_0$ & Lambert root; $z_0^2/(4K)$ & Part~I's $s_0$; Part~II's $y_0$ & Part~I's starting index $n_0$ in $N(y)$\\
$F_M$, $r_M$ & logarithmic model and its root (Section~\ref{lcp:tr:sec:inverse}) & $r_M$ plays the role of Part~II's $x_M$ & Part~II's $F_M$, a non-logarithmic shifted model; Part~I's $F_3$; the radius $r$\\
$S(z)$, $H(z)$ & $\sum u_nz^n/p_n$ and $S/P$ (Section~\ref{lcp:tr:sec:lambert}) & $S$ is Part~II's $H(z)$; $H$ is the bracket of \eqref{lcp:ao:eq:Hsolution} & Part~II's $S=1-B$ and $H$\\
$a$, $b$ & $\pm\ii\sqrt q$ (Section~\ref{lcp:tr:sec:lambert}) & the same in Part~II & $a_n$; $b_r$\\
\end{longtable}}

\subsection{Unshifted and shifted coefficients}\label{lcp:tr:sec:shift}
Part~III expands in powers of $n^{-1/2}$, Part~II in powers of
$N^{\mathrm{II}}=n-1/6$. The following remark converts between them.

\begin{remark}[write, 5 October 2026: the conversion]\label{lcp:tr:rem:shift}
Put $x=n^{-1/2}$, so that $N^{\mathrm{II}}=n(1-x^2/6)$, and let $c_s^{\mathrm{II}}$ be
Part~II's coefficients \eqref{lcp:ao:eq:firstc} (with $\kappa=K$ and
$C^{\mathrm{II}}=\delta$). Then the coefficients $c_r$ of
Theorem~\ref{lcp:tr:thm:main} are the Taylor coefficients of
\begin{equation}\label{lcp:tr:eq:shift}
 \sum_{r\ge0}c_rx^r=\sum_{s\ge0}c_s^{\mathrm{II}}x^s
 \Bigl(1-\frac{x^2}6\Bigr)^{-(3+s)/2}
 \exp\Bigl(\frac{2\sqrt K}{x}\Bigl(\sqrt{1-\frac{x^2}6}-1\Bigr)\Bigr),
\end{equation}
where the exponent $-\sqrt Kx/6-\sqrt Kx^3/144-\cdots$ is a power series
in $x$ without constant term; $c_r$ depends only on
$c_0^{\mathrm{II}},\ldots,c_r^{\mathrm{II}}$. For example
\[
 c_1=c_1^{\mathrm{II}}-\frac{\sqrt K}6,\qquad
 c_2=c_2^{\mathrm{II}}-\frac{\sqrt K}6c_1^{\mathrm{II}}+\frac K{72}+\frac14.
\]
At the level of the real expansions the same shift reads
$\sum_jd_jt^j=4\ee^{-t/6}\sum_mb_m^{\mathrm{II}}t^m$, with Part~II's $b_m$
of \eqref{lcp:ao:eq:bdef}: Part~II absorbs the factor $\ee^{-z/6}$ of the eta
quotient into $\ee^{nz}$, since $\ee^{nz-z/6}=\ee^{N^{\mathrm{II}}z}$.
\end{remark}
\begin{proof}
Fix $m$. As $n\to\infty$,
\[
 (N^{\mathrm{II}})^{-(3+s)/2}=n^{-3/2}x^s\Bigl(1-\frac{x^2}6\Bigr)^{-(3+s)/2},\qquad
 \ee^{2\sqrt{KN^{\mathrm{II}}}}=\ee^{2\sqrt{Kn}}
 \exp\bigl(2\sqrt K(\sqrt{N^{\mathrm{II}}}-\sqrt n)\bigr),
\]
with $2\sqrt K(\sqrt{N^{\mathrm{II}}}-\sqrt n)$ equal to the exponent in
\eqref{lcp:tr:eq:shift}. Every factor is analytic at $x=0$, and
$\OO((N^{\mathrm{II}})^{-m/2})=\OO(x^m)$. Hence Theorem~\ref{lcp:ao:thm:main}
with $M=m$ gives $a_n=\delta n^{-3/2}\ee^{2\sqrt{Kn}}(\sum_{r<m}\tilde c_rx^r+\OO(x^m))$,
with $\tilde c_r$ the Taylor coefficients of the right side of
\eqref{lcp:tr:eq:shift}. Theorem~\ref{lcp:tr:thm:main} gives the same
expansion with $c_r$. An asymptotic expansion in powers of $x$ is unique,
so $\tilde c_r=c_r$ for $r<m$, and $m$ is arbitrary. The examples are the
first two Taylor coefficients, using
$(1-x^2/6)^{-3/2}=1+x^2/4+\OO(x^4)$. The last statement follows from
$h=B^{\mathrm{II}}=D^{\mathrm{II}}$ (Proposition~\ref{lcp:tr:prop:T} and
\eqref{lcp:ao:eq:PBD}) and \eqref{lcp:tr:eq:dalgorithm}.
\end{proof}
The formula was also checked symbolically for $r\le4$ (exact SymPy
arithmetic). The values are:
\begin{center}\small
\begin{tabular}{@{}r l l l l@{}}
\toprule
$r$ & $c_r$ (Part III, unshifted) & value & $c_r^{\mathrm{II}}$ (Part II, shifted) & value\\
\midrule
0 & $1$ & $1$ & $1$ & $1$\\
1 & $-\frac{\sqrt K}6-\frac3{2\sqrt K}=-d$ & $-1.03410521762627$ & $-\frac3{2\sqrt K}$ & $-0.64874695296076$\\
2 & $\frac12+\frac3{4K}-\frac{35K}{72}=e$ & $-1.95847649289674$ & $\frac3{4K}-\frac K2$ & $-2.53272698896975$\\
3 & $\frac{755K^{3/2}}{1296}+\frac{175\sqrt K}{72}-\frac5{8\sqrt K}=f$ & $12.5504457791462$ & $\frac{\sqrt K(K+5)}2$ & $11.9607911105086$\\
4 & $\frac5{16K}-\frac{175}{32}-\frac{3775K}{864}-\frac{29375K^2}{31104}$ & $-55.7596598453122$ & $-\frac{45}8-\frac{15K}4-\frac{41K^2}{48}$ & $-50.0848008875789$\\
\bottomrule
\end{tabular}
\end{center}
The manuscript prints $c_1$ and $c_2$ (Theorem~\ref{lcp:tr:thm:main}) and
$b_r=c_r\beta^r$ for $r\le4$ (Section~\ref{lcp:tr:sec:inverse}); the closed
form of $c_4$ above is $b_4/\beta^4$. The column $c_1,c_2,c_3$ is Part~I's
$-d,e,f$ of \eqref{lcp:eq:d}--\eqref{lcp:eq:f}.

\subsection{The manuscript's abstract}
\begin{quote}\small
The area generating function for L-convex polyominoes admits an exact decomposition into an explicit eta quotient times the square of a regular multiplier, together with a positive remainder that is exponentially smaller on the positive real axis. The decomposition follows from a product generating function for an auxiliary continuant recurrence and a duality of its ordinary generating function. It proves a complete real-axis singular expansion. A positive representation with a common finite partition factor supplies uniform minor-arc damping; an elementary interpolation lemma transports the real expansion into the saddle window. Together these arguments prove the coefficient equivalent conjectured by Guttmann and Kot\v e\v sovec and extend it to every algebraic order. We give effective formulas for the coefficients, logarithmic inversion with an integer-threshold qualification, and exact numerical checks against the published coefficient data.
\end{quote}

\section{The enumeration problem and the result}\label{lcp:tr:sec:main}
An L-convex polyomino is a polyomino in which any two cells can be joined inside the polyomino by a path with at most one right-angle bend. We use the fixed area-counting convention of OEIS A126764 and include the empty object. Write
\[
 A(q)=\sum_{n\ge0}a_nq^n=1+q+2q^2+6q^3+15q^4+\cdots.
\]
Castiglione, Frosini, Munarini, Restivo and Rinaldi~\cite{castiglione} proved the exact generating function
\begin{equation}\label{lcp:tr:eq:original}
 A(q)=1+\sum_{k\ge0}\frac{q^{k+1}f_k(q)}{(q;q)_k^2(1-q^{k+1})},
\end{equation}
where
\begin{equation}\label{lcp:tr:eq:frec}
 f_0=1,\qquad f_1=1+2q-q^2,\qquad
 f_k=2f_{k-1}-(1-q^k)^2f_{k-2}.
\end{equation}
Here $(a;q)_n=\prod_{j=0}^{n-1}(1-aq^j)$ and $(a;q)_\infty$ denotes the corresponding convergent infinite product when $|q|<1$.
Guttmann and Kot\v e\v sovec~\cite{gk} numerically conjectured
\begin{equation}\label{lcp:tr:eq:leading}
 a_n\sim\frac{13\sqrt2}{768}\,n^{-3/2}
       \exp\!\left(\pi\sqrt{\frac{13n}{6}}\right).
\end{equation}
Their introduction and the OEIS entry give the constant in \eqref{lcp:tr:eq:leading}; the concluding display in the paper prints its reciprocal. The exact data and the derivation below select \eqref{lcp:tr:eq:leading}.

\begin{writenote}[Added 5 October 2026, batch~101]
The statement about the reciprocal was checked at intake on 5~October 2026
against arXiv:2109.09928v3: the final display of Section~2 of~\cite{gk} is
$3\cdot2^8\ee^{\pi\sqrt{13n/6}}/(13\sqrt2\,n^{3/2})$, whose amplitude
$3\cdot2^8/(13\sqrt2)$ is $1/\delta$, while the decimal constant
quoted in that section is $\delta=0.0239385\ldots$; the paper's notation
$c$ versus $1/c$ around its equation~(22) is inconsistent in the same way.
The amplitude in its equation~(5) and in the OEIS entry is $\delta$. This
also settles the unchecked statement of Part~II (note in
Section~\ref{lcp:ao:sec:main}).
\end{writenote}


Set
\begin{equation}\label{lcp:tr:eq:constants}
 K=\frac{13\pi^2}{24},\qquad C=\frac{1}{4(2\pi)^{3/2}},\qquad
 \delta=\frac{CK}{2\sqrt\pi}=\frac{13\sqrt2}{768}.
\end{equation}
\begin{theorem}\label{lcp:tr:thm:main}
For every fixed positive integer $m$,
\begin{equation}\label{lcp:tr:eq:allcoeff}
 a_n=\delta n^{-3/2}\ee^{2\sqrt{Kn}}
 \left(\sum_{r=0}^{m-1}c_r n^{-r/2}+O_m(n^{-m/2})\right),
 \qquad c_0=1.
\end{equation}
The constants $c_r$ are given effectively by
\begin{equation}\label{lcp:tr:eq:cr}
 c_r=\sum_{j=0}^r d_jK^{j/2}
 \frac{(-1)^{r-j}}{(r-j)!(16\sqrt K)^{r-j}}
 \prod_{\ell=1}^{r-j}\bigl((2j+5)^2-(2\ell-1)^2\bigr),
\end{equation}
where the rational numbers $d_j$ are the coefficients in
\begin{equation}\label{lcp:tr:eq:realexp}
 A(\ee^{-t})=Ct^{3/2}\ee^{K/t}
 \left(\sum_{j=0}^{J-1}d_jt^j+O_J(t^J)\right),\qquad t\downarrow0,
\end{equation}
valid for every fixed $J$. An empty product in \eqref{lcp:tr:eq:cr} is one. In particular,
\[
 c_1=-\frac{\sqrt K}{6}-\frac{3}{2\sqrt K},\qquad
 c_2=-\frac{35K}{72}+\frac12+\frac{3}{4K}.
\]
\end{theorem}
The theorem is an all-orders Poincar\'e expansion of the dominant exponential scale. It does not assert a complete transseries of secondary exponential contributions. Such contributions remain a separate problem.

\begin{writenote}[Added 5 October 2026, batch~101]
Theorem~\ref{lcp:tr:thm:main} is the same theorem as Part~II's
Theorem~\ref{lcp:ao:thm:main}, written in the unshifted index $n$ instead
of $N^{\mathrm{II}}=n-1/6$; Remark~\ref{lcp:tr:rem:shift} converts the
coefficients, and its table shows $c_1,c_2,c_3=-d,e,f$ of Part~I's
Theorem~\ref{lcp:thm:main}. Here $C$ is the real-axis amplitude
$1/(4(2\pi)^{3/2})$, not Part~II's $C$, which is $\delta$. What is new in
Part~III is the proof: Section~\ref{lcp:tr:sec:transfer} carries out the
coefficient transfer in full (Section~\ref{lcp:tr:sec:cross}).
\end{writenote}


\section{An exact product decomposition}\label{lcp:tr:sec:product}
Throughout this section $0<q<1$. Put $p_n=(q;q)_n$, $p_\infty=(q;q)_\infty$, and extend $u_n=f_{n-1}$ by $u_0=1$. Thus
\begin{equation}\label{lcp:tr:eq:urec}
 u_0=u_1=1,\qquad u_{n+1}=2u_n-(1-q^n)^2u_{n-1}.
\end{equation}
Introduce an auxiliary sequence with the same recurrence,
\begin{equation}\label{lcp:tr:eq:Qrec}
 Q_0=1,\quad Q_1=2,\qquad Q_{n+1}=2Q_n-(1-q^n)^2Q_{n-1}.
\end{equation}
These are a specialization of the Al-Salam--Chihara polynomials, but the proof below needs only \eqref{lcp:tr:eq:Qrec}.

\begin{writenote}[Added 5 October 2026, batch~101]
This sentence is proved in Remark~\ref{lcp:tr:rem:asc} at the end of this
section.
\end{writenote}


\subsection{The product generating function and its slope}
Dividing \eqref{lcp:tr:eq:Qrec} by $p_n$ and summing gives
\[
 (1-z)^2P(z)=(1+qz^2)P(qz),\qquad
 P(z)=\sum_{n\ge0}\frac{Q_nz^n}{p_n},\quad P(0)=1.
\]
Iteration therefore yields
\begin{equation}\label{lcp:tr:eq:P}
 P(z)=\frac{(-qz^2;q^2)_\infty}{(z;q)_\infty^2}\qquad(|z|<1).
\end{equation}
Let $V_n=2q^n-q^{2n}$. The recurrence implies
\[
 \Delta Q_{n+1}=\Delta Q_n+V_nQ_{n-1},\qquad \Delta Q_1=1.
\]
Hence $Q_n\ge n+1$ and its increments increase. Also $Q_{n-1}\le n\Delta Q_n$, so
\[
 \Delta Q_{n+1}\le(1+nV_n)\Delta Q_n.
\]
Since $\sum nV_n<\infty$, the increments have a finite limit $D>0$, and
\begin{equation}\label{lcp:tr:eq:Qbounds}
 n+1\le Q_n\le D(n+1),\qquad Q_n/n\longrightarrow D.
\end{equation}
Comparison of the double pole of \eqref{lcp:tr:eq:P} with $Q_n/p_n\sim Dn/p_\infty$ gives
\begin{equation}\label{lcp:tr:eq:D}
 D=\frac{(-q;q^2)_\infty}{p_\infty}.
\end{equation}
This pole comparison is an elementary Abelian statement: if $b_n/n\to L$, then $(1-z)^2\sum b_nz^n\to L$ as $z\uparrow1$.

\subsection{Reduction of order}
A discrete Wronskian calculation, initialized by $u_0/Q_0=1$ and $u_1/Q_1=1/2$, gives
\begin{equation}\label{lcp:tr:eq:reduction}
 \frac{u_n}{Q_n}=1-\sum_{j=0}^{n-1}\frac{p_j^2}{Q_jQ_{j+1}}.
\end{equation}
Define
\begin{equation}\label{lcp:tr:eq:h}
 h(q)=1-\sum_{j\ge0}\frac{p_j^2}{Q_jQ_{j+1}}.
\end{equation}
For every $J\ge0$, \eqref{lcp:tr:eq:Qbounds} gives the useful explicit bound
\begin{equation}\label{lcp:tr:eq:htail}
 0\le\sum_{j\ge J}\frac{p_j^2}{Q_jQ_{j+1}}
 \le p_J^2\sum_{j\ge J}\frac{1}{(j+1)(j+2)}
 =\frac{p_J^2}{J+1}.
\end{equation}
In particular, $u_n/Q_n\to h$. To see strict positivity, observe that $u_n$ is positive and convex, with $\Delta u_1=0$ and $\Delta u_2=V_1>0$. Its increasing increments are bounded because $u_n\le Q_n$ by \eqref{lcp:tr:eq:reduction}. They therefore have a finite positive limit $D_u$, and $h=D_u/D>0$.

\subsection{Duality and the square of the multiplier}
Let
\[
 B_Q(z)=\sum_{n\ge0}\frac{Q_nz^n}{p_n^2},\qquad
 A_Q=B_Q(q)-B_Q(q^2),\qquad
 M=\frac{D}{p_\infty^2}=\frac{(-q;q^2)_\infty}{p_\infty^3}.
\]
The normalized recurrence has generating equation
\begin{equation}\label{lcp:tr:eq:BQ}
 (1-z)^2B_Q(z)=2B_Q(qz)-B_Q(q^2z).
\end{equation}
Its double-pole coefficient at $z=1$ is $M$, so
\begin{equation}\label{lcp:tr:eq:Mboundary}
 M=2B_Q(q)-B_Q(q^2).
\end{equation}
For $r\ge0$, the sequence $w_r=B_Q(q^{r+1})$ satisfies the same recurrence as $Q_r$ and $u_r$. Matching $r=0,1$ in \eqref{lcp:tr:eq:Mboundary} proves
\[
 w_r=M u_r-A_QQ_r.
\]
As $r\to\infty$, $w_r\to1$ whereas $Q_r\to\infty$. Division by $Q_r$ yields
\begin{equation}\label{lcp:tr:eq:duality}
 A_Q=Mh,\qquad u_r-hQ_r=\frac{B_Q(q^{r+1})}{M}.
\end{equation}

\begin{proposition}\label{lcp:tr:prop:factor}
The following decomposition and bounds are exact for $0<q<1$:
\begin{align}
 A(q)&=1+M(q)h(q)^2+R(q),\label{lcp:tr:eq:factor}\\
 R(q)&=\frac{1}{M(q)}\sum_{n\ge1}
       \frac{q^nB_Q(q^{n+1})}{p_{n-1}p_n},\label{lcp:tr:eq:Rformula}\\
 0\le R(q)&\le D(q)\frac{q}{1-q^2},\qquad
 \frac{R(q)}{M(q)}\le\frac{q p_\infty^2}{1-q^2}.\label{lcp:tr:eq:Rbound}
\end{align}
\end{proposition}
\begin{proof}
By \eqref{lcp:tr:eq:original},
$A-1=\sum_{n\ge1}q^n(1-q^n)u_n/p_n^2$.
Substitute \eqref{lcp:tr:eq:duality} and use $A_Q=Mh$ to obtain \eqref{lcp:tr:eq:factor}--\eqref{lcp:tr:eq:Rformula}. Alternatively, the remainder in \eqref{lcp:tr:eq:reduction} and \eqref{lcp:tr:eq:htail} give
\[
 0\le u_n-hQ_n\le\frac{Q_np_n^2}{n+1}\le Dp_n^2.
\]
Multiplying by $q^n(1-q^n)/p_n^2$ and summing gives
$R\le D\sum_{n\ge1}q^n(1-q^n)=Dq/(1-q^2)$.
\end{proof}

\begin{writenote}[Added 5 October 2026, batch~101]
The identity \eqref{lcp:tr:eq:factor} is Part~II's \eqref{lcp:ao:eq:exactmain}:
Part~III's $M$ is Part~II's $P$ (Part~I's $R^{\mathrm I}$), its $h$ is
Part~II's $B=D$ (Part~I's $L^{\mathrm I}/R^{\mathrm I}$), and its $R$ is
Part~II's $R^{\mathrm{II}}-1=E(q)/(q;q)_\infty$; the last equality was checked at
$q=0.3$ and $0.6$ to 40 digits. Part~II obtains the identity by Euler,
Heine and Fine transformations; the duality argument
\eqref{lcp:tr:eq:BQ}--\eqref{lcp:tr:eq:duality}, the positive series
\eqref{lcp:tr:eq:Rformula} and the bounds \eqref{lcp:tr:eq:Rbound} are new.
In Part~I's letters, $M=2C^Q_1-C^Q_2$ is the identity $R^{\mathrm I}=2C^Q_1-C^Q_2$
of Section~4.2, and \eqref{lcp:tr:eq:reduction} is the Wronskian
\eqref{lcp:eq:wronskian} divided by $Q_nQ_{n+1}$.
\end{writenote}


\subsection{An explicit alternating series}
The multiplier also has an explicit alternating-series expression.
\begin{proposition}\label{lcp:tr:prop:T}
For $0<q<1$,
\begin{equation}\label{lcp:tr:eq:T}
 h(q)=\frac{1}{1+q}\sum_{j\ge0}
 \frac{(-1)^j q^{j(j-1)/2}(q;q)_j}{(-q^3;q^2)_j}.
\end{equation}
The absolute values of its summands decrease, and the error after $j=J-1$ is at most the absolute $j=J$ summand, divided by $1+q$.
\end{proposition}
\begin{proof}
Euler's expansion of $(q^{n+1};q)_\infty$ gives
\[
 A_Q=p_\infty^{-1}\sum_{j\ge0}\alpha_j
 \bigl(P(q^{j+1})-P(q^{j+2})\bigr),\qquad
 \alpha_j=\frac{(-1)^jq^{j(j+1)/2}}{p_j}.
\]
All sums are absolutely convergent for fixed $q$. Shifting the second sum and using
$\alpha_j-\alpha_{j-1}=q^{-j}\alpha_j$ gives
\[
 A_Q=\frac{P(q)}{p_\infty}
 \sum_{j\ge0}\frac{(-1)^jq^{j(j-1)/2}p_j}{(-q^3;q^2)_j},
\]
since $P(q^{j+1})/P(q)=p_j^2/(-q^3;q^2)_j$.
Finally $P(q)/p_\infty=M/(1+q)$ and $A_Q=Mh$. Consecutive absolute summands have ratio
\[
 \frac{q^j(1-q^{j+1})}{1+q^{2j+3}}<1,
\]
which proves the remainder assertion.
\end{proof}

\begin{writenote}[Added 5 October 2026, batch~101]
Since $(-q;q^2)_{j+1}=(1+q)(-q^3;q^2)_j$, the series \eqref{lcp:tr:eq:T} is
Part~II's series $D$ of \eqref{lcp:ao:eq:PBD}, so
Proposition~\ref{lcp:tr:prop:T} gives $h=D^{\mathrm{II}}$; with Part~II's
Proposition~\ref{lcp:ao:prop:square} and Remark~\ref{lcp:ao:rem:LR}, $h$ is
the partial theta function \eqref{lcp:ao:eq:theta} and Part~I's
$D^{\mathrm I}=L^{\mathrm I}/R^{\mathrm I}$. The proof here (Euler's expansion of
$A_Q$) differs from Part~II's (Fine's transformation). The alternating error
bound is new.
\end{writenote}

\begin{remark}[write, 5 October 2026: the Al-Salam--Chihara specialization]\label{lcp:tr:rem:asc}
In the normalization of Koekoek, Lesky and Swarttouw~\cite[\S14.8]{lcp-tr-kls},
the Al-Salam--Chihara polynomials $Q_n(x;a,b\,|\,q)$ satisfy $Q_0=1$,
$Q_1=2x-(a+b)$ and
\[
 2xQ_n=Q_{n+1}+(a+b)q^nQ_n+(1-q^n)(1-abq^{n-1})Q_{n-1}\qquad(n\ge1).
\]
For $0<q<1$ and either choice of $\sqrt q$,
\[
 Q_n(q)=Q_n\bigl(1;\ii\sqrt q,-\ii\sqrt q\,\big|\,q\bigr)\qquad(n\ge0).
\]
Indeed, with $a=\ii\sqrt q$ and $b=-\ii\sqrt q$ one has $a+b=0$ and $ab=q$,
so $(1-q^n)(1-abq^{n-1})=(1-q^n)^2$; at $x=1$ the recurrence becomes
$Q_{n+1}=2Q_n-(1-q^n)^2Q_{n-1}$ with $Q_0=1$ and $Q_1=2$, which is
\eqref{lcp:tr:eq:Qrec}. Both sequences are determined by these data, so they
coincide. (The recurrence above is the Askey--Wilson recurrence with $c=d=0$,
rescaled by $2^n$ from its monic form.) As a check independent of the
recurrence, the basic hypergeometric definition
$Q_n(x;a,b\,|\,q)=a^{-n}(ab;q)_n\,{}_3\phi_2(q^{-n},a\ee^{\ii\theta},a\ee^{-\ii\theta};ab,0;q,q)$,
$x=\cos\theta$, was evaluated at $x=1$, $q=0.37$, $0\le n\le9$ to 40 digits
and agrees with \eqref{lcp:tr:eq:Qrec}. The point $x=1$ is an endpoint of
the orthogonality interval. The identification concerns $Q_n$ as a
function of $x$; it says nothing yet about the zeros of $Q_m$ as a
polynomial in $q$, which Part~I's research questions~1 and~4 ask about.
The solution $u_n$ of the same recurrence, with $u_1=1$ instead of $2$, is
not covered by this identification.
\end{remark}


\section{The complete real axis expansion}\label{lcp:tr:sec:real}
For fixed $J$, the bound \eqref{lcp:tr:eq:htail} is $O((1-q)^{2J})$. Each finite term is rational and regular at $q=1$, because $Q_j(1)=2^j$. Thus \eqref{lcp:tr:eq:h} proves a complete expansion of $h$ in integer powers of $1-q$. The alternating formula \eqref{lcp:tr:eq:T} provides a second algorithm, with remainder $O((1-q)^J)$ after $J$ terms.

The product $M$ is the eta quotient
\begin{equation}\label{lcp:tr:eq:eta}
 M(q)=\frac{(q^2;q^2)_\infty^2}
 {(q;q)_\infty^4(q^4;q^4)_\infty}.
\end{equation}
The standard modular product transformation, equivalently the modular inversion of Dedekind's eta function~\cite{lcp-tr-dlmf}, yields
\begin{equation}\label{lcp:tr:eq:Mexp}
 M(\ee^{-t})=\left(\frac{t}{2\pi}\right)^{3/2}
 \exp\!\left(\frac{K}{t}-\frac{t}{6}\right)
 \left(1+O(\ee^{-\pi^2/t})\right).
\end{equation}
Indeed $(\ee^{-at};\ee^{-at})_\infty
=\sqrt{2\pi/(at)}\exp(-\pi^2/(6at)+at/24)
(1+O(\ee^{-4\pi^2/(at)}))$ for fixed $a>0$.
The bound \eqref{lcp:tr:eq:Rbound} shows that $R/M$ is exponentially small; $1/M$ is likewise exponentially small. Hence \eqref{lcp:tr:eq:factor} proves \eqref{lcp:tr:eq:realexp}, with the explicit formal coefficient rule
\begin{equation}\label{lcp:tr:eq:dalgorithm}
 \sum_{j\ge0}d_jt^j\;\sim\;4\ee^{-t/6}h(\ee^{-t})^2.
\end{equation}
This rule means finite asymptotic truncation, rather than convergence of the displayed formal series.
For $d=1-q$, the first terms are
\begin{align*}
 h(q)={}&\frac12-\frac{d^2}{8}-\frac{5d^4}{32}
 +\frac{d^5}{8}-\frac{65d^6}{128}
 +\frac{37d^7}{32}-\frac{2053d^8}{512}+O(d^9),\\
 4h(\ee^{-t})^2={}&1-\frac{t^2}{2}+\frac{t^3}{2}
 -\frac{41t^4}{48}+\frac{7t^5}{4}
 -\frac{6317t^6}{1440}+\frac{1529t^7}{120}
 -\frac{1702513t^8}{40320}+O(t^9).
\end{align*}
The normalized coefficients are
\begin{equation}\label{lcp:tr:eq:dvalues}
\begin{aligned}
 d_0&=1,&d_1&=-\frac16,&d_2&=-\frac{35}{72},\\
 d_3&=\frac{755}{1296},&d_4&=-\frac{29375}{31104},&
 d_5&=\frac{1772639}{933120},\\
 d_6&=-\frac{31514551}{6718464},&d_7&=\frac{3808743271}{282175488},&
 d_8&=-\frac{85931514901}{1934917632}.
\end{aligned}
\end{equation}
A real-axis expansion by itself does not justify pointwise coefficient extraction. The next section supplies the missing analytic control explicitly.

\begin{writenote}[Added 5 October 2026, batch~101]
This section re-proves Part~I's Theorem~\ref{lcp:thm:radial} by a third
route: $\sum_jd_jt^j=\ee^{-t/6}\mathcal F(t)$ with $\mathcal F$ of
\eqref{lcp:eq:Fintro}, and the expansion of $h$ in $d=1-q$ is Part~I's
expansion of $\mathcal D$ in $w=1-q$ (Section~8.1), coefficient for
coefficient. Part~II's coefficients \eqref{lcp:ao:eq:Dexp} are those of
$h(\ee^{-z})$ and $h(\ee^{-z})^2$. The expansion \eqref{lcp:tr:eq:Mexp} is
\eqref{lcp:eq:Rlocal} and \eqref{lcp:ao:eq:Plocal}, here with the explicit
remainder exponent $\pi^2/t$ from the factor $(q^4;q^4)_\infty$.
\end{writenote}


\section{From real asymptotics to coefficients}\label{lcp:tr:sec:transfer}
The proof uses positivity, a common finite partition factor, and local analytic interpolation. It supplies the needed analytic control directly; the pitfalls of omitting complex hypotheses from Tauberian arguments are discussed in~\cite{lcp-tr-bjm}.
\subsection{A positive representation with a common partition
factor}\label{lcp:tr:a-positive-representation-with-a-common-partition-factor}

Put \(u_0=u_1=1\) and \(u_n=f_{n-1}\) for \(n\ge1\). The defining
recurrence becomes

\[
u_n=2u_{n-1}-(1-q^{n-1})^2u_{n-2}\quad(n\ge2).
\]

Define \(v_n=u_n/(q;q)_n\), \(v_0=1\), and \(w_n=v_n-v_{n-1}\). Then

\[
(1-q^n)v_n=2v_{n-1}-(1-q^{n-1})v_{n-2},
\]

and therefore

\begin{equation}\label{lcp:tr:tr:2}
(1-q^n)w_n=(1+q^n)w_{n-1}+(q^n+q^{n-1})v_{n-2}.
\end{equation}

Since \(w_1=q/(1-q)\), induction proves that all \(v_n\) and \(w_n\)
have nonnegative coefficients. The exact area generating function is

\begin{equation}\label{lcp:tr:tr:3}
A(q)=1+\sum_{n\ge1}\frac{q^nv_n(q)}{(q;q)_{n-1}}.
\end{equation}

For any integer \(N\), split \eqref{lcp:tr:tr:3} into \(L_N\), containing its
constant term and \(n\le N\), and \(H_N\), containing \(n\ge N+1\). Both
have nonnegative coefficients, and

\begin{equation}\label{lcp:tr:tr:4}
H_N(q)=\frac{G_N(q)}{(q;q)_N},\qquad G_N\text{ has nonnegative coefficients}.
\end{equation}

All series here converge in the open unit disk. For completeness, put
\(a_j=2r^j-r^{2j}\) for a fixed real \(r<1\). Summing the recurrence
twice gives \[
u_n(r)=1+\sum_{j=1}^{n-1}(n-j)a_ju_{j-1}(r).
\] For \(m_n=u_n(r)/(n+1)\), it follows that
\(m_n\le1+\sum_{j=1}^{n-1}ja_jm_{j-1}\). Discrete Gronwall bounds this
by \(\prod_{j\ge1}(1+ja_j)<\infty\). Thus \(u_n(r)=O_r(n+1)\), giving
convergence of \eqref{lcp:tr:tr:3}. Nonnegative coefficients give absolute
convergence on smaller complex disks.

\subsection{The low-index part is exponentially
smaller}\label{lcp:tr:the-low-index-part-is-exponentially-smaller}

Take \(r=e^{-t}\) and \(N=\lfloor1/t\rfloor\). Positivity and the
original recurrence give \(0<u_n(r)\le2^n\). Hence

\[
L_N(r)\le 1+N\frac{2^N}{(r;r)_N^2}.
\]

The elementary bound \(1-e^{-x}\ge xe^{-x}\) gives

\[
(r;r)_N\ge t^NN!e^{-tN(N+1)/2}.
\]

Using \(N!\ge(N/e)^N\) and \(tN=1+O(t)\), we obtain

\begin{equation}\label{lcp:tr:tr:5}
L_N(r)=O\!\left(t^{-1}\exp\left(\frac{E}{t}\right)\right),
 \qquad E=3+\log2<K=\frac{13\pi^2}{24}.
\end{equation}

Changing the absolute constant accommodates the harmless
\(O(1)\) in the exponential. Because \(L_N\) is
coefficientwise nonnegative, the same upper bound applies to
\(|L_N(re^{i\theta})|\).

\subsection{Uniform minor-arc damping}\label{lcp:tr:uniform-minor-arc-damping}

From \eqref{lcp:tr:tr:4},

\begin{equation}\label{lcp:tr:tr:6}
|H_N(re^{i\theta})|\le H_N(r)\,D_N(t,\theta),\qquad
D_N(t,\theta)=\prod_{j=1}^N\frac{1-r^j}{|1-r^je^{ij\theta}|}.
\end{equation}

There is an absolute \(c>0\), independent of sufficiently small \(t\),
such that

\begin{equation}\label{lcp:tr:tr:7}
D_N(t,\theta)\le\exp\left[-c\min\left(\frac{\theta^2}{t^3},\frac1t\right)\right]
\quad(-\pi\le\theta\le\pi).
\end{equation}

Here is a proof. Restrict the logarithm of the product to
\(N/2\le j\le N\). On this range \(r^j\) is bounded above and below by
fixed constants strictly between zero and one. Thus

\[
\log D_N\le-c_1\sum_{N/2\le j\le N}(1-\cos j\theta).
\]

The sum is bounded below by \(c_2N\min\{1,N^2\theta^2\}\):

\begin{itemize}[itemsep=0pt]
\item
  If \(|\theta|\le1/N\), use \(1-\cos x\gg x^2\) for \(|x|\le1\).
\item
  For \(1/N\le|\theta|\le C_0/N\), put \(\lambda=N|\theta|\). Uniform
  Riemann-sum convergence on the compact interval
  \(1\le\lambda\le C_0\), and strict positivity of
  \(\int_{1/2}^1(1-\cos\lambda x)\,dx\), give the claim.
\item
  For \(C_0/N\le|\theta|\le\pi\), the geometric-sum bound
  \(|\sum e^{ij\theta}|\le1/|\sin(\theta/2)|\le\pi/|\theta|\), with a
  sufficiently large fixed \(C_0\), gives the claim.
\end{itemize}

Consequently,

\begin{equation}\label{lcp:tr:tr:8}
|A(re^{i\theta})|\le L_N(r)+A(r)\exp\left[-c\min\left(\frac{\theta^2}{t^3},\frac1t\right)\right].
\end{equation}

This is the needed quantitative exclusion of secondary arcs, including
neighborhoods of other roots of unity.

\begin{writenote}[Added 5 October 2026, batch~101]
Compare Part~I's \eqref{lcp:eq:minor} and Part~II's
Lemma~\ref{lcp:ao:lem:minor}: the same three-range argument, there for the
infinite product $R^{\mathrm I}=P^{\mathrm{II}}$, here for the finite
partition factor $(q;q)_N$ common to every term of $H_N$. Together with
the trivial bound $|L_N(r\ee^{\ii\theta})|\le L_N(r)$, the bound
\eqref{lcp:tr:tr:8} needs no estimate of $h$, of $M$ or of the remainder
$R$ off the real axis; it replaces Part~II's whole-circle
Lemma~\ref{lcp:ao:lem:global}. In \eqref{lcp:tr:tr:7} the three ranges are
complete as written: on $N/2\le j\le N$ one has $\ee^{-1}\le r^j\le
\ee^{-(1-t)/2}$, and every factor of $D_N$ is at most one, since
$|1-r^j\ee^{\ii j\theta}|^2=(1-r^j)^2+2r^j(1-\cos j\theta)$.
\end{writenote}


\subsection{Real all-orders expansions propagate to the entire saddle
window}\label{lcp:tr:real-all-orders-expansions-propagate-to-the-entire-saddle-window}

We use the following elementary analytic lemma.

\textbf{Interpolation lemma.} If \(F\) is holomorphic on a neighborhood
of \(|w|\le4\), \(|F(w)|\le M\) there, and \(|F(x)|\le\varepsilon\) for
\(-1\le x\le1\), then, for every positive integer \(k\),

\begin{equation}\label{lcp:tr:tr:9}
\sup_{|w|\le1}|F(w)|\le4^{k-1}\varepsilon+\frac43M\left(\frac23\right)^k.
\end{equation}

To prove this, interpolate at the \(k\) roots of the Chebyshev
polynomial \(T_k\). Let \(\omega_k\) be the monic polynomial with those
roots. On \(|w|\le1\), \(|\omega_k(w)|\le2^k\); on \(|\zeta|=4\),
\(|\omega_k(\zeta)|\ge3^k\). The contour remainder formula gives the
second term in \eqref{lcp:tr:tr:9}. Moreover
\(|\omega_k'(x_j)|=2^{1-k}|T_k'(x_j)|\ge2^{1-k}k\); bounding the other
\(k-1\) numerator factors by 2 gives a total Lagrange interpolation norm
at most \(4^{k-1}\), which proves the first term.

Fix any positive constant \(B\), and set

\[
R_t=t^{3/2}\sqrt{B\log(1/t)},\qquad z=t+R_tw.
\]

For any fixed truncation index \(L\), apply \eqref{lcp:tr:tr:9} to

\[
F_t(w)=C^{-1}z^{-3/2}e^{-K/z}A(e^{-z})-\sum_{j=0}^{L-1}d_jz^j.
\]

The powers use the branch analytic in the right half-plane. Nonnegative
coefficients imply \(|A(e^{-z})|\le A(e^{-\Re z})\). For \(z=x+iy\) with
\(|w|\le4\), \(x\asymp t\), and the real leading asymptotic in
\eqref{lcp:tr:eq:realexp} gives

\begin{equation}\label{lcp:tr:tr:10}
|z^{-3/2}e^{-K/z}A(e^{-z})|
\ll\exp\left(\frac{Ky^2}{x(x^2+y^2)}\right)
\le\exp(O(R_t^2/t^3))=t^{-O(B)}.
\end{equation}

The polynomial subtraction remains bounded for fixed \(L\). On the real
interval \(-1\le w\le1\), \eqref{lcp:tr:eq:realexp} gives \(F_t(w)=O_L(t^L)\),
uniformly. In \eqref{lcp:tr:tr:9}, choose \(k=\lceil D\log(1/t)\rceil\), with
\(D\) sufficiently large; then choose \(L\) still larger. For any
prescribed \(P\), this makes both terms \(O(t^P)\). Dropping the excess
polynomial terms proves, for every fixed \(P\),

\begin{equation}\label{lcp:tr:tr:11}
A(e^{-z})=Cz^{3/2}e^{K/z}
\left(\sum_{j=0}^{P-1}d_jz^j+O_P(t^P)\right)
\end{equation}

uniformly on \(|z-t|\le R_t\). Constants may depend on \(B,P\), which is
exactly what is needed. This step only requires a complete real-axis
expansion, analyticity, and positivity.

\subsection{Cauchy saddle and all-orders
conclusion}\label{lcp:tr:cauchy-saddle-and-all-orders-conclusion}

Set \(t=\sqrt{K/n}\). Cauchy's formula reads

\[
a_n=\frac1{2\pi}\int_{-\pi}^{\pi}A(e^{-t+i\theta})e^{nt-in\theta}\,d\theta.
\]

For any requested algebraic accuracy, choose \(B\) sufficiently large.
Equations \eqref{lcp:tr:tr:5} and \eqref{lcp:tr:tr:8} show that \(|\theta|>R_t\)
contributes less than the main scale by that power of \(t\). Indeed the
high-index part has factor at most \(e^{-cB\log(1/t)}\), while the
low-index part is exponentially smaller by \(e^{-(K-E)/t}\). The main
scale is \(t^3e^{2K/t}\), up to a fixed constant, and the crude length
bound loses only the fixed power \(t^{-3/2}\) relative to
\(A(r)e^{nt}\).

On \(|\theta|\le R_t\), \eqref{lcp:tr:tr:11} applies. In particular, \[
|e^{K/z+nz}|=\exp\left(\frac{2K}{t}-\frac{K\theta^2}{t(t^2+\theta^2)}\right)
\le e^{2K/t}e^{-K\theta^2/(2t^3)}
\] for all sufficiently small \(t\). Thus a bracket error \(O(t^P)\) in
\eqref{lcp:tr:tr:11} integrates to an error \(O(t^P)\) relative to the main
coefficient scale, without a logarithmic loss. Expanding the saddle
phase \(K/z+nz\) about \(z=t\), with \(\theta=t^{3/2}u\), gives Gaussian
factor \(e^{-Ku^2}\). Taylor expansion to arbitrary fixed order and
Gaussian tail bounds justify integration term by term. Equivalently, the
contribution of each \(C d_j z^{3/2+j}e^{K/z}\) is represented
asymptotically by

\begin{equation}\label{lcp:tr:tr:12}
C d_j\left(\frac Kn\right)^{5/4+j/2}I_{-5/2-j}(2\sqrt{Kn}).
\end{equation}

Formula \eqref{lcp:tr:tr:12} is an asymptotic shorthand for the saddle
expansion, rather than a claim of an exact coefficient identity. The
standard large-positive-argument Bessel expansion~\cite{lcp-tr-dlmf} yields

\begin{equation}\label{lcp:tr:tr:13}
a_n\sim\frac{CK}{2\sqrt\pi}\,n^{-3/2}e^{2\sqrt{Kn}}
\sum_{r\ge0}c_rn^{-r/2},\qquad c_0=1,
\end{equation}

where

\begin{equation}\label{lcp:tr:tr:14}
c_r=\sum_{j=0}^r d_jK^{j/2}
\frac{(-1)^{r-j}}{(r-j)!(16\sqrt K)^{r-j}}
\prod_{\ell=1}^{r-j}\left((2j+5)^2-(2\ell-1)^2\right).
\end{equation}

An empty product is 1. Since the Bessel orders are half-integral, the
product is zero whenever \(r-j>j+2\); this can simplify calculation. For
example,

\[
c_1=d_1\sqrt K-\frac{3}{2\sqrt K},\qquad
c_2=d_2K-3d_1+\frac{3}{4K}.
\]

Precisely, after retaining \(r=0,\ldots,m-1\), the error in the bracket
in \eqref{lcp:tr:tr:13} is \(O_m(n^{-m/2})\). Finally,

\[
\frac{CK}{2\sqrt\pi}=\frac{13\sqrt2}{768},\qquad
2\sqrt{Kn}=\pi\sqrt{13n/6},
\]

which proves Theorem~\ref{lcp:tr:thm:main}.

\begin{writenote}[Added 5 October 2026, batch~101]
The step from the saddle window to \eqref{lcp:tr:tr:13} is stated above as
``Taylor expansion \ldots\ and Gaussian tail bounds'' and, equivalently, by
the shorthand \eqref{lcp:tr:tr:12}. Lemma~\ref{lcp:tr:lem:hankel} and the
paragraph after it carry it out exactly: on the window, each term is a
Hankel integral equal to the Bessel function in \eqref{lcp:tr:tr:12}, up to
an error that is a large power of $t$ once the constant $B$ is large. This
is Part~II's Hankel--Bessel step \eqref{lcp:ao:eq:Bessel}, with the contour
deformation written out.
\end{writenote}

\begin{lemma}[write, 5 October 2026: the saddle segment is part of a Hankel contour]\label{lcp:tr:lem:hankel}
Let $\alpha$ be real, $n\ge K$, $t=\sqrt{K/n}$, $\lambda=\sqrt{Kn}=K/t$
(so $\lambda\ge1$), and $0<\rho\le1$. Let $S$ be the segment $z=t+\ii y$,
$|y|\le\rho t$, oriented upwards. Then
\begin{equation}\label{lcp:tr:eq:hankel}
 \frac1{2\pi\ii}\int_Sz^\alpha\ee^{K/z+nz}\dd z
 =t^{\alpha+1}\bigl(I_{-\alpha-1}(2\lambda)+\varepsilon\bigr),\qquad
 |\varepsilon|\le2^{(|\alpha|+1)/2}\ee^{2\lambda-\lambda\rho^2/2}
 +\Gamma(|\alpha|+1)\,\ee^{1-\lambda}.
\end{equation}
\end{lemma}
\begin{proof}
Substitute $z=ts$. Since $K/z=\lambda/s$ and $nz=\lambda s$, the left side
of \eqref{lcp:tr:eq:hankel} is $t^{\alpha+1}\Phi_S$ with
$\Phi_\Gamma=\frac1{2\pi\ii}\int_\Gamma s^\alpha\ee^{\lambda(s+1/s)}\dd s$ and
$S$ now the segment $s=1+\ii\sigma$, $|\sigma|\le\rho$; powers are principal
in the plane cut along $(-\infty,0]$. Put $\rho_1=\sqrt{1+\rho^2}$ and
$\psi_0=\arctan\rho$, so that the ends of $S$ are $\rho_1\ee^{\pm\ii\psi_0}$.
Let $\mathcal H$ run from $-\infty$ to $-\rho_1$ along the lower edge of the
cut, along $|s|=\rho_1$ from $\arg s=-\pi$ to $-\psi_0$, up $S$, along
$|s|=\rho_1$ from $\psi_0$ to $\pi$, and back to $-\infty$ along the upper
edge. On $\mathcal H$ one has $|s|\ge1$, so
$\ee^{\lambda/s}=\sum_{k\ge0}\lambda^ks^{-k}/k!$ converges uniformly there,
and $\int_{\mathcal H}|s^{\alpha-k}\ee^{\lambda s}||\dd s|<\infty$. Integrating
term by term and using Hankel's formula
$\frac1{2\pi\ii}\int_{\mathcal H}s^{-w}\ee^{\lambda s}\dd s=\lambda^{w-1}/\Gamma(w)$
(for $\lambda>0$ and every complex $w$; the contour may be deformed to the
standard one since the integrand decays on the edges),
\[
 \Phi_{\mathcal H}=\sum_{k\ge0}\frac{\lambda^k}{k!}\,
 \frac{\lambda^{k-\alpha-1}}{\Gamma(k-\alpha)}
 =\sum_{k\ge0}\frac{\lambda^{2k-\alpha-1}}{k!\,\Gamma(k-\alpha)}
 =I_{-\alpha-1}(2\lambda),
\]
by the power series of $I_\nu$ with $\nu=-\alpha-1$ (terms with
$1/\Gamma(k-\alpha)=0$ vanish on both sides). It remains to bound the
pieces of $\mathcal H$ outside $S$. On the arcs $s=\rho_1\ee^{\ii\psi}$,
$\psi_0\le|\psi|\le\pi$,
\[
 \Re\Bigl(s+\frac1s\Bigr)=\Bigl(\rho_1+\frac1{\rho_1}\Bigr)\cos\psi
 \le\Bigl(\rho_1+\frac1{\rho_1}\Bigr)\frac1{\rho_1}
 =2-\frac{\rho^2}{1+\rho^2}\le2-\frac{\rho^2}2,
\]
$|s^\alpha|\le\rho_1^{|\alpha|}\le2^{|\alpha|/2}$, and the two arcs have total
length at most $2\pi\rho_1\le2\sqrt2\,\pi$; after division by $2\pi$ they
contribute at most $2^{(|\alpha|+1)/2}\ee^{\lambda(2-\rho^2/2)}$. On the
edges $s=x\ee^{\pm\ii\pi}$, $x\ge\rho_1\ge1$, the modulus of the integrand is
$x^\alpha\ee^{-\lambda(x+1/x)}\le x^{|\alpha|}\ee^{-x}\ee^{-(\lambda-1)x}
\le x^{|\alpha|}\ee^{-x}\ee^{1-\lambda}$; the two edges contribute at most
$\frac1\pi\Gamma(|\alpha|+1)\ee^{1-\lambda}$. Hence
$|\Phi_S-I_{-\alpha-1}(2\lambda)|$ is bounded as stated.
\end{proof}

\paragraph{The extraction step, completed (write, 5 October 2026).}
Take $\rho=R_t/t=\sqrt{tB\log(1/t)}$, which is at most one for small $t$.
Then $\lambda\rho^2=KB\log(1/t)$, so the error in \eqref{lcp:tr:eq:hankel}
is at most $c_\alpha t^{KB/2}\ee^{2\lambda}$, while
$I_{-\alpha-1}(2\lambda)\asymp\ee^{2\lambda}\lambda^{-1/2}\asymp\ee^{2\lambda}t^{1/2}$.
In Cauchy's formula the window $|\theta|\le R_t$ is the segment $S$
(with $z=t-\ii\theta$ and $a_n=\frac1{2\pi\ii}\int A(\ee^{-z})\ee^{nz}\dd z$
along the upward segment). By \eqref{lcp:tr:tr:11} with $P$ terms, the
window contributes
\[
 \sum_{j<P}Cd_j\,t^{5/2+j}\bigl(I_{-5/2-j}(2\lambda)+\varepsilon_j\bigr)
 +\OO\bigl(t^{P+3}\ee^{2K/t}\bigr),
\]
where the last term is the integrated bracket error, by the Gaussian bound
for $|\ee^{K/z+nz}|$ displayed above and $|z|^{3/2}\le2t^{3/2}$. The
$\varepsilon_j$ contribute $\OO(t^{3+KB/2-1/2}\ee^{2K/t})$, which is
$\OO(t^{P+3}\ee^{2K/t})$ once $KB/2\ge P+1/2$. The main scale is
$t^3\ee^{2K/t}$. Finally, for half-integral order the expansion
\cite[10.40.1]{lcp-tr-dlmf} terminates:
$I_{-5/2-j}(x)=\ee^x(2\pi x)^{-1/2}\bigl(\sum_{k=0}^{j+2}(-1)^ka_k x^{-k}+\OO(x^{-J})\bigr)$
for every $J$, with $a_k=\prod_{\ell=1}^k((2j+5)^2-(2\ell-1)^2)/(k!\,8^k)$
\cite[10.17.1]{lcp-tr-dlmf}. With $x=2\lambda$ and
$t^{5/2+j}\lambda^{-1/2}=K^{1+j/2}n^{-3/2-j/2}$, the $j$th term is
\[
 \frac{CK}{2\sqrt\pi}\,n^{-3/2}\ee^{2\sqrt{Kn}}\,d_jK^{j/2}n^{-j/2}
 \sum_{k=0}^{j+2}\frac{(-1)^k\prod_{\ell=1}^k((2j+5)^2-(2\ell-1)^2)}{k!\,(16\sqrt K)^k}\,n^{-k/2}
 \bigl(1+\OO(n^{-J})\bigr).
\]
Collecting $j+k=r$ gives \eqref{lcp:tr:tr:14}; the terms with $r\ge P$ and
all errors are $\OO(t^P)=\OO(n^{-P/2})$ relative to the main scale. This is
Theorem~\ref{lcp:tr:thm:main} with $m=P$.



% Self-contained mathematical contribution.  Requires amsmath/amssymb.
% A(q) is the generating function; \mathcal A is a constant.
\section{Coefficient expansion and exact inversion}\label{lcp:tr:sec:inverse}
Put
\[
 K=\frac{13\pi^2}{24},\qquad \beta=2\sqrt K=\pi\sqrt{\frac{13}{6}},
 \qquad \delta=\frac{13\sqrt2}{768},\qquad z=\beta\sqrt n,
 \qquad \mathcal A=\delta\beta^3=\frac{169\sqrt{39}\,\pi^3}{13824}.
\]
The coefficient-transfer theorem, applied to the all-orders real-axis expansion,
gives, for every fixed integer $R\ge0$,
\begin{equation}\label{lcp:tr:eq:coef-full}
 a_n=\delta n^{-3/2}e^{\beta\sqrt n}
 \left(\sum_{r=0}^{R}\frac{b_r}{z^r}+O(z^{-R-1})\right).
\end{equation}
Here $b_0=1$.  If $d_j$ denotes the coefficient of $t^j$ in the normalized
real-axis expansion, the finite coefficient rule is
\begin{equation}\label{lcp:tr:eq:b-rule}
 b_r=\sum_{j=\max(0,\lceil(r-2)/2\rceil)}^r
 d_j(2K)^j\frac{(-1)^{r-j}(r+2)!}
 {2^{r-j}(r-j)!(2j+2-r)!}.
\end{equation}
In particular,
\begin{align*}
 b_1&=-3-\frac K3,\\
 b_2&=3+2K-\frac{35}{18}K^2,\\
 b_3&=-5K+\frac{175}{9}K^2+\frac{755}{162}K^3,\\
 b_4&=5K-\frac{175}{2}K^2-\frac{3775}{54}K^3
             -\frac{29375}{1944}K^4.
\end{align*}
Thus, in the alternative convention with powers $n^{-1/2}$, the correction
coefficient is $c_r=b_r/\beta^r$.

For completeness, the rule follows from the dominant saddle contributions
\[
 C\sum_j d_j(K/n)^{(j+5/2)/2}I_{-j-5/2}(2\sqrt{Kn}),
 \qquad C=\frac1{8\sqrt2\,\pi^{3/2}}.
\]
For $m=j+2$, the dominant exponential part of the half-integral Bessel expansion is
\[
 I_{-m-1/2}(z)\sim \frac{e^z}{\sqrt{2\pi z}}
 \sum_{k=0}^{m}\frac{(-1)^k(m+k)!}{k!(m-k)!(2z)^k}.
\]
Collecting $j+k=r$ gives \eqref{lcp:tr:eq:b-rule}. The leading constant is
$CK/(2\sqrt\pi)=\delta$. These Bessel expressions describe the dominant saddle:
they do not claim a convergent infinite Bessel sum or determine exponentially
small contributions from other parts of the contour.  The coefficient-transfer
proof is needed to establish \eqref{lcp:tr:eq:coef-full}; a positive real-axis asymptotic
by itself would not suffice. An independent Gaussian-saddle expansion checks
\eqref{lcp:tr:eq:b-rule} through order $z^{-4}$ in the supplied symbolic code.

\subsection{The logarithmic expansion}
Define $\lambda_r$ by the formal series
\[
 \log\left(1+\sum_{r\ge1}b_r u^r\right)=\sum_{r\ge1}\lambda_r u^r.
\]
A finite recursion, requiring no convergence assumption, is
\[
 \lambda_r=b_r-\frac1r\sum_{k=1}^{r-1}k\lambda_k b_{r-k}.
\]
The first four coefficients are
\begin{align*}
 \lambda_1&=-3-\frac K3,\\
 \lambda_2&=-2K^2+K-\frac32,\\
 \lambda_3&=4K^3+\frac{251}{18}K^2-K,\\
 \lambda_4&=-\frac{47}{3}K^4-\frac{148}{3}K^3
             -\frac{125}{3}K^2-K+\frac94.
\end{align*}
For every fixed $R$ this gives
\begin{equation}\label{lcp:tr:eq:log-full}
 \log a_n=z-3\log z+\log\mathcal A
             +\sum_{r=1}^{R}\frac{\lambda_r}{z^r}+O(z^{-R-1}).
\end{equation}

\subsection{Strict increase and the meaning of an inverse}
We count L-convex polyominoes up to translation. L-convexity implies column
convexity: two cells in the same column can be joined by a path with at most one
turn only if the intervening vertical segment lies in the polyomino. The same
argument gives row convexity, so each occupied row has an interval of columns.
These row intervals are pairwise nested. Indeed, if two were incomparable,
choose a cell in each interval lying outside the other. Neither of the two
possible one-turn corner cells is present, contradicting L-convexity.

Choose the topmost widest row and append one cell to its right end. Nestedness
means that the chosen row meets every occupied column. The new cell can therefore
be joined to any old cell by following the chosen row and then the old cell's
column. Column convexity proves that this path stays inside. Old connections
are unchanged, so the result remains L-convex. The added cell is the sole cell
in a new rightmost column; deleting that column recovers the original
translation class. This is an injection from area $n$ to area $n+1$. The vertical
column of area $n+1$ is outside its image. Consequently
\[
 a_{n+1}\ge a_n+1\quad(n\ge1).
\]
The empty polyomino is a separate initial case, with $a_0=a_1=1$.
Thus for $y>1$ in the exact range there is a unique positive integer $n$ with
$a_n=y$. For arbitrary real $y>1$ define the threshold
\[
 N(y)=\min\{n\ge1:a_n\ge y\}.
\]
Exact-range inversion and this step-valued threshold must be distinguished.

\begin{writenote}[Added 5 October 2026, batch~101]
Part~II's Corollary~\ref{lcp:ao:cor:brackets} uses the same injection,
mirrored (a cell adjoined on the left of the row), and concludes only that
$a_n$ is nondecreasing, with strict increase for large $n$ from the
asymptotics. The observation that the vertical bar is not in the image
sharpens this to $a_{n+1}\ge a_n+1$ for every $n\ge1$. In the shipped data
$a_{n+1}-a_n\ge1$ for $1\le n<8000$, with equality only at $n=1$
($a_2=2=a_1+1$).
\end{writenote}


\subsection{A Lambert model and its corrections}
The leading smooth model is $y=\mathcal A e^z/z^3$.  On its increasing branch,
\[
 z_0(y)=-3W_{-1}\left(-\frac13(\mathcal A/y)^{1/3}\right),
 \qquad n_0(y)=\frac{z_0(y)^2}{4K}.
\]
It is real for $y\ge\mathcal A e^3/27$, and $z_0\ge3$.
On the exact range $y=a_n\to\infty$, inversion of \eqref{lcp:tr:eq:log-full} gives
\begin{align*}
 z={}&z_0-\frac{\lambda_1}{z_0}
       -\frac{3\lambda_1+\lambda_2}{z_0^2}
       -\frac{\lambda_1^2+9\lambda_1+3\lambda_2+\lambda_3}{z_0^3}
       +O(z_0^{-4}),\\
 n={}&\frac1{4K}\left[z_0^2-2\lambda_1
       -\frac{6\lambda_1+2\lambda_2}{z_0}
       -\frac{\lambda_1^2+18\lambda_1+6\lambda_2+2\lambda_3}{z_0^2}
       +O(z_0^{-3})\right].
\end{align*}
There is a finite recursive extension to every order: substitute
$z=z_0+\sum_{r\ge1}Q_r z_0^{-r}$ in
\[
 (z-z_0)-3\log(z/z_0)+\sum_{r\ge1}\lambda_r z^{-r}=0
\]
and cancel powers successively. In particular,
\[
 n-n_0(a_n)\longrightarrow
 \eta=-\frac{\lambda_1}{2K}=\frac16+\frac{36}{13\pi^2}
 \approx0.447248.
\]
The bound $0<\eta<1/2$ follows already from $\pi>3$. Thus nearest-integer rounding of $n_0(a_n)$ recovers $n$ for all
sufficiently large $n$. This is an eventual, non-effective statement until
explicit remainder constants are provided; it does not justify such rounding
for arbitrary $y$ or at a claimed finite numerical cutoff.

\begin{remark}[write, 5 October 2026: the limit $\eta$]\label{lcp:tr:rem:eta}
The limit is proved as follows. Theorem~\ref{lcp:tr:thm:main} with $m=2$
gives \eqref{lcp:tr:eq:log-full} with $R=1$:
$\log a_n=g(z)+\log\mathcal A+\lambda_1/z+\OO(z^{-2})$, where $z=\beta\sqrt n$
and $g(x)=x-3\log x$. For large $n$, $a_n\ge\mathcal A\ee^3/27$, and
$z_0=z_0(a_n)\ge3$ satisfies $a_n=\mathcal A\ee^{z_0}/z_0^3$, that is
$\log a_n=g(z_0)+\log\mathcal A$; this is the defining equation of $W_{-1}$
after the substitution $w=-z_0/3$. Hence $g(z)-g(z_0)=-\lambda_1/z+\OO(z^{-2})$.
Since $g$ is increasing on $[3,\infty)$ and $g(z)\to\infty$, also
$z_0\to\infty$, and $g'\ge1/2$ on $[6,\infty)$ gives $z-z_0=\OO(1/z)$. By
the mean value theorem $g(z)-g(z_0)=(z-z_0)(1-3/\xi)$ with $\xi$ between
$z$ and $z_0$, so $z-z_0=-\lambda_1/z+\OO(z^{-2})=-\lambda_1/z_0+\OO(z_0^{-2})$
(the sign agrees with the first correction of the display above). Therefore
\[
 n-n_0(a_n)=\frac{z^2-z_0^2}{4K}=\frac{(z-z_0)(z+z_0)}{4K}
 =-\frac{\lambda_1}{2K}+\OO(n^{-1/2}),
\]
and $-\lambda_1/(2K)=(3+K/3)/(2K)=1/6+36/(13\pi^2)=\eta$. Finally
$\eta<1/2$ if and only if $\pi^2>108/13$, which holds since $\pi>3$; so
$n_0(a_n)\in(n-1/2,n+1/2)$ for all large $n$, and nearest-integer rounding
recovers $n$. Part~I's table of the leading inverse in Section~10.3
(``leading inverse minus $n$'', from $-0.5986$ at $n=50$ to $-0.4762$ at
$n=2000$) is $-(n-n_0(a_n))$ at the same points, since Part~I's $s_0$ is
$z_0$; it shows the approach to $-\eta=-0.44725\ldots$ from below.
\end{remark}


\subsection{All orders power log inversion}
Put
\[
 Y=\log y,\qquad T=3\log Y-\log\mathcal A.
\]
For $y=a_n\to\infty$ the complete inverse has the power--log form
\begin{equation}\label{lcp:tr:eq:z-inverse}
 z=Y+T+\sum_{r=1}^{M}\frac{P_r(T)}{Y^r}
       +O\left(\frac{(\log Y)^{M+1}}{Y^{M+1}}\right)
\end{equation}
for every fixed $M\ge0$. For $M=0$ the remainder is understood as
$O((\log Y)/Y)$. The first four polynomials are
\begin{align*}
 P_1(T)={}&3T-\lambda_1,\\
 P_2(T)={}&-\frac32T^2+(9+\lambda_1)T-3\lambda_1-\lambda_2,\\
 P_3(T)={}&T^3-\left(\lambda_1+\frac{27}{2}\right)T^2
 +(27+9\lambda_1+2\lambda_2)T
 -\lambda_1^2-9\lambda_1-3\lambda_2-\lambda_3,\\
 P_4(T)={}&-\frac34T^4+\left(\lambda_1+\frac{33}{2}\right)T^3
 -\left(\frac{33}{2}\lambda_1+3\lambda_2+81\right)T^2\\
 &+(3\lambda_1^2+54\lambda_1+15\lambda_2+3\lambda_3+81)T\\
 &-\frac{15}{2}\lambda_1^2-3\lambda_1\lambda_2
 -27\lambda_1-9\lambda_2-3\lambda_3-\lambda_4.
\end{align*}
Here is a finite rule for all orders. Let
$V_{m-1}(u)=T+\sum_{r=1}^{m-1}P_r(T)u^r$. Then
\[
 P_m(T)=-[u^m]\left\{
 V_{m-1}(u)-T-3\log(1+uV_{m-1}(u))
 +\sum_{j=1}^{m}\lambda_j u^j(1+uV_{m-1}(u))^{-j}\right\}.
\]
Each use requires only finite Taylor expansions. This also proves by induction
that $P_m$ is a polynomial of degree at most $m$. Substitution cancels the
logarithmic equation to the asserted order. The finite smooth logarithmic model
has derivative $1+O(z^{-1})$ with respect to $z$, so the mean value theorem
converts the residual into the inverse remainder in \eqref{lcp:tr:eq:z-inverse}.
No derivative of an unspecified asymptotic remainder is taken.

Squaring gives a particularly useful version directly for the area:
\begin{align*}
 n=\frac1{4K}\Bigg[{}&Y^2+2TY+T^2+6T-2\lambda_1\\
 &+\frac{3T^2+18T-6\lambda_1-2\lambda_2}{Y}\\
 &+\frac{-T^3+(54+6\lambda_1+2\lambda_2)T
 -\lambda_1^2-18\lambda_1-6\lambda_2-2\lambda_3}{Y^2}\Bigg]
 +O\left(\frac{(\log Y)^4}{Y^3}\right).
\end{align*}
For example, retaining only the first line gives an $o(1)$ error, namely
$O((\log Y)^2/Y)$, and therefore eventually recovers the exact-range inverse
by nearest-integer rounding. These are power--log asymptotic expansions, often
called logarithmic transseries; no assertion of convergence or of complete
exponentially small sectors is being made.

\subsection{Rigorous threshold brackets and limitations}
For fixed $M\ge1$, define
\[
 F_M(x)=\beta\sqrt x-\frac32\log x+\log\delta
         +\sum_{r=1}^{M}\frac{\lambda_r}{(\beta\sqrt x)^r}.
\]
This is strictly increasing for all sufficiently large $x$. Let $r_M(y)$ be its
large real root $F_M(r_M(y))=\log y$. Equation \eqref{lcp:tr:eq:log-full}, together with
$F_M'(x)\sim\beta/(2\sqrt x)$, shows that for some constants $C_M,Y_M>0$,
\begin{equation}\label{lcp:tr:eq:threshold-bracket}
 \left\lceil r_M(y)-C_M(\log y)^{-M}\right\rceil
 \ \le N(y)\le\ 
 \left\lceil r_M(y)+C_M(\log y)^{-M}\right\rceil
 \qquad(\log y\ge Y_M).
\end{equation}
To justify the bracket uniformly, use the logarithmic remainder on a fixed-size
neighborhood of $r_M(y)$, where $x\asymp(\log y)^2$. Increasing or decreasing $x$
by $C_M(\log y)^{-M}$ then dominates the error in $\log a_x$ at integer arguments.
The upper comparison produces a coefficient at least $y$; the lower comparison
excludes every smaller integer. Strict coefficient monotonicity completes the
argument.

On the exact range, the stronger unrounded conclusion is
$n=r_M(a_n)+O((\log a_n)^{-M})$. For general $y$, a step-valued threshold cannot
have an unqualified $o(1)$ expansion by a smooth function. The bracket is often
a single integer, and otherwise gives two consecutive possibilities once its
width is less than one. An explicit power--log approximation may replace
$r_M$, using its own remainder bound: for example the displayed expansion
through $Y^{-2}$ approximates $r_3(y)$ (and likewise $r_M(y)$ for $M\ge3$) with error $O((\log Y)^4/Y^3)$ in the corresponding bracket. The roots $r_1$ and $r_2$ omit terms needed at this accuracy.
The constants and the onset $Y_M$ in these asymptotic assertions have not been
made explicit. Thus the results prove eventual exact-range recovery and
asymptotic threshold bounds, not certified finite-data cutoffs.

\begin{remark}[write, 5 October 2026: step-valued thresholds]\label{lcp:tr:rem:step}
The sentence above that a step-valued threshold has no unqualified $o(1)$
expansion by a smooth function holds for every function $g$ continuous on
some $[y_1,\infty)$: $N(y)-g(y)$ does not tend to zero. Proof: by the strict
increase $a_{n+1}\ge a_n+1$ ($n\ge1$), $N(a_n)=n$ and $N(y)=n+1$ for
$a_n<y\le a_{n+1}$. If $N(y)-g(y)\to0$, choose $n$ with $a_n\ge y_1$ and
$|N(y)-g(y)|<1/3$ for all $y\ge a_n$. Then $g(a_n)<n+1/3$, while
$g(y)>n+2/3$ for $a_n<y\le a_{n+1}$; letting $y\downarrow a_n$, continuity
gives $g(a_n)\ge n+2/3$, a contradiction.
\end{remark}

\begin{writenote}[Added 5 October 2026, batch~101: an instance of repository results]
As for Parts~I and~II (the notes in Sections~9.2 and~\ref{lcp:ao:sec:inverse}),
this inversion is an instance of the transseries-and-inversion
volume~\cite{lcp-tai}, and no novelty is claimed for the method: the core
$z_0-3\log z_0=\log(y/\mathcal A)$ is \texttt{p0:thm:lambert-core} with
$a=1$, $b=-3$ and the branch $W_{-1}$; the corrected Lambert expansion and
the power--log inverse \eqref{lcp:tr:eq:z-inverse} are the reversion of
\texttt{p0:thm:perturbed-inversion}, in the Lambert-centred form
\texttt{p0:thm:lambert-centered}; and the brackets
\eqref{lcp:tr:eq:threshold-bracket} are part~(2) of
\texttt{p0:thm:staircase}. They are Part~II's
Corollary~\ref{lcp:ao:cor:brackets} with the order shifted by one: the
model $F_M$ here contains $\lambda_1,\ldots,\lambda_M$, that is $M+1$
coefficients, and its width $(\log y)^{-M}\asymp x^{-M/2}$ is Part~II's
$\delta_{M+1}$. Part~III's $z_0$ is Part~II's $y_0$ \eqref{lcp:ao:eq:W} and
Part~I's $s_0$ at the same target, and its $\lambda_r$ are the unshifted
counterparts of Part~II's $\ell_r$ \eqref{lcp:ao:eq:ell}. New here are the
limit $\eta$ and the eventual rounding rule of
Remark~\ref{lcp:tr:rem:eta}, which Part~I (rounding of the interpolation
$\nu$) and Part~II (brackets) do not state.
\end{writenote}



\section{Exact checks and numerical diagnostics}\label{lcp:tr:sec:checks}
The accompanying code makes three logically distinct checks. First, an integer coefficient recurrence reconstructs all $2001$ coefficients $a_0,\ldots,a_{2000}$ in the public OEIS b-file~\cite{bfile}, with exact agreement. Second, a literal expansion of \eqref{lcp:tr:eq:original}--\eqref{lcp:tr:eq:frec} independently agrees through $n=400$. Third, exact recurrence generation extends the sequence through $n=8000$. The additional terms are generated data, not independently verified OEIS data.

For reproducibility, the retained normalized b-file snapshot (one index/value pair per line) has SHA-256
\begin{center}\small\texttt{d036795c5c67b45836e4dbd76036ef0af520}\\
\texttt{baf4751aaf88421f853e5cedf027}.
\end{center}
The normalized polynomial recurrence used for exact computation is
\[
 b_n(q)=\frac{2b_{n-1}(q)-b_{n-2}(q)}{(1-q^n)^2},\qquad
 b_{-1}(q)=b_0(q)=1,
\]
where $b_n(q)=u_n(q)/(q;q)_n^2$ and the symbol $b_n(q)$ here denotes a generating series, distinct from the numerical asymptotic constants $b_r$ in the preceding section. Since
\[
 A(q)=1+\sum_{n\ge1}q^n(1-q^n)b_n(q),
\]
truncated integer power series suffice. Multiplication by $(1-q^n)^{-2}$ uses the exact coefficients $j+1$ at powers $q^{nj}$.

The leading normalized ratios are
\begin{center}
\begin{tabular}{rr}
\toprule
$n$ & $a_n/(\delta n^{-3/2}\ee^{\beta\sqrt n})$\\
\midrule
100 & 0.88579013894172337144\\
500 & 0.95077637637757003576\\
1000 & 0.96568896870264176064\\
2000 & 0.97602525341243592074\\
4000 & 0.983206124724057642\\
8000 & 0.988210257959352301\\
\bottomrule
\end{tabular}
\end{center}
Table~\ref{lcp:tr:tab:errors} uses relative prediction error $\widehat a_n/a_n-1$ after retaining the stated order in $z^{-1}$, where $z=\beta\sqrt n$.
\begin{table}[ht]
\centering
\begin{tabular}{rrr}
\toprule
Retained correction & $n=2000$ & $n=8000$\\
\midrule
$z^{-1}$ & $+8.723657094\cdot10^{-4}$ & $+2.308154832\cdot10^{-4}$\\
$z^{-2}$ & $-1.309262065\cdot10^{-4}$ & $-1.691475405\cdot10^{-5}$\\
$z^{-4}$ & $-1.443558234\cdot10^{-6}$ & $-4.735525978\cdot10^{-8}$\\
$z^{-8}$ & $-1.184484504\cdot10^{-9}$ & $-2.517837511\cdot10^{-12}$\\
\bottomrule
\end{tabular}
\caption{Exact-data diagnostics for the algebraic coefficient expansion.}\label{lcp:tr:tab:errors}
\end{table}

The leading Lambert inverse biases $n-n_0(a_n)$ at $n=2000,4000,8000$ are, respectively,
\[
 0.47624742230330,\qquad0.46806132666914,\qquad0.46212708952736,
\]
consistent with the proven limit $0.447248405983909726\ldots$. The fourth-order corrected Lambert expression has bias approximately $-1.44654\cdot10^{-6}$ at $n=8000$; the power--log inverse through $Y^{-3}$ has bias approximately $-7.80091\cdot10^{-6}$ there.

A finite scan finds that nearest-integer rounding of the leading Lambert model succeeds for every $556\le n\le8000$, fails for $2\le n\le555$, and has no real large-branch value at $n=1$. This is a diagnostic observation only. It does not prove that $556$ is a global cutoff; the theorem proves eventual recovery without an effective onset.

The source package includes exact coefficient generation, the normalized b-file snapshot, high-precision diagnostics, a symbolic expansion of the regular multiplier, and an independent Gaussian-saddle calculation. The scripts use explicit checks that remain active under Python's optimized mode. Numerical agreement is supplementary evidence; the asymptotic theorem rests on the estimates in Sections~2--4.

\begin{writenote}[Added 5 October 2026, batch~101: files and rerun]
In the repository the scripts are
\texttt{code/\allowbreak 03-transfer-*}, the recorded outputs are
\texttt{data/\allowbreak 03-transfer-*.json}, and the exact-checks note is
\texttt{03-transfer-exact-checks-README.md}; the report's README maps every
delivered name and gives a rerun recipe on a scratch copy in the delivered
layout. The b-file snapshot (\texttt{b126764\_web.txt}) is not shipped: it is
byte-identical to the live OEIS b-file (checked at intake, 5~October 2026)
and equals Part~II's fixture \texttt{data/\allowbreak 02-all-orders-coefficients-2000.txt}
plus a final newline. All 2001 terms of that b-file therefore agree with the
coefficients computed in each of the three parts. The terms $2001\le n\le8000$
were also reproduced at intake by Part~I's \texttt{code/verify.py}, a
different implementation, so they rest on two programs, though still on
the one cited generating function. The OEIS data are licensed CC BY-SA 4.0.
\end{writenote}


\section{Scope of the result and further questions}\label{lcp:tr:sec:context}
The exact area generating function is established in the 2007 enumeration paper~\cite{castiglione}; the 2023 numerical study~\cite{gk} explicitly presents the leading equivalent as a conjecture. The argument here supplies an exact product decomposition, proves that equivalent, and gives its all-orders algebraic refinement and inverse expansions. A targeted public-literature search through 2 October 2026 found no subsequent published proof of this equivalent. That is a bounded search statement, not a claim of exhaustive coverage of unpublished work. In particular, the independently derived area generating function in W.~R.~G.~James's 2005 Melbourne thesis, cited by the 2007 paper, was not checked in full.

A separate unpublished manuscript, \emph{L-convex polyominoes and a colored-partition comparison}, prepared for Vladimir Reshetnikov with ChatGPT and dated 1 October 2026~\cite{lcp-tr-earlier}, also proves the leading equivalent and three coefficient corrections, and develops an all-orders radial expansion and inverse asymptotics. It uses the same auxiliary product and eta quotient, and its radial coefficients through $t^8$ agree with those above. Its coefficientwise colored-partition comparison is a different route to finite-order coefficient estimates; that manuscript explicitly leaves the all-orders pointwise coefficient expansion outside its proved scope. The present common-partition-factor minor-arc bound and Chebyshev interpolation into the saddle window supply that additional coefficient-transfer step. Accordingly, the distinction claimed here is the all-orders pointwise coefficient theorem, rather than the first proof of the leading equivalent.

\begin{remark}[write, 5 October 2026: a stale novelty sentence]\label{lcp:tr:rem:novelty}
The sentence ``Accordingly, the distinction claimed here is the all-orders
pointwise coefficient theorem, rather than the first proof of the leading
equivalent'' was accurate against Part~I's manuscript, which leaves that
theorem open (Section~8.2, research question~1, the claim ledger). It is
stale against this report. Part~II, a manuscript dated 1~October 2026 and
written into the report on 2~October 2026, proves the all-orders
coefficient theorem (Theorem~\ref{lcp:ao:thm:main}); Part~III's manuscript
is dated a day later and does not know it. The corrected reading:
Theorem~\ref{lcp:tr:thm:main} is the same theorem as
Theorem~\ref{lcp:ao:thm:main} (Remark~\ref{lcp:tr:rem:shift}), and the
distinction of Part~III is its route, a fully written coefficient transfer
(Section~\ref{lcp:tr:sec:transfer} with Lemma~\ref{lcp:tr:lem:hankel}) that
needs no complex estimate of $h$, of $M$ or of the remainder, where
Part~II sketches its whole-circle Lemmas~\ref{lcp:ao:lem:global}
and~\ref{lcp:ao:lem:minor}. The manuscript's statement that it is not the
first proof of the leading equivalent stands; that equivalent is now proved
three times in this report (Section~\ref{lcp:tr:sec:cross}). The same
sentence appears in the delivery README (not shipped) and in the key
\texttt{known\_related\_unpublished\_manuscript} of the shipped
\texttt{data/\allowbreak 03-transfer-PROVENANCE.json}, which stays
byte-identical.
\end{remark}


The proof is specific to L-convex area enumeration. It does not depend on the distinct 2024 Z-convex area series, whose data and ensuing analysis were later flagged as incorrect on \href{https://blogs.unimelb.edu.au/tony-guttmann/sample-page/publications/publications-2020/}{Tony Guttmann's publication list}. Nor does a perimeter generating function resolve the area problem.

\begin{writenote}[Added 5 October 2026, batch~101]
Checked at intake on 5~October 2026: Tony Guttmann's publication list
annotates the 2024 paper \emph{Asymptotics of Z-convex polyominoes}
(cited here as~\cite{zconvex}) with the remark that its series data are
wrong, so the subsequent analysis is too. The remark is the author's own;
no erratum was examined. Part~I cites that paper only as a research target
(question~7 of Section~11, with a dated note) and Part~II for the
conjectural status of the L-convex asymptotic, which the note does not
affect.
\end{writenote}


Several natural extensions remain.
\begin{enumerate}
\item Determine the genuinely exponentially small terms beyond \eqref{lcp:tr:eq:allcoeff}. The exact remainder \eqref{lcp:tr:eq:Rformula}, modular corrections to \eqref{lcp:tr:eq:Mexp}, and neighborhoods of roots of unity can all contribute. The minor-arc estimate proves their irrelevance to every algebraic order but does not identify their individual leading terms.
\item Make remainder constants effective. This would turn the eventual inverse-rounding results into certified finite cutoffs and provide rigorous finite-$n$ error bars for the coefficient expansion.
\item Study the large-order growth of the rational $d_j$. The finite recurrence for $h$ and the alternating series \eqref{lcp:tr:eq:T} make such a study computationally accessible, but factorial growth, Borel singularities, and optimal truncation are not proved here.
\item Introduce a width or height parameter in the auxiliary recurrence, and investigate joint area--shape asymptotics. The double-pole product and its dual recurrence suggest a route, but a bivariate theorem requires its own uniform estimates.
\end{enumerate}

Two other sequence families remain distinct possible research directions rather than consequences of this result. The normalized growth constant for $000$-avoiding ascent sequences, OEIS A202058, is conjectural in the 2022 numerical study~\cite{lcp-tr-conway}; its finer stretched-exponential scale is not sufficiently settled to assume the present method applies. For OEIS A229741 and the late coefficients A260879, factorial-series composition already falls under general asymptotic machinery~\cite{lcp-tr-borinsky}; a defensible new target would concern late-term growth or optimal truncation, with an additional focused prior-art check.

\begin{writenote}[Added 5 October 2026, batch~101]
Under Vladimir's standing rule these questions, and the two sequence
families of the preceding paragraph, are restated with their status in
Section~\ref{lcp:tr:sec:further}, where they are matched with Part~I's
research questions and Part~II's open problems. The preceding paragraph is
partly stale against the repository: the growth constant of A202058 is
proved, and a stretched-exponential scale excluded, in the report
\nolinkurl{a202058-ascent-000-growth}, and the late-coefficient equivalent of
A260879 is proved in the transseries collection
(Section~\ref{lcp:tr:sec:further}, item~5, gives the details).
\end{writenote}


\section{A Lambert series for the multiplier (the manuscript's Appendix~A)}\label{lcp:tr:sec:lambert}
This identity is not needed for the proof of the main theorem, but can be useful for further analytic work. Let
\[
 S(z)=\sum_{n\ge0}\frac{u_nz^n}{p_n},\qquad S(z)=P(z)H(z).
\]
The recurrence gives $(1-z)^2S(z)=(1+qz^2)S(qz)-z$. Division by the homogeneous product and iteration yield
\[
 H(z)=1-\sum_{j\ge0}
 \frac{q^jz(q^{j+1}z;q)_\infty^2}{(-q^{2j+1}z^2;q^2)_\infty}.
\]
Since $H(1)=h$, this gives
\[
 h=1-\frac{p_\infty^2}{(-q;q^2)_\infty}
 \sum_{j\ge0}\frac{q^j(-q;q^2)_j}{p_j^2}.
\]
Take $a=i\sqrt q$, $b=-i\sqrt q$. Heine's first transformation~\cite{lcp-tr-dlmf}, with parameters $a,b,c=q$ and argument $q$, cancels the infinite products and gives
\begin{equation}\label{lcp:tr:eq:LambertH}
 h(q)=1-\sum_{n\ge0}\frac{b^n}{1-aq^n}.
\end{equation}
The series converges absolutely for $0<|q|<1$ with either choice of $\sqrt q$; swapping the square root interchanges $a,b$ and leaves the value unchanged by the symmetric double-series expansion. The value at $q=0$ is obtained by continuation. Thus \eqref{lcp:tr:eq:LambertH} is an analytic representation in the open unit disk, not a choice-dependent extension. Its use for detailed secondary exponential analysis would still require appropriate uniform estimates.

\begin{writenote}[Added 5 October 2026, batch~101]
This is Part~II's \eqref{lcp:ao:eq:Lambert}, by the same Heine
transformation; Part~II's Lemma~\ref{lcp:ao:lem:H} gives the formula for
$H(z)$. Here $S(z)$ is Part~II's $H(z)$ (both are
$\sum f_{n-1}z^n/(q;q)_n$), $P(z)$ is Part~II's $G(z)$, and Part~III's $H(z)$
is the bracket of \eqref{lcp:ao:eq:Hsolution}. Part~II goes on to the
partial theta form \eqref{lcp:ao:eq:theta}, so that
$h=1-\sum_{k\ge0}q^{k(k+1)}(1-q^{2k+1})/(1+q^{2k+1})$.
\end{writenote}

\section{Part III: further questions and research}\label{lcp:tr:sec:further}
This section was written when Part~III was added. Under Vladimir's
standing rule of 4~October 2026, every claim of the manuscript that it does
not prove is stated here as an open question, with its source, the
argument sketched for it and what is missing. Two sketched statements of
the manuscript are proved in this write and are listed for the record. No
claim of the manuscript was found to be false; one sentence is stale
(Remark~\ref{lcp:tr:rem:novelty}).
\begin{enumerate}[leftmargin=*,label=\textbf{\arabic*.},itemsep=0.5em]
\item \textbf{Exponentially small terms.} (Section~\ref{lcp:tr:sec:context},
 question~1.) Identify the exponentially small contributions beyond
 \eqref{lcp:tr:eq:allcoeff}: from the remainder $R$ of
 \eqref{lcp:tr:eq:Rformula}, from modular corrections to
 \eqref{lcp:tr:eq:Mexp}, and from neighbourhoods of other roots of unity.
 The minor-arc bound \eqref{lcp:tr:tr:8} proves only that they do not affect
 any algebraic order; it does not give their leading terms. The real bound
 $R/M\le qp_\infty^2/(1-q^2)$ of \eqref{lcp:tr:eq:Rbound} is a starting
 point on the radius, not on the circle. This is Part~II's open problem~2
 and Part~I's research question~5.
\item \textbf{Effective constants and the onset 556.}
 (Section~\ref{lcp:tr:sec:context}, question~2, and
 Section~\ref{lcp:tr:sec:checks}.) Make the constants of
 Section~\ref{lcp:tr:sec:transfer} explicit: $c$ in \eqref{lcp:tr:tr:7},
 the implied constants in \eqref{lcp:tr:tr:5} and \eqref{lcp:tr:tr:10}, and
 the remainder constants of the real expansion. That would give certified
 error bars for $a_n$ and turn the eventual rounding rule of
 Remark~\ref{lcp:tr:rem:eta} into a theorem with an explicit starting
 index. Concretely: is nearest-integer rounding of $n_0(a_n)$ correct for
 every $n\ge556$? The finite scan supports it through $n=8000$: the bias
 $n-n_0(a_n)$ is $0.4999927$ at $n=556$ and $0.4621$ at $n=8000$, its
 smallest value in the scan, but nothing proves that it stays below $1/2$
 beyond the scanned range. This is Part~II's open
 problem~1 and Part~I's research question~2.
\item \textbf{Large-order behaviour.} (Section~\ref{lcp:tr:sec:context},
 question~3.) Determine the growth of the rational $d_j$ (and so of $c_r$):
 factorial growth, Borel singularities, optimal truncation. The tail bound
 \eqref{lcp:tr:eq:htail} and the alternating series \eqref{lcp:tr:eq:T}
 make the computation accessible; nothing is proved. This is Part~II's open
 problem~3.
\item \textbf{Area and shape.} (Section~\ref{lcp:tr:sec:context},
 question~4.) Introduce a width or height parameter into the auxiliary
 recurrence and prove joint area--shape asymptotics. The source suggests
 the double-pole product \eqref{lcp:tr:eq:P} and the duality
 \eqref{lcp:tr:eq:duality} as a route; a bivariate theorem needs its own
 uniform estimates. This continues Part~II's open problem~4 and Part~I's
 research question~7 in another direction.
\item \textbf{Two other sequence families.} (Section~\ref{lcp:tr:sec:context}.)
 The growth constant of $000$-avoiding ascent sequences (OEIS A202058) is
 conjectural in~\cite{lcp-tr-conway}, and its stretched-exponential scale is
 not settled enough to assume that the present method applies. For OEIS
 A229741 and the late coefficients A260879, the general machinery
 of~\cite{lcp-tr-borinsky} already covers factorial-series composition; a new
 target would concern late-term growth or optimal truncation, after a
 focused prior-art check. The source offers both as directions, not
 results.

 \emph{Status in the repository (dated 5~October 2026, this write).} Both
 families are now treated elsewhere in the repository, so the item is
 re-scoped. The report
 \nolinkurl{oeis-sequence-asymptotics/a202058-ascent-000-growth} proves the
 normalized growth constant $\mu=8/(3\pi^2)$ of A202058 conjectured
 in~\cite{lcp-tr-conway} (its Part~I, Theorem~1.1), excludes every
 stretched-exponential factor $\exp(cn^\sigma)$ with $c,\sigma>0$ (its
 Corollary~9.2), and proves $\limsup_n\log c_n/(\log n)^2=2/3$ for
 $c_n=a_n\mu^{-n}/n!$ (its Part~III); the pointwise law
 $\log c_n\sim\tfrac23(\log n)^2$ is open there. The volume \emph{Late coefficients of the factorially forced
 Catalan recurrence} of the series-and-transseries collection (directory
 \texttt{Late\_\allowbreak Coefficients\_\allowbreak Factorially\_\allowbreak Forced\_\allowbreak Catalan\_\allowbreak Recurrence}
 under \nolinkurl{Analysis/Transseries/docs/series-and-transseries})
 proves the late-coefficient equivalent conjectured in A260879,
 $c_k\sim\Gamma(k)/(\log2)^k$, with every fixed algebraic correction, and
 asserts no optimal truncation or Borel continuation. What remains of the
 source's suggestion is the pointwise law for A202058 and optimal
 truncation for A229741 and A260879; neither is connected with the method
 of this report.
\item \textbf{The finite denominators.} (This write, from
 Remark~\ref{lcp:tr:rem:asc}.) $Q_m(q)$ is the Al-Salam--Chihara polynomial
 $Q_m(1;\ii\sqrt q,-\ii\sqrt q\,|\,q)$. Does that identification help to
 locate the zeros of $Q_m$ as polynomials in $q$, the question left open by
 Part~I's research questions~1 and~4? Nothing is claimed.
\end{enumerate}
\paragraph{Sketched statements proved in this write.}
\begin{itemize}[leftmargin=*,itemsep=0.15em]
\item The remark that $Q_n$ is a specialization of the Al-Salam--Chihara
 polynomials, stated without proof in Section~\ref{lcp:tr:sec:product}:
 Remark~\ref{lcp:tr:rem:asc}.
\item The side remark of Section~\ref{lcp:tr:sec:inverse} that a
 step-valued threshold has no unqualified $o(1)$ expansion by a smooth
 function: Remark~\ref{lcp:tr:rem:step}, for every continuous function.
\end{itemize}
The two steps that the manuscript states briefly, the limit $\eta$ and the
passage from the saddle window to Bessel functions, are written out in
Remark~\ref{lcp:tr:rem:eta} and Lemma~\ref{lcp:tr:lem:hankel}.

\section{Provenance of Part III, and where the write had to choose}\label{lcp:tr:sec:provenance}
This section was written when Part~III was added.

\paragraph{The manuscript.} Part~III
(Sections~\ref{lcp:tr:sec:main}--\ref{lcp:tr:sec:lambert}) is bundle
Report~95, ``All Order Area Asymptotics for L Convex Polyominoes'', dated
2~October 2026, batch~101 of the intake: archive
\texttt{L\_\allowbreak Convex\_\allowbreak Polyomino\_\allowbreak Area\_\allowbreak Asymptotics\_\allowbreak Source.zip},
arrival commit \texttt{60f54ea06}, placed in \texttt{f7c612c72}. It
contributed the duality proof of the decomposition, the transfer proof of
the all-orders theorem, the strict increase, the limit $\eta$, the
inverse in $z=\beta\sqrt n$ and $Y=\log y$, and exact coefficients through
$n=8000$. It names no ProveIt commit, so it has no pin. The report is now
built from three manuscripts; Parts~I and~II are described in
Section~\ref{lcp:ao:sec:provenance}.

\paragraph{Where the write had to choose.}
\begin{itemize}[leftmargin=*,itemsep=0.15em]
\item \emph{Placement.} Part~III is an addition to this report, not a
 separate report: it re-proves the report's main theorem by a third route.
 It is printed after Part~II's closing Section~23 and before Part~I's
 appendices, so no number of Parts~I or~II moves. The manuscript's
 Appendix~A is printed as Section~\ref{lcp:tr:sec:lambert}.
\item \emph{Letters.} The manuscript's letters are kept, with a clash table
 (Section~\ref{lcp:tr:sec:notation}) and superscripts $\mathrm I$,
 $\mathrm{II}$ in write notes. Its macro \verb|\ee| is the report's; its
 \verb|\dd| (a spaced differential, unused in its text) and its unused
 \verb|\N|, \verb|\Z|, \verb|\R|, \verb|\qP| were dropped; the pandoc
 command \verb|\tightlist| was replaced by an \verb|itemize| option.
\item \emph{Labels.} The manuscript's 54 labels carry the prefix
 \texttt{lcp:tr:} (so its \texttt{tr:2} is \texttt{lcp:tr:tr:2}); the write
 added section labels, two tables, the conversion remark and equation,
 four further remarks, a lemma and its equation, all with the same prefix.
\item \emph{Duplicated results.} Results proved in Parts~I or~II keep
 Part~III's proof as well, with cross-references
 (Section~\ref{lcp:tr:sec:cross}); no proof was shortened. Part~III's
 coefficients are unshifted, and Remark~\ref{lcp:tr:rem:shift} converts them.
\item \emph{Bibliography.} The manuscript's entries for
 Castiglione--Frosini--Munarini--Restivo--Rinaldi, Guttmann--Kot\v e\v sovec
 and OEIS A126764 cite the same works as \cite{castiglione,gk,oeis,bfile}
 and were merged into them (its page reference for the first is noted in
 that entry). Its three DLMF entries became one entry \cite{lcp-tr-dlmf};
 its entries for Bringmann--Jennings-Shaffer--Mahlburg, Conway et al.\ and
 Borinsky are \cite{lcp-tr-bjm,lcp-tr-conway,lcp-tr-borinsky}; its entry for
 Part~I's manuscript is \cite{lcp-tr-earlier}, with a pointer to Part~I. The
 write added \cite{lcp-tr-kls} for Remark~\ref{lcp:tr:rem:asc}.
\item \emph{Additions.} Section~\ref{lcp:tr:sec:intro}, the two closing
 sections, the remarks and the lemma headed ``write'', and the
 \textsf{[write]} notes are editorial. No statement, proof or number of the
 manuscript was changed.
\item \emph{Not shipped.} The manuscript's \LaTeX{} sources, all printed
 here: the main file \texttt{article.tex}, its inputs
 \texttt{transfer.tex}, \texttt{inverse.tex}, \texttt{diagnostics.tex},
 \texttt{context.tex} and \texttt{references.tex}, and the expanded
 \texttt{article\_\allowbreak standalone.tex}. Also not shipped: its 16-page PDF, its
 delivery README, \texttt{verify\_package.py} (a SHA-256 check against
 \texttt{PROVENANCE.json}), \texttt{build\_local.sh} (a PDF build of the
 unshipped sources), the b-file snapshot (a copy of OEIS data that
 Part~II's fixture already carries) and a duplicate of the symbolic results.
 The archive survives in the arrival commit; retrieve it with
 \texttt{git show 60f54ea06:\allowbreak docs/\allowbreak incoming/\allowbreak L\_Convex\_\allowbreak Polyomino\_\allowbreak Area\_\allowbreak Asymptotics\_\allowbreak Source.zip}.
\end{itemize}


\appendix
\section{Reproduction commands and files}
The archive contains this standalone \LaTeX{} source, the compiled PDF,
two verification scripts, exact coefficient data, and their verification
logs. Build the document with
\begin{verbatim}
pdflatex -interaction=nonstopmode -halt-on-error lconvex_asymptotics.tex
pdflatex -interaction=nonstopmode -halt-on-error lconvex_asymptotics.tex
\end{verbatim}
Run the checks from the archive directory with
\begin{verbatim}
python verify.py --nmax 2000 --out checks
python verify_symbolic.py
\end{verbatim}
The first command requires only the Python standard library.
The second additionally requires SymPy and mpmath.
The quadratic number of arithmetic operations in the coefficient checker
is not a quadratic bit-complexity claim: the integer sizes grow with $n$.
The code uses $\OO(N)$ large-integer storage for each of a fixed number of
truncated series. No external data download is required for reproduction.

\begin{writenote}[Added 2 October 2026, batch~75]
In the repository the source is \texttt{article.tex}, the scripts are in
\texttt{code/}, and the contents of the delivered \texttt{checks/}
directory are in \texttt{data/}; the compiled PDF of the archive is not
shipped. \texttt{verify.py} writes into \texttt{--out} relative to the
working directory, while \texttt{verify\_symbolic.py} reads
\texttt{coefficients.csv} from, and writes into, a \texttt{checks/}
directory beside itself. The commands above therefore work only in the
delivered layout; the report's README recreates it on a scratch copy. On
2~October~2026 that recipe reproduced \texttt{coefficients.csv} and both
tables byte for byte and the other outputs apart from line endings.
\end{writenote}

\section{A short review checklist}
A reader auditing the proof can concentrate on four potential failure
points. First, retain the slope term when $z\to1$ in
\eqref{lcp:eq:Cdiff}; setting $(1-z)^2C(z)$ to zero would give a false
identity. Second, distinguish coefficientwise positivity from positivity
on the real interval; the proof of \eqref{lcp:eq:majorants} establishes the
stronger statements explicitly. Third, the finite-cut error is estimated
on the \emph{whole} coefficient circle using positive tails and the
controlled full-circle integral of $Q_3^{-2}$. Fourth, the real-axis all-orders
series is not used as a substitute for this complex estimate.
These checks isolate the substantive logical steps of the argument.

\begingroup
\raggedright
\begin{thebibliography}{9}
\bibitem{oeis}
The OEIS Foundation Inc., \emph{A126764: Number of L-convex polyominoes
with $n$ cells}. Definition, fixed-polyomino convention, and
Guttmann--Kot\v e\v sovec conjecture dated June 9, 2021.
\url{https://oeis.org/A126764}; internal entry
\url{https://oeis.org/A126764/internal}. Accessed October 1, 2026.

\bibitem{castiglione}
G.~Castiglione, A.~Frosini, E.~Munarini, A.~Restivo, and S.~Rinaldi,
\emph{Combinatorial aspects of L-convex polyominoes},
European Journal of Combinatorics \textbf{28} (2007), 1724--1741.
\url{https://doi.org/10.1016/j.ejc.2006.06.020}.
The exact recurrence used here is also explicitly reproduced in
\cite{gk}, equation (1.1).
Part~III cites its equations (24)--(25), pp.~1738--1739 (note added in
batch~101).

\bibitem{gk}
A.~J.~Guttmann and V.~Kot\v e\v sovec,
\emph{A numerical study of L-convex polyominoes and 201-avoiding
ascent sequences}, S\'eminaire Lotharingien de Combinatoire,
\textbf{87B} (2023), Lattice Path Combinatorics special issue, article 2, 11 pp.;
arXiv:2109.09928, version 3 (2023).
\url{https://arxiv.org/abs/2109.09928}.
Journal page: \url{https://www.mat.univie.ac.at/~slc/wpapers/lattpath/2.html}.

\bibitem{bfile}
V.~Kot\v e\v sovec, \emph{Table of $n,a(n)$ for $n=0,\ldots,2000$},
b-file for OEIS A126764.
\url{https://oeis.org/A126764/b126764.txt}.
Selected anchor values were checked against this source.

\bibitem{dlmf}
NIST Digital Library of Mathematical Functions,
\emph{Section 27.14: Unrestricted partitions}, especially
27.14.12--27.14.14 on Dedekind's eta function and its transformation.
\url{https://dlmf.nist.gov/27.14}. Accessed October 1, 2026.

\bibitem{zconvex}
A.~J.~Guttmann and P.~Massazza,
\emph{Asymptotics of Z-convex polyominoes},
RAIRO--Theoretical Informatics and Applications (2024).
\url{https://www.rairo-ita.org/articles/ita/pdf/2024/01/ita220046.pdf}.
Cited as a further research target, not as an input to the proof.
Part~II cites its p.~19, which repeats that the L-convex asymptotic is
conjectural (note added in batch~77P5).
In batch~101 (checked 5~October 2026) the first author's publication
list was found to annotate this paper: its series data are wrong, and so
is the analysis based on them (note in Section~\ref{lcp:tr:sec:context}).

\bibitem{proveit}
V.~Reshetnikov, \emph{ProveIt}, mathematical research and formalization
repository. \url{https://github.com/VladimirReshetnikov/ProveIt}.
Repository overview and targeted A126764 search inspected October 1, 2026.

\bibitem{lcp-tai}
\emph{Transseries and inversion}, a volume of the ProveIt repository:
\nolinkurl{Analysis/Transseries/docs/series-and-transseries/Transseries_And_Inversion/transseries_and_inversion.tex}
(theorems \texttt{p0:thm:lambert-core}, \texttt{p0:thm:lambert-centered},
\texttt{p0:thm:staircase}; Section \texttt{p3:sec:lambert},
\texttt{p3:thm:core-correction}, \texttt{p3:thm:staircase}). Added in
writing, batch~75. Part~II also cites \texttt{p0:thm:perturbed-inversion}
(batch~77P5).

\bibitem{lcp-klarner}
\emph{Finite-prefix corrections and dual certificates for polyomino growth},
a report of the ProveIt research-report collection:
\nolinkurl{SetTheory/Cardinals/docs/reports/enumerative-combinatorics/polyomino-growth-finite-prefix-corrections}.
Added in writing, batch~75.

\bibitem{lcp-ao-dlmf}
NIST Digital Library of Mathematical Functions, version 1.2.8, 15 September 2026. Heine and Fine transformations, \href{https://dlmf.nist.gov/17.6}{\S17.6}; eta transformation, \href{https://dlmf.nist.gov/23.18}{\S23.18}; modified Bessel asymptotics, \href{https://dlmf.nist.gov/10.40}{\S10.40}. Inspected 1 October 2026. Cited by Part~II (batch~77P5).
\bibitem{lcp-tr-earlier}
\emph{L-convex polyominoes and a colored-partition comparison: A proof of
the A126764 asymptotic, three correction terms, an all-orders radial
expansion, and inverse asymptotics}. Separate unpublished research
manuscript prepared for Vladimir Reshetnikov with ChatGPT, 1 October 2026,
19 pp. Cited by Part~III as printed; this manuscript is Part~I of the
present report (note added in batch~101).

\bibitem{lcp-tr-dlmf}
NIST Digital Library of Mathematical Functions: \S23.18, modular
transformations, and \S27.14, Dedekind eta products,
\url{https://dlmf.nist.gov/23.18}, \url{https://dlmf.nist.gov/27.14};
equations 10.40.1 and 10.17.1, modified Bessel asymptotics and coefficient
polynomials, \url{https://dlmf.nist.gov/10.40.E1},
\url{https://dlmf.nist.gov/10.17.E1}; equation 17.6.6, Heine's first
transformation, \url{https://dlmf.nist.gov/17.6.E6}. Cited by Part~III (the
manuscript's three DLMF entries, merged in batch~101).

\bibitem{lcp-tr-bjm}
K.~Bringmann, C.~Jennings-Shaffer and K.~Mahlburg, \emph{On a Tauberian
theorem of Ingham and Euler--Maclaurin summation}.
\url{https://arxiv.org/abs/1910.03036}. See the complex hypothesis and
counterexamples discussed in Theorem~1.1 and Section~3.2. Cited by
Part~III.

\bibitem{lcp-tr-conway}
A.~R.~Conway, M.~Conway, A.~Elvey Price and A.~J.~Guttmann,
\emph{Pattern-avoiding ascent sequences of length 3}, Electronic Journal of
Combinatorics \textbf{29}(4) (2022), P4.25.
\href{https://doi.org/10.37236/11266}{doi:10.37236/11266};
\url{https://arxiv.org/abs/2111.01279}. Cited by Part~III.

\bibitem{lcp-tr-borinsky}
M.~Borinsky, \emph{Generating asymptotics for factorially divergent
sequences}, Electronic Journal of Combinatorics \textbf{25}(4) (2018),
P4.1. \url{https://arxiv.org/abs/1603.01236}. Cited by Part~III.

\bibitem{lcp-tr-kls}
R.~Koekoek, P.~A.~Lesky and R.~F.~Swarttouw, \emph{Hypergeometric
Orthogonal Polynomials and Their $q$-Analogues}, Springer Monographs in
Mathematics, Springer, Berlin, 2010; \S14.8, the Al-Salam--Chihara
polynomials. Added in writing, batch~101 (Remark~\ref{lcp:tr:rem:asc}).
\end{thebibliography}
\endgroup
\end{document}
